%% Template for the submission to:
%%   The Annals of Statistics [AOS]
%%
%%%%%%%%%%%%%%%%%%%%%%%%%%%%%%%%%%%%%%%%%%%%%%
%% In this template, the places where you   %%
%% need to fill in your information are     %%
%% indicated by '???'.                      %%
%%                                          %%
%% Please do not use \input{...} to include %%
%% other tex files. Submit your LaTeX       %%
%% manuscript as one .tex document.         %%
%%%%%%%%%%%%%%%%%%%%%%%%%%%%%%%%%%%%%%%%%%%%%%

\documentclass[bj,authoryear]{imsart}

%% Packages
\RequirePackage{amsthm,amsmath,amsfonts,amssymb,booktabs,array}
%\RequirePackage[numbers]{natbib}
\RequirePackage[authoryear]{natbib}%% uncomment this for author-year citations
%\RequirePackage[colorlinks,citecolor=blue,urlcolor=blue]{hyperref}%% uncomment this for coloring bibliography citations and linked URLs
\RequirePackage{graphicx}%% uncomment this for including figures
\RequirePackage{makecell}

\startlocaldefs
%%%%%%%%%%%%%%%%%%%%%%%%%%%%%%%%%%%%%%%%%%%%%%
%%                                          %%
%% Uncomment next line to change            %%
%% the type of equation numbering           %%
%%                                          %%
%%%%%%%%%%%%%%%%%%%%%%%%%%%%%%%%%%%%%%%%%%%%%%
%\numberwithin{equation}{section}
%%%%%%%%%%%%%%%%%%%%%%%%%%%%%%%%%%%%%%%%%%%%%%
%%                                          %%
%% For Axiom, Claim, Corollary, Hypothesis, %%
%% Lemma, Theorem, Proposition              %%
%% use \theoremstyle{plain}                 %%
%%                                          %%
%%%%%%%%%%%%%%%%%%%%%%%%%%%%%%%%%%%%%%%%%%%%%%
\theoremstyle{plain}
\newtheorem{thm}{\textsc{Theorem}}[section]
\newtheorem{prop}[thm]{\textsc{Proposition}}
\newtheorem{cor}[thm]{\textsc{Corollary}}
\newtheorem{lem}[thm]{\textsc{Lemma}}
%%%%%%%%%%%%%%%%%%%%%%%%%%%%%%%%%%%%%%%%%%%%%%
%%                                          %%
%% For Assumption, Definition, Example,     %%
%% Notation, Property, Remark, Fact         %%
%% use \theoremstyle{remark}                %%
%%                                          %%
%%%%%%%%%%%%%%%%%%%%%%%%%%%%%%%%%%%%%%%%%%%%%%
%\theoremstyle{remark}
%\newtheorem{???}{???}
%\newtheorem*{???}{???}
%\newtheorem{???}{???}[???]
%\newtheorem{???}[???]{???}
%%%%%%%%%%%%%%%%%%%%%%%%%%%%%%%%%%%%%%%%%%%%%%
%% Please put your definitions here:        %%
%%%%%%%%%%%%%%%%%%%%%%%%%%%%%%%%%%%%%%%%%%%%%%
\newcommand{\david}[1]{{\small{\color{blue} \sf $\heartsuit$ [David: #1]}}}
\newcommand{\piotr}[1]{{\small{\color{red} \sf $\clubsuit$ Piotr: #1}}}
\newcommand{\cpz}[1]{{\color{red} #1}}
\newcommand{\paul}[1]{{\small{\color{green} \sf $\spadesuit$ Paul: #1}}}
\newcommand{\E}{\mathbb{E}}
\newcommand{\N}{\mathbb{N}}
\newcommand{\R}{\mathbb{R}}
\newcommand{\I}{\mathbb{I}}
\newcommand{\M}{\mathcal{M}}
\newcommand{\T}{\mathcal{T}}
\newcommand{\bs}{\boldsymbol}
\newcolumntype{M}[1]{>{\centering\arraybackslash}m{#1}}
\endlocaldefs

\begin{document}

\begin{frontmatter}
%%%%%%%%%%%%%%%%%%%%%%%%%%%%%%%%%%%%%%%%%%%%%%
%%                                          %%
%% Enter the title of your article here     %%
%%                                          %%
%%%%%%%%%%%%%%%%%%%%%%%%%%%%%%%%%%%%%%%%%%%%%%
\title{Improving variable selection properties by leveraging external data}
%\title{A sample article title with some additional note\thanksref{T1}}
\runtitle{Improving variable selection properties by leveraging external data}
%\thankstext{T1}{A sample of additional note to the title.}

\begin{aug}
%%%%%%%%%%%%%%%%%%%%%%%%%%%%%%%%%%%%%%%%%%%%%%%
%% Only one address is permitted per author. %%
%% Only division, organization and e-mail is %%
%% included in the address.                  %%
%% Additional information can be included in %%
%% the Acknowledgments section if necessary. %%
%% ORCID can be inserted by command:         %%
%% \orcid{0000-0000-0000-0000}               %%
%%%%%%%%%%%%%%%%%%%%%%%%%%%%%%%%%%%%%%%%%%%%%%%
\author[A]{\fnms{Paul}~\snm{Rognon-Vael}\ead[label=e1]{paul.rognon@gmail.com}}
\author[A]{\fnms{David}~\snm{Rossell}\ead[label=e2]{rosselldavid@gmail.com}}
\author[B]{\fnms{Piotr}~\snm{Zwiernik}\ead[label=e3]{piotr.zwiernik@utoronto.ca}}
%%%%%%%%%%%%%%%%%%%%%%%%%%%%%%%%%%%%%%%%%%%%%%
%% Addresses                                %%
%%%%%%%%%%%%%%%%%%%%%%%%%%%%%%%%%%%%%%%%%%%%%%
\address[A]{Department of Economics and Business, Universitat Pompeu Fabra \printead[presep={ ,\ }]{e1,e2}}
\address[B]{Department of Statistical Sciences, University or Toronto
\printead[presep={,\ }]{e3}}
\end{aug}

\begin{abstract}
Sparse high-dimensional signal recovery is only possible under certain conditions on the number of parameters, sample size, signal strength and underlying sparsity. We show that leveraging external information, as possible with data integration or transfer learning, allows to push these mathematical limits. Specifically, we consider external information that allows splitting parameters into blocks, first in a simplified case, the Gaussian sequence model, and then in the general linear regression setting. We show how external information dependent, block-based, $\ell_0$ penalties attain model selection consistency under milder conditions than standard $\ell_0$ penalties, and they also attain faster model recovery rates. We first provide results for oracle-based $\ell_0$ penalties that have access to perfect sparsity and signal strength information. Subsequently, we propose an empirical Bayes data analysis method that does not require oracle information and for which efficient computation is possible via standard MCMC techniques.

Our results provide a mathematical basis to justify the use of data integration methods in high-dimensional structural learning.

% \david{Below a possible abstract update}
% \color{blue}
% Sparse high-dimensional signal recovery is only possible under certain conditions on the number of parameters, sample size, signal strength and underlying sparsity (number of truly non-zero parameters). We show that these mathematical limits can be pushed when one has external information that allow splitting parameters into blocks. This occurs in many applications, including data integration and transfer learning tasks. We consider the Gaussian sequence model and linear regression, and show how block-based criteria attain model selection consistency under milder conditions than standard criteria, and they also attain faster model recovery rates. We first provide results for oracle-based criteria that have access to perfect sparsity and signal strength information, and subsequently an empirical Bayes-motivated block L0 penalty that does not require oracle information.
% \color{black}
% \paul{I like it. I am not sure readers will understand "block-based criteria", what about "block-based L0 penalty"? }\piotr{Nice. And getting there. :)}
\end{abstract}



\begin{keyword}[class=MSC]
\kwd[Primary ]{62F07}
\kwd[; secondary ]{62C12}
\kwd{62R07}
\end{keyword}

\begin{keyword}
\kwd{data integration}
\kwd{high-dimensional statistics}
\end{keyword}
\end{frontmatter}
%%%%%%%%%%%%%%%%%%%%%%%%%%%%%%%%%%%%%%%%%%%%%%
%% Please use \tableofcontents for articles %%
%% with 50 pages and more                   %%
%%%%%%%%%%%%%%%%%%%%%%%%%%%%%%%%%%%%%%%%%%%%%%
%\tableofcontents

%%%%%%%%%%%%%%%%%%%%%%%%%%%%%%%%%%%%%%%%%%%%%%
%%%% Main text entry area:

%\paul{\Large THIS IS THE CURRENT VERSION}

% \david{Just as a thought, how do you feel about "external data" instead of "external information"? And "using external data" over "integrating external information"? I'm thinking of this because the term "data" is used in transfer learning, a keyword that we should probably use.}


\paul{cite "Learning from a lot: Empirical Bayes for
high-dimensional model-based prediction"}
\paul{Add Info theoretic necessary conditions + Correct table comparison}
\paul{Add new lin reg necessary condition}
\paul{Correct proofS of nec betamin conditionS in GSM, no need for $\tau_j<\beta_{\min,j}^*$}
\paul{correct table of scalings}
\section{Introduction} 

High-dimensional inference theory relies on assumptions regarding sparsity and signal strength which,
although mathematically necessary, can be too strong in practice (e.g., see \cite{giannone2021economic}).
Our motivation is that in many applications one has external information regarding each parameter, e.g. its magnitude or its likelihood of being non-zero, that can be leveraged to relax said assumptions and enhance inference.
%In many modern applications, the sample size is insufficient to reliably do inference. This limitation has motivated the development of high-dimensional techniques that frequently rely on strict assumptions on sparsity and signal strength. Yet, often, external information is available and can be leveraged to enhance inference.
%Our motivation stems from the observation that, often, 
In particular, external information can guide our decisions regarding which parameters to include in a regression model. 
In a data integration setting, this information originates from previous datasets or similar selection problems (e.g., studying related cancer types). More generally, the information may also originate from each variable's inherent nature (e.g., clinical history vs. genomic markers, sociodemographics versus job history), or meta-covariates (e.g., functional annotations on genes), etc.
%This information may originate from each variable's inherent nature (e.g., clinical history versus genomic markers, sociodemographics versus job history),
%information from previous datasets or similar selection problems (e.g., studying related cancer types), or meta-covariates
%or external data that distinguishes relevant variables from irrelevant ones 
%(e.g., functional annotations on genes), for example. %, \paul{or from similar selection problems (e.g. transfer learning)}. 
We investigate this concept in the Gaussian sequence model and in linear regression, where external information partitions variables into blocks with potentially distinct characteristics. We show that said information allows pushing the mathematical conditions under which consistent model recovery is possible, and improving the associated rates.

Using external information is often advocated within the data integration and transfer learning literature, as joint learning can lead to more accurate inference than analyzing datasets separately. Indeed, numerous applied works employed external information to guide inference. For instance, \cite{stingo:2011, cassese:2014} proposed Bayesian variable selection methods for gene expression, where prior probabilities for non-zero coefficients depend on biological knowledge and meta-covariates. Additionally, \cite{chen_tinghuei:2021} predicted disease outcomes by allowing LASSO penalties to depend on functional annotation categories. Beyond regression, node and edge covariates have been incorporated in \cite{peterson} and \cite{jewson2023graphical} to drive edge inclusion in Gaussian graphical models, while \cite{schiavon2022generalized} used meta-covariates to determine non-zero loadings in factor models. In causal analysis, the inclusion of control covariates may be driven by their degree of association with the covariates of interest (referred to as treatments) 
%guide the inclusion of control covariates based on their degree of association with the treatments 
\citep{belloni2014inference,antonelli2021bayesian}. Collectively, empirical evidence consistently demonstrates improved structural learning when suitable external data is integrated. However, despite this empirical success, a theoretical framework explaining precisely why and how this occurs is not currently available.

%We show that incorporating external information allows pushing the theoretical limits under which consistent variable selection is possible, understood as recovering the true model with probability going to 1 as the sample size $n$ and the number of variables $p$ grow. 
Previous literature extensively explored how high-dimensional variable selection is constrained by inherent characteristics of the data, such as the number of samples $n$ and parameters $p$, the signal strength, the correlation between variables and the number of variables truly associated to the outcome \citep{buhlmannvandegeerbook,handbookbayesvarsel,wainwrighthighdimbook}. Here we analyse how conditions for consistency can be relaxed in the presence of external information. To make our ideas concrete, consider variable selection in the Gaussian linear regression
\begin{equation}\label{eq:linearmodel}
	\bs y \;=\; \bs X \bs \beta^* + \bs\epsilon,\qquad \bs \epsilon \sim N(0,\sigma^2 I_n),
\end{equation}
where $\bs X\in \R^{n\times p}$, $\bs y\in \R^n$, $\sigma >0$, and $\bs \beta \in \mathbb{R}^p$ are the data-generating parameters. A large class of methods operate by penalizing the size of an estimated $\hat{\bs  \beta}$. For instance, penalized likelihood methods optimize the log-likelihood plus a penalty term driven by the $\ell_q$ "norm" of $\hat{\bs  \beta}$ for $q\in [0,1]$  or folded concave penalties \citep{lasso,bertsimas}. 
In Bayesian settings, variable selection is often based on posterior model probabilities that are directly connected to $\ell_0$ penalties \citep{BIC,EBIC,rossell2022concentration}. 

Our interest is in the setting where the external information partitions the $p$ variables into $b$ blocks denoted $B_j \subset \{1,\ldots,p\}$, $j=1,\ldots,b$ such that key characteristics, like the level of sparsity or signal strength, is thought to potentially vary across the $B_j$'s. Assume that for each variable $i=1,\ldots,n$ we have external information $z_i$ that partitions variables into $b$ blocks,
%we have access to a vector of external information $z_i$ that partitions variables $i=1,\ldots,p$ into $b$ blocks, 
i.e. a partition function 
$$\Pi: z_i \to B_j \in \{1,...b\}.$$ 
We consider $\ell_0$ penalties that depend on the external information $\bs z=(z_1,\ldots,z_p)$ through the induced partition. More precisely, we introduce block $\ell_0$ penalties that allow modulating the strength of the penalty in each block. Unlike standard $\ell_0$ penalties, such as BIC \citep{BIC} or EBIC \citep{EBIC}, block penalties are non-exchangeable in the sense that the penalty for adding a variable $i$ may depend on its block $B_j=B_j(z_i)$. 

Our focus on $\ell_0$ penalties is motivated by their superior variable selection properties (e.g. see \cite{infotheowainwright}) %XXXX \paul{Do you have good references that would discuss that David? I am not finding them} 
which makes them particularly suitable to investigate the benefits of incorporating external information. 
Although our goal is to study fundamental properties of structural learning,
we remark that advances in optimization and MCMC methods made $\ell_0$ penalization more computationally tractable: it can be solved exactly for $p$ in the hundreds \citep{bertsimas} and with probability going to 1 with linear cost in $p$ using MCMC \citep{Yang2016,zhou2022} (under mild conditions). Moreover, our results in the Gaussian sequence model apply to essentially any penalized likelihood or Bayesian method, including $\ell_1$ penalties, because in that setting selection operates by thresholding \citep{davidomiros}. 

We mainly discuss external informed penalties that grow linearly in the model size. The linearity assumption allows useful connections with $\ell_1$ penalties and Bayesian methods. We show that our results are tight with respect to known limits on exact support recovery \citep{butucea2018,infotheowainwright}. Nonlinear $\ell_0$ penalties have however been shown to be optimal for estimation and prediction in the Gaussian sequence model \citep{Wu2013} %\paul{That's Harrisson Zhou...} 
and linear regression \citep{Bunea2007}. In the Supplement, we give support recovery guarantees and rates of convergence for nonlinear, externally informed, block penalties.  
\smallskip

\noindent \textbf{Our contributions:} 
%Our main contributions are as follows. 
We show that external information-dependent, block-based, $\ell_0$ penalties soften the theoretical conditions for consistent model recovery understood as recovering the support of $\bs \beta^*$ %exactly the true model 
with probability going to 1 as $n$ and $p$ grow. We consider first the sequence model, a simplified setting where we obtain very tight results. %which simplicity allows to capture the essence of the benefits of external information-dependent penalties. 
These relate to existing literature, our goal is to characterize precisely the benefits of external information.
Second, we study 
%We then extend our results to 
linear regression under arbitrary design, which is our main contribution. In both settings, we
show that an oracle may take advantage of the external information %to set block-based penalties 
so that variable selection consistency is either attained where otherwise it would not be possible, or is attained at a faster rate. Our analysis highlights in particular how leveraging external information weakens (potentially overly stringent) conditions on signal strength for support recovery.
Finally, we propose empirical Bayes data-analysis  %variable selection 
procedures that realize the theoretical benefits without requiring an oracle
(see \cite{petrone:2014,castillo2020spike} for background on empirical Bayes in Bayesian model selection). 
In our examples, these methods run in seconds using MCMC. Another contribution of independent interest are new tight necessary and sufficient conditions for variable selection consistency in linear regression under arbitrary design and fixed support.
\smallskip

\noindent \textbf{Related work:} Our work has connections with multiple hypothesis testing ideas. In this line of research, \cite{genovese2006} proposed a false discovery rate (FDR) procedure in which p-values are weighted based on prior information. They show that if the weights are positively associated to the null hypotheses being false, their procedure improves power. Subsequent works discussed oracle choices of weights either based on external information or derived from the data \citep{roederwasserman,basu}. Recently, \cite{ramdas2019prior} included prior information in a group-based FDR control. 
Relative to this work, we study variable selection consistency in high-dimensional regression as well as the associated conditions on sparsity and signal strength.

Also related to our work, \cite{scarlett2012compressed} studied compressed sensing when prior information allows splitting parameters into blocks, where one knows the true proportion of non-zero parameters in each block. They show that sparse signal recovery with block-based penalties requires smaller $n$ than with exchangeable penalties. A key difference with our work is their assuming independence across covariates, which renders the results inapplicable to regression. Also, they do not consider the sequence model, nor that the proportions of non-zero parameters are unknown in practice.
\smallskip

\noindent \textbf{Organization:} %The paper is organized as follows. 
Section~\ref{sec:intro_blockpen} introduces block $\ell_0$ penalization. Section~\ref{sec:properties_orth} presents its
model selection properties and benefits in the Gaussian sequence model, in an oracle setting where the true sparsity and betamin conditions are known for all blocks. Section~\ref{sec:properties_legression} studies block $\ell_0$ penalties in linear regression, and shows analogous benefits to those in Section~\ref{sec:properties_orth}. These results can be extended to a wide class of Bayesian variable selection methods. Section~\ref{sec:informed_l0} presents data-based 
procedures, motivated by empirical Bayes, that achieve the improved model selection consistency and rates without requiring an oracle. Section~\ref{sec:simulation} shows empirical examples and Section~\ref{sec:discussion} concludes. Proofs are gathered in the Supplement.
\smallskip

\noindent \textbf{Notation:}  
We denote by $\bs \beta\in \R^p$ the parameters of interest 
and by $\bs \beta^*$ their true values. Let  $V=\{1,\ldots,p\}$. For any $A\subseteq V$ and any vector $\bs x\in \R^p$, $\bs x_A$ denotes the subvector of $\bs x$ with entries corresponding to indices in $A$. For any matrix $\bs X \in \R^{n\times p}$, $\bs X_A$ denotes the submatrix of $\bs X$ obtained by selecting the \emph{columns} with indices in $A$. We denote by
$S=\left\{i\in V:\,\,  \beta^*_i \neq 0 \right\}$ the true support of $\bs \beta^*$, its size is $s=|S|$, and hence $p-s$ is the number of truly inactive parameters.
Denote $\beta_{\min}^*=\min_{i \in S}|\beta^*_i|$ the smallest true signal.
We assume that $V$ is partitioned into a fixed number $b$ of disjoint blocks $B_j \subseteq V$ for $j=1,\ldots,b$.
We denote by $S_j=\left\{i\in B_j:\,\,  \beta^*_i \neq 0\right\}$ the set of active parameters in block $B_j$,  $s_j=|S_j|$ its size, $p_j-s_j$ the number of inactive parameters in the block,
and $\beta_{\min,j}^*=\min_{i \in S_j}|\beta^*_i|$.
Given sequences $f(n) > 0$ and $g(n) >0$, $f(n)=O(g(n))$ means that there exists a constant $c<\infty$ such that $f(n) \leq c g(n)$ for all $n \geq n_0$ and some fixed $n_0$, $f(n)=o(g(n))$ means that $\lim_{n\to\infty}f(n)/g(n)=0$, and $f(n)=\Theta(g(n))$ means that $f(n) = O(g(n))$ and $g(n) = O(f(n))$. For any set $A$, $A^C$ denotes the complement of $A$. \smallskip

\noindent\textbf{Our asymptotic regime:} Although not explicitly denoted, $s_j$, $p_j$ and $\beta_{\min,j}^*$ are functions of $n$, and so are $s$, $p$ and $\beta_{\min}^*$. We study asymptotic regimes where \\[5pt]
     \hypertarget{lab:A1}{(A1)}\quad $n \,\to \, \infty$, and the number of blocks $b$ is constant. \\[5pt]
     \hypertarget{lab:A2}{(A2)}\quad For all $j$, $p_j-s_j \,\to \, \infty$.\\[5pt]
Assumption~\hyperlink{lab:A1}{A1} can be slightly relaxed to allow $b$ to grow slowly with $p$, but we assume a constant $b$ for simplicity.
Similarly, Assumption~\hyperlink{lab:A2}{A2} is a mild assumption that can be relaxed but simplifies the exposition. We provide results for both the case where the $s_j\geq1$ are fixed and $s_j \rightarrow \infty$. We do not make any assumption on the asymptotic regime linking $n$ and $p$, but our main interest is in $n= o(p)$ settings. Assumptions~\hyperlink{lab:A1}{A1}--\hyperlink{lab:A2}{A2} describe the general framework of our results, in each of our results below we specify precisely what %we precise for each result which 
assumptions are needed.
%Assumption~\hyperlink{lab:A2}{A2} implies that $p-s$ goes to infinity but we do not make any assumption on the asymptotic regime linking $n$ and $p$. We provide results for both cases $s_j\geq1$ fixed or diverging.

\section{Variable selection via informed block penalization}\label{sec:intro_blockpen}

Consider a set of candidate models $\M$ given by subsets $M\subseteq V$ and their corresponding coordinate subspaces 
%\david{I'm confused, we just said $\M$ but below we have $\mathcal{L}_M$? Should it be $\M$ everywhere?} \paul{$\M$ is the set of models, $\mathcal{L}_M$ is the parameter space for any $\hat{\beta}_M$: $\hat{\beta}_i$ is 0 if $i\not \in M$, it is any real number if $i\not \in M$}
$$
\mathcal{L}_M\;:=\;\{\beta\in \R^p: \beta_i=0\mbox{ if }i\notin M\}.
$$
Consider a standard $\ell_0$ selection procedure with penalty $\eta(M)$ that depends on $M$ linearly through its cardinality, $\eta(M)=\kappa |M|$ for some $\kappa>0$. The selected model is
\begin{equation}\label{eq:stdl0}
\hat{S}= \arg\max_{M\in \M} \left\{\max_{\beta\in \mathcal{L}_M}\ell(\bs y; \bs \beta)\;-\;\kappa |M|\right\},
\end{equation}
where $\ell(\bs y; \bs \beta)$ is the log-likelihood function. 
We study a externally-informed block procedure that penalizes differently the blocks $B_1,\ldots,B_b$ induced by the external information,
\begin{equation}\label{eq:Mhat}
  \hat S^b\;\in\;\arg\max_{M\in \M} \left\{\max_{\beta\in \mathcal{L}_M}\ell(\bs y; \bs \beta)-\sum_{j=1}^b\kappa_j |M_j| \right\},
\end{equation}
where $\kappa_1>0,\ldots,\kappa_b>0$. We denote a model by $M=M_1\cup\cdots\cup M_b$,
where $M_j \subseteq B_j$ are %and with $M_j:=M\cap B_j$ represents 
the selected parameters in block $j$.
Note that $\hat{S}^1$ is the standard $\ell_0$ selector in \eqref{eq:stdl0}.

\section{Gaussian sequence model}\label{sec:properties_orth}

In this section we discuss the properties of the externally-informed $\hat S^b$ in the sequence model. It is a popular simplified model for high-dimensional inference, and it is also central to non-parametric statistics. See \cite{johnstone}, Chapter 3 for a complete introduction. In our case, the study of the sequence model allows to capture the essence of the benefits of $\hat S^b$ over standard selectors before moving to the setting of interest, linear regression.

We start by introducing the sequence model, its connection to orthogonal linear regression and discuss that $\hat S^b$ simplifies to performing block-based thresholding. %performs block-wise thresholding in that setting. 
We then provide the variable selection properties of $\hat S^b$, discuss its benefits in lessening the conditions for consistent variable selection and in its convergence rates, and illustrate said advantages under different asymptotic regimes.

In this section, we focus on regimes where, in addition to Assumptions~\hyperlink{lab:A1}{A1}--\hyperlink{lab:A2}{A2}, we have\\[5pt]
%we study the selection properties of $\hat{S}^b$ under assumption \\[5pt]
\hypertarget{lab:A3}{(A3)}\quad $\text{For all } j, \;\; s_j \,\to \, \infty$.\\[5pt]
That is, we assume a diverging number of active signals $s_j$ in every block. %Alternatively, the $s_j$'s can be assumed to be finite. 
We present here results for $\hat S^b$ under Assumption~\hyperlink{lab:A3}{A3} to simplify the exposition, and because we derive a novel necessary condition on signal strength under that assumption. We obtained analogous results when the $s_j$'s are finite (section S5 of the Supplement). Note some results in the current section do not require Assumption~\hyperlink{lab:A3}{A3}, we specify in each result which assumptions are needed. 

\subsection{Sequence model and thresholding}\label{sec:GSMgen}

The Gaussian sequence model assumes:
\begin{equation}\label{eq:GSM}
	\bs y \;=\; \sqrt{n} \bs \beta^* + \bs\epsilon,\qquad \bs \epsilon \sim N(0,\sigma^2 I_p),
\end{equation}
where $\bs y \in \mathbb{R}^p$ and without loss of generality we set $\sigma^2=1$ to streamline notation. 

Let $\tilde{\bs \beta}=\bs y/\sqrt{n} \sim N(\bs \beta^*,\tfrac{1}{n}I_p)$ be the MLE under \eqref{eq:GSM}. The next basic result states that %, by orthogonality, 
the block-wise $\ell_0$ penalty reduces to thresholding $\tilde{\bs \beta}$, %$\tilde{\bs \beta}=(\tilde \beta_i)$, 
with block-dependent thresholds.
\begin{prop}\label{prop:isthresh}
	In the sequence model \eqref{eq:GSM}, let $\hat{S}^b$ and $\kappa_1,\ldots,\kappa_b$ defined in \eqref{eq:Mhat}. Then $\hat S^b=\hat S^b_1\cup \cdots \cup \hat S^b_b$, where, for each $j=1,\ldots,b$,
	$$
	\hat S_j^b\;:=\;\left\{i\in B_j:\; |\tilde{\beta}_i|> \sqrt{\frac{2\kappa_j}{n}}\right\}.
	$$ 
\end{prop}
Other popular penalties
%Most popular variable selection procedures 
also take this block-based thresholding form. It is the case of the LASSO and adaptive LASSO when letting the penalty vary by block: see Lemma~S1.1 %in Auxiliary results 
in the Supplement. %Auxiliary results in the Supplement~S0. 
It is also the case of Bayesian procedures under most standard priors \citep{davidomiros} where one sets a different prior inclusion probability in each block. Thus, equivalently, to study the block-informed penalization for the Gaussian sequence model, we study generic thresholding model selectors of the form
\begin{equation}\label{eq:thresholdS}
\hat S_j^{b}\;:=\;\{i\in B_j: |\tilde{\beta}_i|>\tau_j\}, \qquad\bs \tau=(\tau_1,\ldots,\tau_b)\in \R^b_{>0}.	
\end{equation}

Results in this section also generalize to orthogonal linear regression with normalized columns, i.e. when $\bs X$ in \eqref{eq:linearmodel} satisfies $\bs X^\top \bs X=n I_p$. In that setting the MLE is also distributed $N(\bs \beta^*,\tfrac{1}{n}I_p)$ and %by orthogonality, 
variable selection with $\hat{S}^b$ also amounts to block thresholding.

\subsection{Selection based on block thresholds}
\label{ssec:block_thresholds}


We study here the statistical performance of the block thresholding operator in \eqref{eq:thresholdS}. Block thresholding was previously studied in the context of parameter estimation for wavelet-based models with equally sized blocks; see \cite{johnstone}, Chapters 7 to 9.
We focus here on its variable selection properties, for arbitrarily-sized blocks. We successively analyze properties relative to conditions for recovery and rates of convergence. 

Consider the following assumptions,\\
\hypertarget{lab:A4}{(A4)}\quad for all $j=1,\ldots,b$ and every sufficiently large $n$,\;\; $\sqrt{n}\tau_j \;\geq\;  \sqrt{2\ln(p_j-s_j)}$,\\[5pt]
\hypertarget{lab:A5}{(A5)}\quad for all $j=1,\ldots,b$ and 
every sufficiently large $n$, \;
	 $\sqrt{n}(\beta_{\min,j}^* - \tau_j) \;\geq\; \sqrt{2\ln(s_j)}$.

\begin{prop}\label{prop:suffnecblockthresh} In the
sequence model \eqref{eq:GSM}, assume \hyperlink{lab:A1}{A1}, \hyperlink{lab:A2}{A2} and \hyperlink{lab:A3}{A3}.
\begin{enumerate}[label=(\roman*)]
    \item If Assumption~\hyperlink{lab:A4}{A4} holds, then 
$\lim_{n \rightarrow \infty} P(\hat{S}^{b}\subseteq S) = 1$.
\item If for some $j\in\{1,\ldots,b\}$\; $
\lim_{n\to \infty} \sqrt{n}\tau_j/\sqrt{2\ln(p_j-s_j)} \;<\;1$,
then $\lim_{n \rightarrow \infty} P(\hat{S}^{b}\subseteq S ) <1$.
\item If Assumption~\hyperlink{lab:A5}{A5} holds, then
$\lim_{n \rightarrow \infty} P(\hat{S}^{b}\supseteq  S)= 1$.
\item Suppose there exists $j\in \{1,\ldots,b\}$ such that  $\beta^*_i=\beta_{\min,j}^*$ for all $i \in S_j$, $s_j/p_j<1$, and $\lim_{n\to \infty} \sqrt{n}\tau_j/\sqrt{2\ln(p_j-s_j)}\geq1$. If, in addition, $\lim_{n\to \infty} \sqrt{n}(\beta_{\min,j}^*-\tau_j)/\sqrt{(\pi/2)\ln(s_j)}\leq 1$
then $\lim_{n \rightarrow \infty} P(\hat S^{b}\supseteq S ) < 1$.
\end{enumerate}
\end{prop}

Proposition~\ref{prop:suffnecblockthresh} (i) to (iii) extend to $\hat{S}^{b}$ results known for the standard thresholding $\hat{S}$ ($\hat{S}^b$ with $b=1$) in orthogonal linear regression \citep{ slope, buhlmannvandegeerbook, wainwrighthighdimbook} and in the sequence model \citep{johnstone}, Chapter 3. Proposition~\ref{prop:suffnecblockthresh} (iv) gives a new necessary condition on signal strength for support recovery in probability in a worst case in which all signals are equal. It shows that in that case sufficient Assumption~\hyperlink{lab:A5}{A5} is essentially necessary up to a constant factor (replacing 2 by $\pi/2$).

% By Proposition~\ref{prop:suffnecIblockthresh},
% a necessary assumption to have no false positives asymptotically
% is
% \begin{equation}\label{cond:necIallblockthresh}
% \lim_{n\to \infty} \frac{\tau_j}{\sqrt{\tfrac{2\ln(p_j-s_j)}{n}}}\;\geq\;1\qquad\mbox{for all }j=1,\ldots,b.	
% \end{equation}

% \bigskip

% \noindent\textbf{False negatives.} 
% False negatives occur when $\hat S^{b}\not\supseteq S$, i.e. for some block $j \in \{1,\ldots,b\}$ and parameter $i\in B_j$,  we have that $\beta^*_i\neq 0$ and $|\tilde\beta_i|\leq \tau_j$.
% Control of false negatives in a block $B_j$ requires bounding the minimum of $s_j=|S_j|$ independent Gaussians with non-zero means $\beta^*_i$, $i\in S_j$. This can be achieved by setting a betamin condition
% on $\beta_{\min,j}^*\;:=\;\min\{|\beta^*_i|: i\in S_j\}$.  
% Intuitively, to have no false negatives (also known as screening property), each $\tau_j$ needs to be smaller than $\beta^*_{\min,j}$.
% More precisely, a sufficient assumption is: \\[5pt]


% \begin{prop}\label{prop:suffnecIIblockthresh}
%  In the sequence model \eqref{eq:GSM}, assume \hyperlink{lab:A1}{A1}--\hyperlink{lab:A3}{A3} and \hyperlink{lab:A5}{A5}. Then
% $\lim_{n \rightarrow \infty} P(\hat{S}^{b}\supseteq  S)= 1$.
% \end{prop}
% Assumption~\hyperlink{lab:A5}{A5} was proved to be sufficient for selection consistency in orthogonal linear regression in \cite{wainwrighthighdimbook} for standard thresholding. Proposition~\ref{prop:suffnecIIblockthresh} generalizes the result to block-thresholding. The study of {\it necessary} assumptions
% is more complicated. Proposition~\ref{prop:necIIblockthresh} below provides first a necessary condition valid for any vector $\bs \beta^*$, and second, a tight condition for a worst-case scenario where $\beta_i^*=\beta_{\min}^*$ for all $i\in S$.


% \bigskip

% \noindent\textbf{Exact model recovery guarantees.} 
Combining Assumption~\hyperlink{lab:A4}{A4} (Proposition~\ref{prop:suffnecblockthresh} (i)) and Assumption~\hyperlink{lab:A5}{A5} (Proposition~\ref{prop:suffnecblockthresh} (iii)), we asymptotically recover the correct support if for all $j=1,\ldots,b$,
\begin{equation*}
  \sqrt{\frac{2\ln(p_j-s_j)}{n}}\;\;\leq\;\;\tau_j\;\;\leq\;\;\beta_{\min,j}^*-\sqrt{\frac{2\ln(s_j)}{n}}
\end{equation*}
for every sufficiently large $n$. In particular, this requires that for all $j=1,\ldots,b$,
\begin{equation}\label{eq:betamin}
    \beta_{\min,j}^* \geq \sqrt{\frac{2\ln(p_j-s_j)}{n}} + \sqrt{\frac{2\ln(s_j)}{n}}.
\end{equation}
% By Propositions~\ref{prop:suffnecIblockthresh} and~\ref{prop:suffnecIIblockthresh} signals as small as the right-hand side in \eqref{eq:betamin} can be recovered in block $B_j$. 
Proposition~\ref{prop:suffnecblockthresh} (ii) and (iv) show that Assumptions~\hyperlink{lab:A4}{A4}--\hyperlink{lab:A5}{A5} are essentially necessary. It follows that betamin condition \eqref{eq:betamin} is near-necessary for selection consistency, in the following sense.
\begin{lem}\label{lem:necbetaminblockthresh}
    If there exists a block $j\in\{1,\ldots,b\}$ with equal non-zero parameter values $\beta^*_i=\beta_{\min,j}^*$ for all $i \in S_j$, where $0<s_j/p_j<1$ and 
    \begin{equation}\label{eq:necbetaminblockthresh}
        \lim_{n\to \infty}\frac{\sqrt{n}\beta_{\min,j}^*}{\sqrt{2\ln(p_j-s_j)}+\sqrt{(\pi/2)\ln(s_j)}}\;<\; 1,
    \end{equation}
    then under Assumptions~\hyperlink{lab:A1}{A1}--\hyperlink{lab:A3}{A3},\;  $\lim_{n \rightarrow \infty} P(\hat{S}^{b} = S) < 1$.
\end{lem}
For standard thresholding $\hat{S}$, \cite{butucea2018} showed that the complementary of \eqref{eq:necbetaminblockthresh} where in the denominator $\pi/2$ is replaced by $2$ is strictly necessary for a stricter definition of selection consistency, the vanishing in expectation of the Hamming loss (the number of false negatives and positives). 

\smallskip

\noindent\textbf{Rate of convergence.} Building upon the
preceding results, we bound the rate of convergence of $P(\hat S^{b}\neq S)$.
\begin{thm}\label{theo:fullrecovGS}
Assume \hyperlink{lab:A4}{A4} and \hyperlink{lab:A5}{A5} and that for all $j=1,\ldots,b$, $p_j-s_j>1$ and $s_j>1$. %\david{Clarify: also \hyperlink{lab:A1}{A1}-\hyperlink{lab:A3}{A3}?} \paul{No, \hyperlink{lab:A1}{A1}-\hyperlink{lab:A3}{A3} are not required for the bound to hold} 
Then, for every sufficiently large $n$,
\begin{equation} \label{eq:rateconvgsm}
P(\hat{S}^{b} \neq S)\;\leq\;\sum_{j=1}^b\,\,\frac{e^{-\tfrac{n}{2}\big[\tau_j^2-\tfrac{2\ln(p_j-s_j)}{n}\big]}}{\sqrt{\pi\ln(p_j-s_j)}}+\frac{e^{-\tfrac{n}{2}\big[(\beta^*_{\min,j}-\tau_j)^2-\tfrac{2\ln(s_j)}{n}\big]}}{\sqrt{\pi\ln(s_j)}}.
\end{equation}
Further,
if $\beta^*_{\min,j}$ satisfy \eqref{eq:betamin} for every block $j$, and the thresholds $\tau_j$ take oracle values
	\begin{equation} \label{eq:optimtauj}
\tau^*_j \;=\; \frac{\beta^*_{\min,j}}{2}+\frac{\ln({p_j/s_j}-1)}{n\beta^*_{\min,j}},
\end{equation}
then $\tau^*_j$ satisfy Assumptions~\hyperlink{lab:A4}{A4} and \hyperlink{lab:A5}{A5} and, for every sufficiently large $n$,
\begin{equation}\label{eq:optimrateconvgsm}
    P(\hat{S}^{b} \neq S) \; \leq \; 2 \sum_{j=1}^be^{-[\frac{n}{8}{\beta^*_{\min,j}}^2-\ln \max\{p_j-s_j,s_j\}]}.
\end{equation}
\end{thm}
In \eqref{eq:rateconvgsm} the first term of each summand decreases exponentially in $n\tau_j^2/2$, while the second term decreases exponentially in $n(\beta_{\min,j}^*-\tau_j)^2/2$. The choice $\tau_j=\tau^*_j$ ensures that both terms are equal and hence approximates the values minimizing \eqref{eq:rateconvgsm}. We refer to these ideal $\tau^*_j$ as {\it oracle thresholds} because they depend on quantities $s_j$ and $\beta_{\min,j}$ that are unknown in practice. 
The bound in \eqref{eq:optimrateconvgsm} closely approximates \eqref{eq:rateconvgsm} for $\tau_j=\tau^*_j$ and is tightest in the worst-case scenario where $\beta_i^*=\beta_{\min,j}^*$ for all $i \in S_j$ and all $j$, in that it approximates the fastest rate achievable in that worst case. In fact, for standard thresholding $\hat{S}$ ($\hat{S}^b$ with $b=1$), Corollary 2.1 of \cite{butucea2018} shows that the choice of threshold $\tau^*=\beta^*_{\min}/2+\ln({p/s-1})/(n\beta^*_{\min})$ is minimax for the Hamming loss up to a constant factor smaller than 2. By independence, it is straightforward that $\tau_j^*$ is minimax, up to a factor 2, for the Hamming loss \textit{in block $B_j$}.

\subsection{Benefits of block thresholds}\label{sec:prosconsblock}

We now examine the benefits of block thresholds. We discuss
two types of benefits: softening the 
conditions for model selection consistency and improving the associated convergence rates. \smallskip

% Intuitively, block penalties allow adapting 
% to varying sparsity $s_j$  or signal strengths $\beta^*_{\min,j}$ across blocks. 
% When the number of inactive parameters $p_j-s_j$ is small one may use smaller thresholds $\tau_j$ and still control false positives, and when $\beta^*_{\min,j}$ is large one may use larger $\tau_j$ and still control false negatives.
% As the blocks become more discriminative, that is, more distinct $p_j-s_j$ and $\beta^*_{\min,j}$ across blocks, the larger the potential benefits.
% Consider a somewhat extreme example with $b=2$ blocks, a large sparse block with $s_1=0$, $p_1-s_1=80$ and a smaller non-sparse block with $s_2=10$, $p_2-s_2=10$, with $n=1$ and $\beta_i^*=4$ for all $i\in S$. 
% In the left panel of Figure~\ref{fig:starter}, a standard threshold selector is used with $\tau=\sqrt{2 \ln(p-s)/ n}$, the minimal threshold such that the family-wise error rate vanishes (dashed line). The large number of truly zero parameters $p-s=90$ forces a high threshold, which leads to a false negative. In the right panel, block thresholds are set to $\tau_j=\sqrt{2 \ln(p_j-s_j)/ n}$.
% Adjusting $\tau_2$ to the lower number of truly zero means in block 2 ($p_2-s_2$) allows recovering the true model. 
% \begin{figure}[ht]
% \begin{minipage}{1\columnwidth}
% \begin{tabular}{M{62mm}M{62mm}}
%         Standard thresholding & Block thresholding \\
%          \includegraphics[scale=1]{images/p.intro.std.pdf} & \includegraphics[scale=1]{images/p.intro.block.pdf} \\
% \end{tabular}
% \end{minipage}
% \caption{In both panels, the 90 white dots are generated from $N(0, 1)$ and the 10 black dots are generated from $N(4, 1)$. Left panel: the dashed line is $\sqrt{2 \ln (p-s)/n}$, the minimal threshold such that the FWER vanishes. Right panel: observations are grouped in two blocks of sizes $p_1=80$ and $p_2=20$. The dashed lines are the minimal block thresholds such that the FWER vanishes}
% \label{fig:starter}
% \end{figure}
% \bigskip

\noindent\textbf{Conditions for variable selection consistency.} 
% In the example, the possibility to lower thresholds in the less sparse blocks %with $\hat S^2$ 
% allowed recovering the smallest active signal. 
% More precisely, 
Assumptions~\hyperlink{lab:A4}{A4}--\hyperlink{lab:A5}{A5} give ranges of threshold values that are sufficient and essentially necessary for asymptotic support recovery. For the standard selector $\hat{S}$, the range for the single threshold $\tau$ is
\begin{equation}\label{eq:rangethreshex}
  \sqrt{\frac{2\ln(p-s)}{n}}\;\;\leq\;\;\tau\;\;\leq\;\;\beta^*_{\min}-\sqrt{\frac{2\ln(s)}{n}}.	
\end{equation}
For the block threshold selector $\hat{S}^{b}$, 
%assuming without loss of generality that the global minimum $\beta_{\min}^*$ is in block $b$, 
the ranges for the $\tau_j$'s are
    \begin{eqnarray}\label{eq:rangeblockthreshex}
  \sqrt{\frac{2\ln(p_j-s_j)}{n}}\;\;&\leq\;\;\tau_j\;\;\leq&\;\;\beta^*_{\min,j}-\sqrt{\frac{2\ln(s_j)}{n}}   .
%  \sqrt{\frac{2\ln(p_1-s_1)}{n}}\;\;&\leq\;\;\tau_1\;\;\leq&\;\;\beta^*_{\min,1}-\sqrt{\frac{2\ln(s_1)}{n}}\nonumber\\
%  &\;\;\vdots \\
%  \sqrt{\frac{2\ln(p_b-s_b)}{n}}\;\;&\leq\;\;\tau_b\;\;\leq&\;\;\beta^*_{\min}-\sqrt{\frac{2\ln(s_b)}{n}} \nonumber.       
    \end{eqnarray}

The ranges in \eqref{eq:rangeblockthreshex} imply that $\hat{S}^{b}$ requires milder conditions for selection consistency than $\hat{S}$. Intuitively, if there exist two blocks such that the ranges in \eqref{eq:rangeblockthreshex} do not overlap, 
then a global threshold $\tau$ cannot possibly satisfy \eqref{eq:rangeblockthreshex} for all $j$ and consistent selection is essentially not possible.
For example, this occurs if the global smallest active signal $\beta_{\min}^*$ is in block $b$ and satisfies  $\beta_{\min}^* - \sqrt{(2/n) \ln(s_b)} < \sqrt{(2/n) \ln(p_1-s_1)}$.
More precisely, Corollary~\ref{cor:selconstimpossible} gives
conditions under which consistent selection is possible with $\hat{S}^b$ but not with $\hat{S}$, in a worst-case setting where all non-zero parameters are equal to $\beta_{\min}^*$.


\begin{cor}\label{cor:selconstimpossible}
Assume \hyperlink{lab:A1}{A1}--\hyperlink{lab:A5}{A5}, that $s/p<1$, and $\beta^*_i=\beta_{\min}^*$ for all $i \in S$. If
\begin{eqnarray*}
    &\lim_{n\to \infty}\frac{\beta_{\min}^*}{\sqrt{\tfrac{2\ln(p-s)}{n}}+\sqrt{\tfrac{\pi}{2}\tfrac{\ln(s)}{n}}}\;<\; 1 
\end{eqnarray*}
then 
$\lim_{n \rightarrow \infty} P(\hat{S} = S) < 1$ and $\lim_{n \rightarrow \infty} P(\hat{S}^{b} = S)=1$.
\end{cor}

\noindent\textbf{Rates of convergence.} We just saw that block penalties can attain consistent model selection where standard penalties cannot. We now discuss differences in the probability of correct selection, when consistent selection is possible for both procedures. We assume that block thresholds 
are set to their oracle values $\tau_j^*$ in \eqref{eq:optimtauj} 
and similarly that the single threshold
takes its oracle value $\tau^*$ defined in an analogous way.
Let ${OR}^b_{orth}$ be the oracle convergence rate for $\tau_j^*$ in \eqref{eq:optimrateconvgsm} and ${OR}_{orth}$ that for $\tau^*$.
Then
\begin{equation}\label{eq:ratioCR}
    \frac{{OR}^b_{orth}}{{OR}_{orth}}\;=\;\sum_{j=1}^be^{-\frac{n}{8}\left({\beta^*_{\min,j}}^2-{\beta_{\min}^*}^2\right)-(\ln \max\{p-s,s\}-\ln \max\{p_j-s_j,s_j\})}
\end{equation}
Since $\beta_{\min,j}^* \geq \beta_{\min}^*$ and $\ln \max\{p-s,s\}\geq\ln \max\{p_j-s_j,s_j\}\big)$, equation \eqref{eq:ratioCR} shows that for every $n$ large enough ${OR}^b_{orth} \leq {OR}_{orth}$ for any partition in blocks, i.e. the oracle convergence rate in \eqref{eq:optimrateconvgsm} for $\tau_j^*$ is never worse than for $\tau^*$. The magnitude of the gain depends on how informative the blocks are. %the discriminative power of the blocks.
For any sparse setting where $s_j<p_j-s_j$ for every $j$, assuming without loss of generality that $\beta_{\min,b}^*=\beta_{\min}^*$, we have
\begin{equation}\label{eq:goodorthratio}
    \frac{{OR}^b_{orth}}{{OR}_{orth}}\;=\; \frac{p_b-s_b}{p-s} \,+\,\sum_{j<b}\frac{p_j-s_j}{p-s}e^{-\frac{n}{8}\left({\beta^*_{\min,j}}^2-{\beta_{\min}^*}^2\right)}
\end{equation}
%\paul{this ratio goes to 1 if the betamins are equal across blocks and for some $j$, $p_j-s_j/p-s\to1$ (happens in our Examples). I guess this is due to simplifications in the bounding.}
There are then two sources of improvement in convergence rates: in any block $j$, a smaller number of truly zero means, that is $p_j-s_j< p-s$, or a larger smallest signal, $\beta^*_{\min,j} > \beta_{\min}^*$. Depending on the sparsity of the block, one of the two effects usually prevails. In blocks where $\beta^*_{\min,j} \approx \beta_{\min}^*$, improvement comes mainly from having a smaller number of truly zero means. In those blocks
one sets $\tau^*_j < \tau^*$, so that one may detect smaller signals. 
If
$\beta^*_{\min,j} \gg \beta_{\min}^*$,
one sets $\tau^*_j > \tau^*$
and the probability of false positives is reduced.

In the particular worst case where all the active signals are equal, $\beta^*_{\min,j}=\beta^*_{\min}$ for all $j$, the ratio 
${OR}^b_{orth}/{OR}_{orth}$ in \eqref{eq:goodorthratio} is exactly one. That is, the oracle rate for the $\tau_j^*$ is not better than for $\tau^*$. This is in line with the minimax analysis of \cite{butucea2018} which shows the near-optimality of $\tau^*$ in the worst case. Under the less rigid assumption of unequal $\beta^*_{\min,j}$'s, ${OR}^b_{orth}/{OR}_{orth}$ is however bounded away from 1 in \eqref{eq:goodorthratio} and the oracle rate for the $\tau_j^*$ is strictly better. This highlights the crucial role that varying signal strength plays in the gains with informed thresholds. The latter cannot be only characterized by the entropy of the distribution of sparsity across blocks as in \cite{scarlett2012compressed}. 
% It is shown that $\tau^*$ is near-minimax over a class of so-called separable selectors that do contain block thresholds. However, the minimized maximum risk is that . We observe indeed that in that case , meaning the $\tau^*_j$ do not outperform $\tau^*$. Our analysis suggests that under the more realistic assumption of unequal active signals, thresholds other than $\tau^*$, like the informed $\tau_j^*$, are preferable. The case where all active signals are equal to $\beta^*_{\min}$ also highlights how a non-uniform distribution of sparsity across blocks alone is insufficient to guarantee improvement in rates of convergence.
% That is, one does not expect benefits when the number of non-zero signals and their magnitudes are similar across blocks. This can be also seen from \eqref{eq:goodorthratio}.

\subsection{Examples with two blocks}\label{sec:examples}

We illustrate our results with four concrete examples contemplating different sparsity regimes and informativeness of the blocks, summarized in Table~\ref{tab:examples}. We focus on a sparse setting where $\ln(s)=o(\ln(p-s))$, with two blocks ($b=2$). We assume that $\beta_{\min}^*$, the smallest non-zero $|\beta_i^*|$ is located in block $2$. 

\begin{table*}[ht]
\caption{Sparsity and block assumptions in Examples 1--4}
\label{tab:examples}
\begin{tabular}{@{}cccccccc@{}}
\hline \\[-1em]
Example & $p-s$ & $p_1-s_1$ & $p_2-s_2$ & $s_1$ 
& $s_2$\\[1pt]
\hline\\[-1em]
$1$    & $3n/2$ & $3n/2-\sqrt{n}$ & $\sqrt{n}$  & $3\ln(n)/2$  & $3\ln(n)/2$ \\[1pt]
$2$    & $e^{n/20}$ & $e^{n/20}-n^2$ & $n^2$  & $3\ln(n)/2$  & $3\ln(n)/2$   \\[1pt]
$3$    & $n$ & $n-\ln(n)$ & $\ln(n)$  & $3\ln(n)/2$  & $3\ln(n)/2$   \\[1pt]
$4$    & $n$ & $n/2$ & $n/2$  & $3\ln(n)/2$  & $3\ln(n)/2$ \\[3pt]
\hline
\end{tabular}
\end{table*}

In Examples 1 to 3, external information is discriminative as it singles out a block $B_1$ that is
sparser than block $B_2$. 
In Example 1 the blocks are 
moderately discriminative, in that $(p_1-s_1)/(p_2-s_2)$ is a power of $n$. In Example 2 they are highly discriminative, since $(p_1-s_1)/(p_2-s_2)$ is exponential in $n$ and in Example 3 they are highly discriminative in that $\bs\beta^*$ is sparse overall but it is non-sparse within block 2 ($s_2 > p_2 - s_2)$. %but $\bs\beta^*_{B_2}$ is non-sparse ($\ln(p_2-s_2)=O(\ln(s_2))$ and $s_2>p_2-s_2$). 
Example 3 also differs from 1 and 2 in that it is less sparse overall. %set in only very mildly sparse settings. 
Example 4 is a non-discriminative
%In Example 4, external information is non-discriminative, the block partition is an average
random guess where each block has half of the inactive parameters. %variables. 
%\paul{It is not super important but I think that what matters in terms of usefulness of the external information is $p-s/p_2-s_2$ but David has been stressing $p_1-s_1/p_2-s_2$ in his edits. My point of view is that one the main gains with $\hat S^b$ is that you can recover smaller signals in the blocks that are less sparse than the global set of variables. If you have a three block situation such: $p-s=e^n$, $p_1-s_1=e^n-n^2-n$, $p_2-s_2=n^2$ and $p_3-s_3=n$, $p_2-s_2/p_3-s_3$ is not telling of the gain with $\hat S^b$. To illustrate how much you can \emph{increase} penalties in some blocks, to me, the relevant ratio would then be $\beta_{\min,j}/\beta_{\min}$. }
%\david{Can we facilitate following why one example is more discriminative than another? To the reader it's not obvious why one should look at $(p-s)/(p_2-s_2)$ to see which example is more discriminative, as opposed to $s_1/p_1$ vs. $s_2/p_2$ (which would be easier to follow) or $(p_1-s_1)/(p_2-s_2)$ (also easier to follow). If we really need to express it in terms of $(p-s)/(p_2-s_2)$ then ok, let's leave it as is -- definitely don't add any further explanation on why we use $(p-s)/(p_2-s_2)$ because that would just make the example harder to read, I think.}

%We then illustrate the benefit of block penalties in terms of rate of convergence with a plot of the ratio ${OR}^b_{orth}/{OR}_{orth}$.

\smallskip
\noindent\textbf{Selection consistency.}
%\noindent\textbf{Valid thresholds for $\hat S$ and $\hat S^b$.}
%To illustrate the benefits in %terms of conditions for selection consistency, 
Figure \ref{fig:boundthresh} plots the range of threshold values $\tau$ and $(\tau_1,\tau_2)$ ensuring selection consistency with $\hat{S}$ and $\hat{S}^{b}$ given in \eqref{eq:rangethreshex} and \eqref{eq:rangeblockthreshex},
%Figure \ref{fig:boundthresh} plots the upper and lower bounds for $\tau$ as in \eqref{eq:rangethreshex} and $\tau_1$ and $\tau_2$ as in \eqref{eq:rangeblockthreshex} 
in the four examples assuming $\beta^*_{\min,1}=2/3$  and $\beta^*_{\min,2}=\beta^*_{\min}=1/10$. 
In Example 2, asymptotic recovery with $\hat{S}$ is not possible in the range of values of $n$ considered while it is with $\hat{S}^{b}$.
In the other examples, asymptotic recovery with $\hat{S}$ is possible but it requires larger $n$ than with $\hat{S}^b$.
Example 1 shows that the gain can be large even when blocks are moderately discriminative, and Example 3 when some blocks are non-sparse.
%Observe that, for $\tau$, the sparser $\bs \beta^*$ is, the larger the $n$ needed to make the upper bound to exceed the lower bound, i.e. to essentially make support recovery possible. 
%The block threshold $\tau_2$ requires smaller $n$, with the gap with respect to $\tau$ increasing as the difference in sparsity between block 2 and $V$ increases. 
The gain in terms of the value of $n$ making recovery possible  is close to null in Example 4 when external information is non-discriminative. 
%In Example 2 on the contrary, asymptotic recovery with $\hat{S}$ is not possible in the range of values of $n$ considered while it is with $\hat{S}^{b}$. Indeed, in that example, $\sqrt{\frac{2\ln(p-s)}{n}}+\sqrt{\frac{\ln(s)}{n}} \to \sqrt{1/10} \approx 0.32$ and $\sqrt{\frac{2\ln(p_2-s_2)}{n}}+\sqrt{\frac{\ln(s_2)}{n}} \to 0$. 
%Example 3 shows that the gain with $\hat{S}^{b}$ can be large even in settings where $V$ is only somewhat sparse.

\begin{figure}[ht]
\begin{minipage}{1\columnwidth}
\setlength{\tabcolsep}{1pt}
\begin{tabular}{M{3mm}M{33mm}M{33mm}M{33mm}M{33mm}}
       \toprule
         & Example 1 & Example 2 & Example 3 & Example 4 \\
        \midrule
        $\tau$ & \includegraphics[scale=1]{images/taus/p.scen1tau.pdf} & \includegraphics[scale=1]{images/taus/p.scen2tau.pdf} & \includegraphics[scale=1]{images/taus/p.scen3tau.pdf} & \includegraphics[scale=1]{images/taus/p.scen4tau.pdf} \\
        $\tau_1$ & \includegraphics[scale=1]{images/taus/p.scen1tau1.pdf} & \includegraphics[scale=1]{images/taus/p.scen2tau1.pdf} & \includegraphics[scale=1]{images/taus/p.scen3tau1.pdf} & \includegraphics[scale=1]{images/taus/p.scen4tau1.pdf} \\
        $\tau_2$ & \includegraphics[scale=1]{images/taus/p.scen1tau2.pdf} & \includegraphics[scale=1]{images/taus/p.scen2tau2.pdf} & \includegraphics[scale=1]{images/taus/p.scen3tau2.pdf} & \includegraphics[scale=1]{images/taus/p.scen4tau2.pdf} \\
        \bottomrule
\end{tabular}
\setlength{\tabcolsep}{6pt}
\end{minipage}
    \caption{%Valid upper bound (solid) and lower bound (dashed) on $\tau$, $\tau_1$ and $\tau_2$ in Examples 1 to 4.
    Smallest (dashed) and largest (solid) value of $\tau$ leading to consistent model recovery in Examples 1 to 4, as given in \eqref{eq:rangethreshex} and \eqref{eq:rangeblockthreshex}.
    Red indicates settings where the interval is empty}
 \label{fig:boundthresh}
\end{figure}
\smallskip
\noindent\textbf{Smallest signal recoverable.} 
We compare the \emph{smallest signal recoverable} by $\hat S$ and $\hat S^b$ while still being selection consistent. We denote those by $\underline{\beta}_{\min,orth}^{*}$ and $\underline{\beta}_{\min,b}^{*}$ respectively. In Assumption~\hyperlink{lab:A5}{A5}, we set the threshold(s) to the lowest value such that the family-wise error rate (FWER) vanishes (per Proposition~\ref{prop:suffnecblockthresh} (i)). Since we assumed that the global minimum $\beta_{\min}^*$ is in block $2$, we get
\begin{equation}\label{eq:betaminbardef}
    \underline{\beta}_{\min,orth}^{*} := \sqrt{\frac{2\ln(p-s)}{n}} + \sqrt{\frac{2\ln(s)}{n}},\qquad
    \underline{\beta}^{*,b}_{\min,orth} := \sqrt{\frac{2\ln(p_2-s_2)}{n}} + \sqrt{\frac{2\ln(s_2)}{n}}
\end{equation} 
%We discuss the ratio $\underline{\beta}^{*,b}_{\min,orth}/\underline{\beta}_{\min,orth}^{*}$ in our four examples. 

The left panel in Figure~\ref{fig:ratios} plots the ratio %of smallest signal recoverable 
$\underline{\beta}^{*,b}_{\min,orth}/\underline{\beta}^{*}_{\min,orth}$. %in our four examples. 
In Example 1, with discriminative blocks, the
smallest signal recoverable with $\hat{S}^{b}$ is asymptotically about 25\% smaller than with the standard selector $\hat{S}$. 
In Example 2 where blocks are even more discriminative, the ratio converges to 0. %: the smallest signal recoverable with $\hat S^b$ is infinitely smaller.
In Example 4, with non-discriminative blocks, the ratio converges to 1, i.e. the benefits in the recoverable signal fade as $n$ grows. 
Example 3 illustrates how highly discriminative blocks can also bring significant benefits in terms of signal recoverable in a regime that is only somewhat sparse.
\begin{figure}[ht]
\begin{minipage}{1\columnwidth}
\begin{tabular}{M{63mm}M{63mm}}
        $\underline{\beta}^{*,b}_{\min,orth}/\underline{\beta}^{*}_{\min,orth}$ & ${OR}^b_{orth}/{OR}_{orth}$ \\
         \includegraphics[scale=1]{images/p.ratio.smallestbeta.pdf} & \includegraphics[scale=1]{images/p.ratio.orconvrate.pdf} \\
\end{tabular}
\end{minipage}
\caption{Ratio of smallest signals recoverable (left) and oracle convergence rates (right) with $\hat S^b$ and with $\hat S$ in Examples 1--4}
\label{fig:ratios}
\end{figure}

\smallskip
\noindent\textbf{Oracle convergence rates.} In the right panel of Figure~\ref{fig:ratios}, we plot the ratio ${OR}^b_{orth}/{OR}_{orth}$ in Examples 1 to 3 where blocks are discriminative. We set $\beta_{\min }^*=\beta_{\min,2}^*=1/5$ and $\beta_{\min,1}^*=1.3\beta_{\min }^*$ to guarantee the recovery is possible with both $\hat S$ and $\hat S^b$. Figure~\ref{fig:ratios} shows that in the three examples the convergence is much faster with $\hat S^b$. %The more discriminative the blocks the faster.


\section{High-dimensional linear regression}
\label{sec:properties_legression}

In Section~\ref{sec:properties_orth} we showed that, in the sequence model, block penalties allow consistent model recovery in settings where it is otherwise not possible (e.g. smaller signals), and improves oracle consistency rates. We now extend the results to linear regression. 
We present sufficient and necessary assumptions, 
betamin conditions and rates for the probability of correct selection for the block informed $\ell_0$-penalized selector $\hat S^b$ in \eqref{eq:Mhat}. We also compare the properties of $\hat S^b$ to those of the standard $\ell_0$ selector $\hat S$. 
Whereas the framework is more involved than in the sequence model and the required proof techniques are different, the results remain conceptually the same. 

%Here $V$ is the set of all variables and $p_j-s_j$ and $s_j$ are respectively the numbers of inactive and active variables in block $B_j$.
%\paul{I removed a sentence that said here "Here $V$ is the set of all variables and $p_j-s_j$ and $s_j$ are respectively the numbers of inactive and active variables in block $B_j$." I don't see why we need to repeat it.} 
We require Assumption~\hyperlink{lab:A1}{A1} across the section, whereas we use \hyperlink{lab:A2}{A2} only in results on necessary conditions for variable selection consistency with $\hat S^b$. Our results on sufficient conditions for selection consistency then hold equally for fixed and diverging $p_j-s_j$ and $s_j$. For simplicity, we assume that $\hat{S}^b$ is unique and that one constrains attention to the set of models $\M\subseteq \mathcal P(V)$ (the power set of $V$),
such that
$\bs X_M$ has full column rank for any $M \in \M$, and that the true support $S$ lies in $\M$. Non-uniqueness of $\hat{S}^b$ and non-full rank models can be accommodated though, at the expense of a slightly more involved treatment.

Our analysis relies on a classical connection
between $\ell_0$ penalties and Bayesian variable selection. 
Section \ref{ssec:regression_proof_strategy} reviews this connection and the proof strategy for our results.
The technical nature of the proof precludes a detailed exposition, we instead only present the most important ideas.
In Section \ref{ssec:regression_mainresult} we state our main theorem
on sufficient conditions for variable selection consistency for $\hat{S}^{b}$ and oracle convergence rates. %as we did for the sequence model. 
To assess the tightness of said sufficient conditions, Section~\ref{ssec:neccond_linreg} gives related {\it necessary} conditions.  %assumptions and betamin conditions for variable selection consistency with $\hat{S}^{b}$. 
Section \ref{ssec:regression_gains} discusses the gains of block penalization in further detail. Section \ref{ssec:nobetaminhyp} gives a general convergence result that, when certain betamin conditions do not hold, one still has guarantees of discarding inactive parameters and detecting sufficiently large active parameters.
%for $\hat{S}^{b}$ under no betamin condition which is used in Section \ref{sec:informed_l0}.
The latter result plays an important role for our data-based procedures in Section \ref{sec:informed_l0}. 


\subsection{Proof strategy}
\label{ssec:regression_proof_strategy}


We state a well-known
reformulation of $\hat S^b$ in linear regression \eqref{eq:linearmodel}. For any model $M\in \M$, let $\tilde{\bs \beta}^{(M)}= (\bs X_{\!\!M}^\top \bs X_{\!\!M})^{-1}\bs X_{\!\!M}^\top\bs y \in \R^{|M|}$ be the MLE under model $M$, and denote
\begin{equation}\label{eq:CM}
    {\rm C}(M)\;: =\;\tfrac{1}{2}\|\bs X_M\tilde{\bs\beta}^{(M)}\|^2-\sum_{j=1}^b\kappa_j |M_j|\quad\text{and}\quad NC(M)\;:=\;\frac{e^{{\rm C}(M)}}{\underset{M' \in \M}{\sum} e^{{\rm C}(M')}}.
\end{equation}
\begin{lem}\label{lem:useful-reformulation} $\hat S^b$ satisfies $\hat S^b\;\in\;\arg\max_{M\in \M} NC(M)$
\end{lem}
By Lemma~\ref{lem:useful-reformulation}, $\hat S^b$ selects
$M\in \M$ with the largest normalized score $NC(M)$. This normalized score can be understood as a pseudo-posterior probability for model $M$ in a Bayesian variable selection framework \citep{BIC}.
Proposition 1 in \cite{rossell2022concentration}, reproduced below as Lemma~\ref{lem:L0toL1convergence}, proves that the expectation of such posterior probabilities bounds the probability of an incorrect model selection $(\hat S^b\neq S)$. 

\begin{lem}\label{lem:L0toL1convergence}
\begin{enumerate}[label=(\roman*)]
\item 
$P(\hat{S}^{b} \neq S) \,\,\leq \,\, 2 \sum_{M \in \M\setminus \{S\}} \E\left(NC(M)\right).$
\item  For any $M,M'\in\M$, such that $M\neq M'$, 
$NC(M)\;\leq \;\big(1+e^{C(M')-C(M)}\big)^{-1}$.
\end{enumerate}
\end{lem}
We use  Lemma~\ref{lem:L0toL1convergence} (i) to show the variable selection consistency of $\hat S^b$ by guaranteeing the vanishing, in expectation, of the sum of the normalized scores. Such vanishing bears resemblance to what is known as \emph{strong selection consistency} property \citep{narisettyhe2014}: the concentration of pseudo-posterior model probabilities on the truth $S$. We prove however $L_1$ convergence instead of convergence in probability.
To bound $\E(NC(M))$ for $M\neq S$, by Lemma~\ref{lem:L0toL1convergence} (ii) with $M'=S$, we may instead bound the expectation of a simple function of $C(S)-C(M)$, the pairwise comparison between each $M$ and the data-generating $S$.
This is achieved by noting that, directly by the definition in \eqref{eq:CM}, $C(S)-C(M)=\tfrac12 L_{SM}+\Delta_{MS}$, where for any two models $M,T\subseteq V$, we denote
\begin{equation}\label{eq:DeltaM}
\Delta_{MT}\;:=\;\sum_{j=1}^b \kappa_j(|M_j|-|T_j|) \quad\text{and} \quad L_{TM}\;:=\;\|\bs X_{T}\tilde{\bs\beta}^{(T)}\|^2-\|\bs X_{M}\tilde{\bs\beta}^{(M)}\|^2.
\end{equation}
%and for any two models $Q,M\subseteq V$, we denote 
%$$L_{QM}\;:=\;\|\bs X_{Q}\tilde{\bs\beta}^{(Q)}\|^2-\|\bs X_{M}\tilde{\bs\beta}^{(M)}\|^2.$$
To lower bound $C(S)-C(M)$ we use that $L_{SM}$ can be expressed in terms of chi-squared variables.
%For any over-fitted model $M \supset S$, we have $\Delta_{MS}>0$ and it is well-known that $-L_{SM} \sim \chi_{|M \setminus S|}^2$. Hence, if one sets large enough $\kappa_j$ (and thus $\Delta_{MS}$), then $NC(M)$ vanishes. For non-overfitted models $M\not\supset S$, 
The idea is to take the union model $Q_S=S\cup M$, and to note that $L_{SM}=L_{Q_SM}-L_{Q_SS}$. We may use the next lemma from \cite{rossell2022concentration} to bound $L_{Q_SM}$, $L_{Q_SS}$, and hence also $L_{SM}$.

\begin{lem}\label{lemma:s7}
    Let $M,Q$ be any two nested models such that $M \subseteq Q$. Then
$$
L_{QM}\;=\;\|\bs X_{Q}\tilde{\bs\beta}^{(Q)}\|^2-\|\bs X_{M}\tilde{\bs\beta}^{(M)}\|^2\;\sim\; \chi_{|Q\setminus M|}^2\left(\mu_{Q M}\right), 
$$
where $\chi_{k}^2(\mu)$ denotes, when $\mu>0$, the noncentral chi-squared distribution with $k$ degrees of freedom and noncentrality parameter $\mu$ and, when $\mu=0$, the chi-squared distribution with $k$ degrees of freedom $\chi_{k}^2$. The parameter $\mu_{QM}$ is given by \begin{equation}\label{eq:muQM}
    \mu_{QM}\;\;:=\;\;\|\big(I_n-P_M\big) \bs X_{Q\setminus M} \bs \beta^*_{Q\setminus M}\|^2
\end{equation}
where $P_M=\bs X_{M}\left(\bs X_{M}^{\top} \bs X_{M}\right)^{-1} \bs X_{M}^{\top}$.
\end{lem}

If $M\supset S$ (over-fitted), then $Q_S=M$ and $\beta^*_{Q_S\setminus S}=\beta^*_{M\setminus S}=0$, because any parameter outside the true support $S$ is by definition $0$. Moreover, $-L_{SM}=L_{Q_SS}$ and by Lemma~\ref{lemma:s7}, $-L_{SM}\sim \chi^2_{|Q_S\setminus S|}$ since $\beta^*_{Q_S\setminus S}=0$. We have $\Delta_{MS}>0$ and if one sets large enough $\kappa_j$ (and thus $\Delta_{MS}$), then $C(S)-C(M)$ is also large and $NC(M)$ vanishes. If $M\subset S$ (under-fitted), then $Q_S=S$, $L_{SM}=L_{Q_SM}$, and by Lemma~\ref{lemma:s7}
$L_{SM}\sim \chi^2_{|Q_S\setminus M|}(\mu_{Q_SM})$, which has expectation $|Q_S\setminus M|+\mu_{Q_SM}$.
If the noncentrality parameter $\mu_{Q_SM}$ is large enough, then $L_{SM}$ and $C(S)-C(M)$ are also large, and $NC(M)$ vanishes. Lemma~\ref{lem:boundrho} below shows that large $\mu_{Q_SM}$ can be achieved by setting a betamin condition and an eigenvalue condition on $\bs X$ involving the following quantity  
 \begin{equation}\label{eq:rhoX}
        \rho(\bs X)\;\;=\;\;\min_{M \in \M:\, M\not \supseteq S} \lambda_{\min}\big(\tfrac{1}{n} \boldsymbol{X}_{S \setminus M}^{\top}\left(I_n- P_M\right) \boldsymbol{X}_{S \setminus M}\big).
    \end{equation}
The quantity $\rho(\bs X)$ is nonnegative and relates to how distinguishable the other models $M\in \M$ are from $S$ \citep{infotheowainwright}.  More specifically, $\tfrac{1}{n} \boldsymbol{X}_{S \setminus M}^{\top}\left(I_n- P_M\right) \boldsymbol{X}_{S \setminus M}$ is the sample covariance matrix of the residuals when regressing $\boldsymbol{X}_{S \setminus M}$ on $\boldsymbol{X}_M$.
%That is, $\rho(\bs X)$ measures how well one can approximate the $\boldsymbol{X}_{S \setminus M}$ with $\boldsymbol{X}_M$.
In an orthonormal case where $\bs X^{\top}\bs X=n\I_p$, then $\rho(\bs X)=1$.  

\begin{lem}\label{lem:boundrho}
 For any $M\in\M$, let $Q_S=S \cup M$. Then the non-centrality parameter satisfies  \eqref{eq:muQM}
    $\mu_{Q_SM}\; \geq \;n\,\rho(\bs X)\,\sum_{j=1}^b|S_j\setminus M_j|\,{\beta_{\min,\,j}^*}^2$.
\end{lem}
If $M\neq S$ is such that $M\not\supset S$ and $M\not\subset S$, simultaneously large enough $\kappa_j$ and $\mu_{Q_S M}$ guarantee that $C(S)-C(M)$ is also large, and that $NC(M)$ vanishes.


\subsection{Sufficient conditions for consistency with block $\ell_0$ penalties}
%\subsection{Selection based on block $\ell_0$ penalty}
\label{ssec:regression_mainresult}

We now state two conditions that are sufficient for asymptotically recovering $S$. 
Building on our previous discussion, we require the block penalties $\kappa_j$ to be large enough,  
and a betamin condition. \\[5pt] 
\hypertarget{lab:A6}{(A6)} \quad For each block $j$, there exists $f_j\to \infty$ (as $n\to \infty$) such that for every sufficiently large $n$,
\begin{equation*}
    \kappa_j \;=\;\ln(p_j-s_j) + f_j
\end{equation*} 
\hypertarget{lab:A7}{(A7)} \quad For each block $j$, there exists $g_j\to \infty$ such that for every sufficiently large $n$,
    \begin{equation*}
        \sqrt{\frac{(1-\gamma) n\rho(\bs X)}{6}}{\beta_{\min,j}^*} \,- \,\sqrt{\kappa_j} \;=\; \sqrt{\ln(s_j)} + g_j .
    \end{equation*}
where $\gamma:=\tfrac12(1+\max_j\ln(p_j-s_j)/\kappa_j)\in (\tfrac12,1)$.\\[5pt] 
These assumptions are similar to Assumptions~\hyperlink{lab:A4}{A4}--\hyperlink{lab:A5}{A5} formulated for the sequence model.
 By Proposition~\ref{prop:isthresh}, in the orthonormal 
 case where $\bs X^{\top}\bs X=n\I_p$, setting block penalties $\kappa_j$ 
 is equivalent to hard-thresholding with block thresholds $\tau_j=\sqrt{2 \kappa_j/n}$.
 Assumptions~\hyperlink{lab:A4}{A4}--\hyperlink{lab:A5}{A5} can then be translated into assumptions on $\kappa_j$, taking $\rho(\bs X)=1$. In that case, Assumption~\hyperlink{lab:A6}{A6} is of the same order as Assumption~\hyperlink{lab:A4}{A4}, up to a term $f_j$ that can grow at an arbitrarily slow rate. %An oracle choice for $f_j$ is presented in Theorem~\ref{theo:suffcondlinearmodel}. 
 Further, if one sets $f_j$ such that $\ln(p_j-s_j)= O(f_j)$ then $1-\gamma > 0$, and then Assumption~\hyperlink{lab:A7}{A7} essentially requires $n \rho(\bs X) \beta_{\min,j}^*$ to be larger than $\sqrt{\kappa_j} + \sqrt{\ln(s_j)}$ (up to constants and a term growing at an arbitrarily slow rate), analogously to Assumption~\hyperlink{lab:A5}{A5}.
% A natural choice stemming from a Bayesian perspective discussed in Section~\ref{sec:informed_l0}, is $f_j=0.5\ln(n)$. %for all $j$. 
%\color{blue} Then, if $p= O(n^k)$ for some $k>0$ we have that $1-\gamma>0$. \color{black}
% The betamin Assumption~\hyperlink{lab:A7}{A7} is the same as \hyperlink{lab:A5}{A5} up to, again, a term $g_j$ increasing at arbitrarily slow rate, and the factor $\sqrt{(1-\gamma)/3}$. \david{In \hyperlink{lab:A7}{A7} I see $(1-\gamma)/6$ not $(1-\gamma)/3$.} This factor is smaller than $1/2$ and depends on the choice of $f_j$'s in \hyperlink{lab:A6}{A6}. \david{"This factor..." is confusing, better state the factor.} It is bounded away from 0 as long as $f_j= \Omega(\ln(p_j-s_j))$ for all $j$. \david{Can we instead say $\ln(p_j-s_j)=O(f_j)$ in the previous sentence? And in the next sentence, say instead: For example, if $f_j=0.5 \ln(n)$ then $(1-\gamma)>0$ as long as $p=O(n^k)$.}
% This is the case for the earlier natural choice \david{"earlier natural choice is confusing", specify} whenever $n^k=\Omega(p)$ for some $k>0$. \paul{essentially NOT ultra high dimension} \david{Can we just say $p=O(n^k)$?} When $f_j= o(\ln(p_j-s_j))$ for some $j$, \hyperlink{lab:A7}{A7} requires $\sqrt{\kappa_j}+\sqrt{\ln(s_j)} = o(n\rho(\bs X)\beta_{min,j})$.
%\david{Earlier we said that $f_j$ can grow at an arbitrarily slow rate, but we now see that $f_j$ has to be at least of order $p_j-s_j$. How about saying: This factor depends on the choice of $f_j$ in \hyperlink{lab:A6}{A6}. For example, by default we take $f_j=...$, in which case $\gamma$ is bounded below by ... Perhaps also relate the choice of $(f_j,g_j)$ to the oracle $\kappa_j^*$'s discussed later in the theorem, like "an ideal choice of $(f_j,g_j)$ leads to oracle penalties $\kappa_j^*$ given in Theorem 4.2 below" (and, if it helps, briefly explain how the oracle $(f_j,g_j)$ depend on $s_j,p_j,\beta_{\min}^*$)} \paul{$f_j$ can grow at any rate.  For any growth rate you will discard asymptotically the overfitting models. When $f_j= o(\ln(p_j-s_j))$ for some $j$, the betamin requirement is stronger than $\Theta(\sqrt{\kappa_j} +\; \sqrt{\ln(s_j)} + g_j)$. To see it, the betamin requirement is essentially: $\sqrt{\frac{ n\rho(\bs X)}{2}}{\beta_{\min,j}^*}  \;=2(\max_j\frac{\ln(p_j-s_j)}{f_j}+1)(\sqrt{\kappa_j} +\; \sqrt{\ln(s_j)}) + g_j$. It is not very intuitive to have such change in betamin condition i know. It is linked to a technicality of the proof. I have tried hard to get rid of it but could not. I have explained more.}
%\david{I get confused by first seeing $f_j=\Omega(\ln(p_j-s_j))$ and later $f_j=o(\ln(p_j-s_j))$.}
Finally, we remark that in our proof one could take a betamin condition that is slightly less strict than Assumption~\hyperlink{lab:A7}{A7}, but we present Assumption~\hyperlink{lab:A7}{A7} here to facilitate comparison to Assumption~\hyperlink{lab:A5}{A5} in the sequence model.  


We can now state our main theorem on the strong variable selection consistency of $\hat{S}^{b}$. The result holds for either fixed or diverging $p_j-s_j$ and $s_j$. %\color{green}holding fixed or diverging numbers of signals:\color{black}
\begin{thm}\label{theo:suffcondlinearmodel}

Under Assumptions~\hyperlink{lab:A1}{A1}, \hyperlink{lab:A6}{A6} and, \hyperlink{lab:A7}{A7}, we have
$$\lim_{n \to \infty} \sum_{M \in \M\setminus \{S\}}\E\left(NC(M)\right)= 0 \qquad \mbox{and} \qquad \lim_{n \to \infty} P(\hat{S}^b = S)= 1.$$
\end{thm}


Theorem~\ref{theo:suffcondlinearmodel} considers a linear penalty across blocks
$\sum_{j=1}^b \kappa_j|M_j|$. This linearity assumption can be relaxed, and in section S5 of the Supplement we give consistency results for nonlinear penalties. 
%Those apply for example to the EBIC penalty \citep{EBIC} and improve on its sufficient betamin conditions for selection consistency given in \cite{ebicdiverging}. 

Prior results support the tightness of the conditions of Theorem~\ref{theo:suffcondlinearmodel}. Firstly, when translated to the orthonormal setting Assumptions~\hyperlink{lab:A6}{A6} and \hyperlink{lab:A7}{A7} are similar to Assumptions~\hyperlink{lab:A4}{A4}--\hyperlink{lab:A5}{A5} where shown to be necessary or near-necessary in Section~\ref{sec:properties_orth}. Secondly, when applied to $\hat{S}$ ($\hat{S}^b$ with $b=1$), Theorem~\ref{theo:suffcondlinearmodel} shows that, for a wide range of asymptotic regimes, standard $\ell_0$ penalties achieve support recovery under the same conditions as an optimal selector that already knows $s$ analyzed in \cite{infotheowainwright}. In some regimes, our sufficient assumptions are even weaker than those for the optimal selector in \cite{infotheowainwright}. In Supplement S7, we precisely discuss the tightness of Assumptions~\hyperlink{lab:A6}{A6} and \hyperlink{lab:A7}{A7}, also showing they match our necessary conditions of Section~\ref{ssec:neccond_linreg} in a wide range of regimes. Moreover, unlike the results in \cite{infotheowainwright}, Theorem~\ref{theo:suffcondlinearmodel} does not make distributional assumptions on $\bs X$ and shows consistency in a $L_1$ sense.

From the proof of Theorem~\ref{theo:suffcondlinearmodel}, we can bound the convergence rate of $P(\hat S^{b}\neq S)$, given in Theorem~\ref{prop:fullrecovreg}. We also give oracle block penalties that approximately optimize the bound, and provide the resulting oracle rate of convergence. 
In the statement, $\delta<1$ and $r>1$ should be understood as being arbitrarily close to 1. 
%\david{The theorem below is really a corollary no?}\paul{It is not really a corollary. It is actually the other way around: in the proof of Theorem~\ref{theo:suffcondlinearmodel} I prove the first bound below and the consistency follows from that bound}
\begin{thm}\label{prop:fullrecovreg}
Assume \hyperlink{lab:A1}{A1}, \hyperlink{lab:A6}{A6}, \hyperlink{lab:A7}{A7}. %\david{Don't we assume \hyperlink{lab:A2}{A2} as in Theorem 4.5?}\paul{No, you don't need \hyperlink{lab:A2}{A2} to state the bound, it is valid for all regimes of $p$} 
Then, for all sufficiently large $n$ and any $\delta \in (0,1)$ and $r>1$,
\begin{equation}\label{eq:rateconvreg}
 P(\hat{S}^{b} \neq S)\;\leq\; 6(2^{2b}-2b)r\,\sum_{j=1}^b\,\, e^{-\tfrac{\delta}{2}\big[\kappa_j-\ln(p_j-s_j)\big]} + 
    e^{-\tfrac{\delta}{2}\big[(\sqrt{\frac{(1-\gamma) n\rho(\bs X)}{6}}{\beta_{\min,j}^*} - \sqrt{\kappa_j})^2 -\ln(s_j) \big]}.
\end{equation}
%where $r_b\to 1$ as $n\to \infty$.
Moreover, suppose that for all $j=1,\ldots,b$ $\lim_{n\to \infty}\big(\sqrt{(1-\gamma) n\rho(\bs X)/3}\;\beta_{\min,j}^*\big)/\big(\sqrt{2\ln(p_j-s_j)} + \sqrt{2\ln(s_j)}\big) >1$ and the $\kappa_j$ are set at the oracle values
	\begin{equation} \label{eq:oraclepenreg}
\sqrt{\kappa^*_j\,} \;=\; \frac12\sqrt{\frac{(1-\gamma) n\rho(\bs X)}{6}}{\beta_{\min,j}^*}+\frac12\sqrt{\frac{6}{(1-\gamma) n\rho(\bs X)}}\frac{1}{\beta_{\min,j}^*}\big(\ln({p_j-s_j})-\ln({s_j})\big).
\end{equation}
Then Assumptions~\hyperlink{lab:A6}{A6} and \hyperlink{lab:A7}{A7} hold. Moreover, if Assumption~\hyperlink{lab:A1}{A1} holds too, then
% Then for all sufficiently large $n$, Assumptions~\hyperlink{lab:A6}{A6} and \hyperlink{lab:A7}{A7} hold and, if Assumption~\hyperlink{lab:A1}{A1} holds too, 
%for any $\delta \in (0,1)$ and $r>1$,
\begin{equation}\label{eq:oraclerateconvreg}
    P(\hat{S}^{b} \neq S) \; \leq \; 12(2^{2b}-2b)r \sum_{j=1}^b e^{-\tfrac{\delta}{2}\big[\frac{(1-\gamma) n\rho(\bs X)}{24}{\beta_{\min,j}^*}^2-\ln \max\{p_j-s_j,s_j\}\big]}.
\end{equation}
%where $r_b\to 1$ as $n\to\infty$.
\end{thm}
In the orthonormal setting, using the equivalence $\tau_j=\sqrt{2 \kappa_j/n}$, the bounds in \eqref{eq:rateconvreg} and \eqref{eq:oraclepenreg} nearly recover the tight bounds in \eqref{eq:rateconvgsm} and \eqref{eq:optimtauj} for the sequence model. 

\subsection{Necessary assumptions for consistency with block $\ell_0$ penalties}\label{ssec:neccond_linreg}

We derive assumptions on $\kappa_j$ and $\beta_{min,j}$ that are necessary for consistent variable selection with $\hat{S}^b$. %Recall that the strength of the block penalties $\kappa_j$ drives the vanishing of the pseudo-posterior probabilities of over-fitted models.
%The block penalties need then to be large enough for the pseudo-posterior probabilities of even the least penalized over-fitted models to vanish. 
For every $j=1,\ldots,b$, denote by $O_j$ the set of models that over-fit by only one variable from block $B_j$:
$$O_j \:=\,\{M\in \M\;|\; M=S\cup\{i\} \text{ where } i\in B_j\setminus S_j\}.$$
The asymptotic recovery of $S$ implies that, $\max_{M\in O_j}NC(M)/NC(S)<1$ with probability going to 1 as $n$ grows, for every $j=1,\ldots,b$. For every $M \in O_j$ the ratio $NC(M)/NC(S)$ grows with $L_{MS}$
%$M_k\in O_j$, the ratio $NC(M_k)/NC(S)$ grows with $L_{M_kS}$ 
which, by Lemma~\ref{lemma:s7}, is $\chi^2_1$ distributed (note that $\bs \beta_{Q_S\setminus S}^*=\bs \beta_{M\setminus S}^*=0$). Then, for every $M \in O_j$, there exists $Z_M \sim N(0,1)$ such that $L_{MS}=Z_{M}^2$.
%$M_k\in O_j$, $Z_{M_k} \sim N(0,1)$ such that $L_{M_kS}=Z_{M_k}^2$. 
%A necessary assumption on $\kappa_j$ follows from properties of the maximum of correlated Gaussian variables.
Let
%Denote the smallest eigenvalue of their correlation matrix:
\begin{align}
    &\underline{\lambda}_j\;:=\;\lambda_{min}(C^j)\;\;\text{where} \;\;C^j_{k,l}=\operatorname{corr}\big(Z_{M_k},Z_{M_l}\big),\;\;\forall M_k,M_l \in O_j
    \nonumber \\
    &\underline{\lambda}\;:=\;\lambda_{min}(C)\;\;\text{where} \;\;C_{k,l}=\operatorname{corr}\big(Z_{M_k},Z_{M_l}\big),\;\;\forall M_k,M_l \in \cup_{j=1}^b O_j.
    \nonumber
\end{align}
%\paul{note it is not the correlation matrix of the $\sqrt{L_{MS}}=|Z_M|$. This result is based on \cite{hartigan} that gives a lower bound for maximum of correlated Gaussians and not maximum of absolute value of correlated Gaussians. The lower bound is probably the same but proving it would require additional lemmas.}
%We denote $\underline{\lambda}$, the smallest eigenvalue of the correlation matrix of the $Z_{M_k}$, $M_k\in\cup_{j=1}^b O_j$.
That is, small $\underline{\lambda}_j$ indicates that truly inactive variables in block $j$ are highly correlated with each other, whereas $\underline{\lambda}_j=1$ that they are uncorrelated. Proposition~\ref{prop:nectypeIreg} below describes how $\underline{\lambda}_j$ relates to a necessary condition for consistency.
In Section \ref{ssec:regression_gains} we discuss how $\underline{\lambda}$ describes settings where $\hat S$ is not consistent but $\hat S^b$ is. 
%A lower bound on the maximum of correlated Gaussians in \cite{hartigan} yields the following result.
\begin{prop}\label{prop:nectypeIreg}
Assume \hyperlink{lab:A1}{A1}-\hyperlink{lab:A2}{A2}.
     If for some $j=1,\ldots,b$,  $\lim_{n\to\infty}\kappa_j/(\underline{\lambda}_j^2\ln(p_j-s_j)) <1$, then  %under Assumptions~\hyperlink{lab:A1}{A1} and \hyperlink{lab:A2}{A2}, 
     $$\lim_{n \to \infty} P\Big(\max_{M\in O_j}\frac{NC(M)}{NC(S)}<1\Big)= 0 \quad\text{and}\quad \lim_{n \to \infty} P(\hat{S}^b=S)= 0.$$
\end{prop}



Proposition~\ref{prop:nectypeIreg} shows that 
a necessary assumption to recover $S$ asymptotically
is 
\begin{equation}\label{cond:necIallblockpen}
\lim_{n\to\infty}\frac{\kappa_j}{\underline{\lambda}_j^2\ln(p_j-s_j)}\;\geq\;1\qquad\mbox{for all }j=1,\ldots,b.
\end{equation}
%When the inactive variables are orthogonal, for every $j=1,\ldots,b$ and $M\in O_j$, the $Z_M$ are independent and $\underline{\lambda}_j=1$. Then
When $\underline{\lambda}_j=1$ as in the orthonormal case, \eqref{cond:necIallblockpen} recovers %Under orthogonal design, \eqref{cond:necIallblockpen} then recovers 
the necessary condition shown in Proposition~\ref{prop:suffnecblockthresh} (ii) for the sequence model. More generally we have $\underline{\lambda}_j \in [0,1]$, and then \eqref{cond:necIallblockpen} is actually a milder condition than that shown in Proposition~\ref{prop:suffnecblockthresh} (ii). That is, $\underline{\lambda}_j=1$ corresponds to the worst case, in terms of controlling false positives.

Consider now an under-fitted model $M\subset S$. %The recovery of $S$ implies that with probability going to 1 as $n$ grows, $NC(M)/NC(S)<1$. 
By definition the ratio $NC(M)/NC(S)$ %By definition of $NC$ the ratio 
grows with block penalties $\kappa_j$ and shrinks with $L_{SM}$, which is distributed $\chi^2_{|S\setminus M|}(\mu_{SM})$ by Lemma~\ref{lemma:s7}.
For the ratio to be small, $\mu_{SM}$ must grow
%Small values of the ratio are attained only if $\mu_{SM}$ grows 
fast enough compared to $\kappa_j$. Lemma~\ref{lem:upperboundnoncentral} shows that $\mu_{SM}$ is bounded by the largest active signals in $S$ that is not in $M$.
\begin{lem}\label{lem:upperboundnoncentral}
    For any $T\subseteq S$, $\mu_{ST} \leq n\,\bar{\lambda}\,\sum_{j=1}^b|S_j\setminus T_j|\,\max_{i\in S_j\setminus T_j}{\beta_{i}^*}^2$  where $\bar{\lambda}\,:=\,\lambda_{\max}\Big(n^{-1} \bs X_{S}^\top\bs X_{S}\Big)$
\end{lem}

It follows that a necessary condition is that the active $\beta_i^*$ that are missing in any underfitted $M$ are not too small. We next formally define what we mean by small and large signals, and also define a subset of intermediate signals 
%introduce a formal distinction between small and large signals in $S$. We also introduce a subset of signals that are not small or large 
that will be used in Section~\ref{ssec:nobetaminhyp}.  For fixed penalties $\kappa_j$'s, and for %and every $j=1,\ldots,b$, 
$\gamma$ and $g_j$ as defined in Assumption~\hyperlink{lab:A7}{A7}, let
\begin{eqnarray}\label{eq:defSes}
    &S_j^{S}(\kappa)\;&:=\;\Bigg\{\beta_i^* \in S_j \,\Big|\,\sqrt{n\bar{\lambda}}|\beta_{i}^*|=o\big(\sqrt{\kappa_j}\big) \Bigg\}\nonumber\\
    &S_j^{L}(\kappa)\;&:=\;\Bigg\{\beta_i^* \in S_j \,\Big|\,\sqrt{\frac{(1-\gamma) n\rho(\bs X)}{6}}|\beta_{i}^*| \,- \,\sqrt{\kappa_j} \;=\; \sqrt{\ln(s_j)} + g_j\Bigg\}\label{eq:defslsi}\\
    &S_j^{I}(\kappa)\;&:=\;S_j \setminus \big(S_j^{L}(\kappa) \cup S_j^{S}(\kappa)\big).
    \nonumber 
\end{eqnarray} \\
The subset $S_j^{S}(\kappa)$ gathers signals in $S_j$ that are small with respect to the penalty $\kappa_j$, $S_j^{L}(\kappa)$ those that are large in that they satisfy Assumption~\hyperlink{lab:A7}{A7}, and $S_j^{I}(\kappa)$ those that are neither large nor small. 
Proposition~\ref{prop:nectypeIIreg} below says that if the set of small signals $S^{S}(\kappa)=\cup_{j=1}^b S_j^{S}(\kappa)$ is not empty, then consistent model recovery is not possible.
%Take $M=S\setminus S^{S}(\kappa)$ where $S^{S}(\kappa)=\cup_{j=1}^b S_j^{S}(\kappa)$. If $S^{S}(\kappa)\neq \emptyset$, the missing active signals in $M$ and thus $\mu_{SM}$ are all small compared to the $\kappa_j$'s. \david{Should it be that "some active signals are small"? (rather than all)} 
%We then establish the following result.
%Using the properties of the noncental $\chi^2$ distribution, we then establish the following convergence of the probabilities of $(NC(M)/NC(S)<1)$ and thus of $(\hat{S}^{b}=S)$.

\begin{prop}\label{prop:nectypeIIreg}
    If $S^{S}(\kappa)\neq\emptyset$ and for all $j=1,\ldots,b$, $\kappa_j\to \infty$, then, under Assumption~\hyperlink{lab:A1}{A1},
    $$\lim_{n \to \infty} P\Big(\frac{NC(S\setminus S^{S}(\kappa))}{NC(S)}<1\Big)= 0 \quad\text{and}\quad \lim_{n \to \infty} P(\hat{S}^b=S)= 0.$$
\end{prop}
%Property~\ref{prop:nectypeIIreg} shows that for any increasing penalty, when $S$ contains signals that are too small with respect to that penalty, the smaller model $S\setminus S^{S}(\kappa)$ is favored at the expense of $S$ and recovery fails. 
It follows that a necessary assumption for variable selection consistency is
\begin{equation}\label{cond:necIIallblockpen}
\lim_{n\to\infty}\frac{\sqrt{n\bar{\lambda}}\beta_{\min,j}^*}{\sqrt{\kappa_j}}\;>\;0\qquad\mbox{for all }j=1,\ldots,b.
\end{equation}
When active variables are orthonormal, %orthogonal and normalized, 
that is $\bs X_S^{\top}\bs X_S = nI$, then $\bar{\lambda}=1$. 
Assumption~\eqref{cond:necIIallblockpen} is then milder and less tight than the corresponding necessary assumption on signal strength shown in Proposition~\ref{prop:suffnecblockthresh} (iv)  in the sequence model. 

An immediate consequence of Proposition~\ref{prop:nectypeIreg} and Proposition~\ref{prop:nectypeIIreg} is the next result.
\begin{cor}\label{cor:necrecovreg}
    If for some $j\in\{1,\ldots,b\}$   $\lim_{n\to\infty}\sqrt{n\bar{\lambda}}\beta_{\min,j}^*/\big(\underline{\lambda}_j\sqrt{\ln(p_j-s_j)}\big) = 0$, then, under Assumption~\hyperlink{lab:A1}{A1} and \hyperlink{lab:A2}{A2},  $\lim_{n \to \infty} P(\hat{S}^b=S) < 1$.
\end{cor}
Corollary~\ref{cor:necrecovreg} yields a necessary assumption for consistency
\begin{equation}\label{eq:betaminnecreg}
\lim_{n\to\infty}\frac{\sqrt{n\bar{\lambda}}\beta_{\min,j}^*}{\underline{\lambda}_j\sqrt{\ln(p_j-s_j)}} > 0\qquad\mbox{for all }j=1,\ldots,b.
\end{equation}
%In an orthonormal design where $\bs X^{\top} \bs X=nI$, %for all $j=1,\ldots,b$, 
%then $\underline{\lambda}_j=1$ and $\bar{\lambda}=1$. 
In the particular case of orthonormality and diverging $s$, the necessary assumption \eqref{eq:betamin} given for the sequence model is then stricter and tighter than \eqref{eq:betaminnecreg}. Assumption \eqref{eq:betaminnecreg} is however a more general necessary betamin condition that applies in all correlated, orthonormal, low and high dimensional settings, for $s$ fixed or diverging. In Supplement S7, we contrast necessary condition~\eqref{eq:betaminnecreg} with the sufficient conditions, assumptions~\hyperlink{lab:A6}{A6} and \hyperlink{lab:A7}{A7}, for standard selector $\hat{S}$. Our analysis shows that they imply the same scaling of $(n,p)$ in a wide range of regimes, confirming their tightness. %\paul{I am happy about this result because I believe there was no such tight necessary assumption for variable selection consistency in regression in the literature. The only ones I found were for LASSO (far less efficient/tight for selection in correlated settings) or compressed sensing}.

\subsection{Benefits of block penalties}
\label{ssec:regression_gains}

We discuss separately the benefits in terms of sufficient conditions for variable selection consistency and those in terms of convergence rate. The results are analogous to those in Section~\ref{sec:properties_orth} for the sequence model. %main message is that benefits discussed in the sequence model in Section~\ref{sec:properties_orth} extend to linear regression. 

\smallskip

\noindent Assumptions~\hyperlink{lab:A6}{A6}--\hyperlink{lab:A7}{A7} give ranges of penalties that are sufficient (but not necessary) for asymptotic support recovery. For simplicity, we restrict our discussion to choices of penalty $\kappa_j$ such that, in Assumption~\hyperlink{lab:A6}{A6}, $\ln(p_j-s_j)=O(f_j)$ for all $j$, so that $\gamma$ is bounded in $(0,1/2)$. For the standard selector $\hat{S}$, the single penalty $\kappa$ is essentially required to satisfy,
for some sequences $f,g\to\infty$ and up to constants
\begin{equation}\label{eq:rangepenex}
  \sqrt{\ln(p-s)}+ f\;\;\leq\;\;\sqrt{\kappa}\;\;\leq\;\;\sqrt{n\rho(\bs X)}\beta^*_{\min}+\sqrt{\ln(s)} + g.	
\end{equation}
For a block selector $\hat{S}^{b}$, the ranges for the $\kappa_j$'s essentially are, for some sequences $f_j,g_j\to\infty$ and up to constants
    \begin{eqnarray}\label{eq:rangeblockpenex}
      \sqrt{\ln(p_j-s_j)}+f_j\;\;&\leq\;\;\sqrt{\kappa_j}\;\;\leq&\;\;\sqrt{n\rho(\bs X)}\beta^*_{\min,j}-\sqrt{\ln(s_j)} + g_j.
      %
%  \sqrt{\ln(p_1-s_1)}\;\;&<\;\;\sqrt{\kappa_1}\;\;<&\;\;\sqrt{n\rho(\bs X)}\beta^*_{\min,1}-\sqrt{\ln(s_1)}\nonumber\\
%  &\;\;\vdots \\
%  \sqrt{\ln(p_b-s_b)}\;\;&<\;\;\sqrt{\kappa_b}\;\;<&\;\;\sqrt{n\rho(\bs X)}\beta^*_{\min}-\sqrt{\ln(s_b)} \nonumber.       
    \end{eqnarray}
Akin to the sequence model, if there exist two blocks such that the ranges in \eqref{eq:rangeblockpenex} do not overlap, then a constant $\kappa$ cannot satisfy \eqref{eq:rangeblockpenex} for both blocks and consistent selection may not be possible.
%by definition, the ranges for $\hat S^b$ are wider than that for $\hat S$. 
That is, sufficient conditions for variable selection consistency are milder with block penalties. Although not discussed here for brevity, block penalties also lead to similar improvements on the necessary conditions discussed in Section \ref{ssec:neccond_linreg}, relative to those standard $\ell_0$ penalties.
Corollary~\ref{cor:selconstimpossiblereg} gives conditions under which consistent selection is possible with $\hat S^b$ but not with $\hat S$.
\begin{cor}\label{cor:selconstimpossiblereg}
Assume \hyperlink{lab:A1}{A1}, \hyperlink{lab:A2}{A2}, \hyperlink{lab:A6}{A6}, \hyperlink{lab:A7}{A7}. If
$\lim_{n\to\infty}\big(\sqrt{n\bar{\lambda}}\beta_{\min}^*\big)/\big(\underline{\lambda}\sqrt{\ln(p-s)}\big) = 0$
then 
$P(\hat{S} = S)\not \to 1$ and $P(\hat{S}^{b} = S)\to 1$.
\end{cor}

Observe that \eqref{eq:rangepenex} and \eqref{eq:rangeblockpenex} are analogous to \eqref{eq:rangethreshex} and \eqref{eq:rangeblockthreshex} for the sequence model, up to rescaling by $\sqrt{2/n}$, the factor $\rho(\bs X)$ and sequences $(f_j,g_j)$ growing arbitrarily slowly. The gains in terms of valid thresholds for consistency discussed in the examples of Section \ref{sec:examples} remain applicable %are essentially the same for 
to linear regression. %The main difference lies in that $\rho(\bs X)$ may be smaller than 1 making the range for $\hat S$ narrower and $\hat S^b$ even more beneficial. 

Let $\underline{\beta}^{*,b}_{\min,reg}$ and $\underline{\beta}^{*}_{\min,reg}$ be the smallest signal recoverable by $\hat{S}^b$ and $\hat{S}$ respectively.
%Consider the smallest signal recoverable by $\hat{S}^b$ and $\hat{S}$ while still being variable selection consistent that we denote $\underline{\beta}^{*,b}_{\min,reg}$ and $\underline{\beta}^{*}_{\min,reg}$. 
Assuming $\beta_{min}^*$ is in block $b$, Assumptions~\hyperlink{lab:A6}{A6}-\hyperlink{lab:A7}{A7} essentially require that $\underline{\beta}^{*,b}_{\min,reg}$ and $\underline{\beta}^{*}_{\min,reg}$ satisfy for some sequences $g,h\to\infty$ and up to constants
%they are essentially required by \hyperlink{lab:A6}{A6} and \hyperlink{lab:A7}{A7} to satisfy, up to bounded factors,
\begin{eqnarray*}
    \underline{\beta}^{*,b}_{\min,reg}&\geq& \sqrt{\frac{\ln(p_b-s_b)}{ n\rho(\bs X)}} + \sqrt{\frac{\ln(s_b)}{ n\rho(\bs X)}} + g,\quad\text{and}  \\
    \underline{\beta}^{*}_{\min,reg}&\geq& \sqrt{\frac{\ln(p-s)}{ n\rho(\bs X)}} + \sqrt{\frac{\ln(s)}{ n\rho(\bs X)}}  + h.
\end{eqnarray*}
These lower bounds are the same as \eqref{eq:betaminbardef} for the sequence model, up to rescaling by $\sqrt{2}$, a factor $\rho(\bs X)$ and $g$ and $h$ which can grow arbitrarily slowly with $n$. Hence, the discussion and examples of the benefits in the smallest recoverable signals in Section~\ref{sec:examples} extend to linear regression.

Let ${OR}^b_{reg}$ be the oracle convergence rate for $\hat{S}^{b}$ in Theorem~\ref{prop:fullrecovreg}, and ${OR}_{reg}$ that for $\hat{S}$. Then we have
\begin{equation*}
    \frac{{OR}^b_{reg}}{{OR}_{reg}}=(2^{2b-1}-b)\sum_{j=1}^be^{-\frac{\delta}{2}\left[\frac{ n\rho(\bs X)}{24}\left((1-\gamma){\beta^*_{\min,j}}^2-(1-\gamma'){\beta_{\min}^*}^2\right)+\ln \max\{p-s,s\}-\ln \max\{p_j-s_j,s_j\}\right]},
\end{equation*}
where $\gamma=\tfrac12(1+\max_j\ln(p_j-s_j)/\kappa_j^*)$ and $\gamma'=\tfrac12(1+\ln(p-s)/\kappa^*)$. The ratio above is essentially the same as \eqref{eq:ratioCR} for the sequence model, up to certain factors.  Specifically, $\delta$, $(1-\gamma)$ and $(1-\gamma')$ are close to 1, whereas $ \rho(\bs X)/24$ slows the convergence rate gains but does not alter the essence of the ratio. % as long as $\rho(\bs X)$ remains bounded above 0. 
The factor $2^{2b-1}$ highlights however a potential limitation: guarantees of gains with block $\ell_0$ penalties deteriorate when one considers a large number of blocks $b$.
We remark that such deterioration may be a consequence of our proof strategy, rather than an inherent limitation of block penalties. Studying cases with %large $b$, including 
$b \to \infty$ is left as future research.

\subsection{Convergence with no betamin condition}\label{ssec:nobetaminhyp}

We now derive a convergence result for pseudo-posterior probabilities $NC(M)$, $M\in\mathcal{M}$ under no assumption on the minimal signal strength. The result generalizes Theorem~\ref{theo:suffcondlinearmodel} and is key to the proofs of Section~\ref{sec:informed_l0}.

%In \eqref{eq:defslsi}, we defined, for a choice of penalties $\kappa_j$, the set of large and intermediate signals. 
Let $\T(\kappa)$ be the set of models that contain all large signals signals in $S^{L}(\kappa)=\cup_{j=1}^b S_j^{L}(\kappa)$, and neither truly inactive parameters in $V\setminus S$ nor small signals in $S^S(\kappa)=\cup_{j=1}^b S^{S}_j(\kappa)$. That is,
\begin{equation} \label{eq:Tkappa}
   \T(\kappa) \;:=\;\Big\{M \in \M \,\big|\, M\,=\, S^L(\kappa)\,\cup\,R,\; R\in \mathcal{P}\big(S^{I}(\kappa)\big)  \Big\},
\end{equation}
where for a set $A$, $ \mathcal{P}(A)$ denotes the power set of $A$.
Theorem~\ref{thm:convtoT} below shows,
%Following a strategy similar to that of the proof of Theorem~\ref{theo:suffcondlinearmodel}, we demonstrate, 
assuming only sufficiently large penalties as given in Assumption~\hyperlink{lab:A6}{A6},  that posterior model pseudo-probabilities concentrate on $\T(\kappa)$.
%the vanishing of the sum of pseudo-probabilities of all models but those in $\T(\kappa)$. 
%The main difference with respect to the proof of Theorem~\ref{theo:suffcondlinearmodel} is that, for any $M\not\in \T(\kappa)$, we bound $\E(NC(M))$ by a function of $C(M)-C(T)$ for a well chosen model $T$ in $\T(\kappa)$. 
\begin{thm}\label{thm:convtoT}
    Assume \hyperlink{lab:A1}{A1}, \hyperlink{lab:A6}{A6}, $|S^{I}(\kappa)|=O(1)$ and $|S^{S}_j(\kappa)|=O(p_j-s_j)$ for every $j=1,\ldots,b$. Then
    $$\lim_{n \to \infty} \sum_{M \in \M\setminus \T(\kappa)} \E\left(NC(M)\right)= 0 \qquad \mbox{and} \qquad
    \lim_{n \to \infty} P(\hat{S}^b \in \T(\kappa))= 1.$$
\end{thm}
%Theorem~\ref{thm:convtoT} shows that, asymptotically, with only assuming
That is, if one sets
sufficiently large penalties, all inactive and all small active signals are discarded. On the other hand, all large signals (relative to the specified penalties) are retained. Intermediate signals in $S^{I}(\kappa)$ may or may not be retained. Assumption $|S^{I}(\kappa)|=O(1)$ was made for simplicity, in fact $|S^{I}(\kappa)|$ can be allowed to increase moderately. 

\section{Data analysis with block $\ell_0$ penalization}\label{sec:informed_l0}

In Sections~\ref{sec:properties_orth} and~\ref{sec:properties_legression}, we derived properties for block penalties where one sets the penalties to oracle values.
We now
propose two data analysis methods anchored in an empirical Bayes perspective
that do not require oracle values, and adapt to the unknown sparsity in the data-generating truth.
The main idea is that, by using a BIC approximation to the marginal likelihood, one may estimate the proportion of truly active variables in each block by the average posterior inclusion probabilities in that block. This provides
%A Bayesian approach leads to a natural estimator of the number of truly active variables in each block and thus, by using empirical Bayes, 
a straightforward way to adapt penalties to sparsity in each block. Relying on the BIC approximation also allows the use of fast Bayesian computational methods that overcome the intractability of $\ell_0$ penalties. 
We remark that, despite their Bayesian motivation, the methods are fully data-dependent and do not require any prior distribution.

\subsection{An estimator of sparsity}\label{ssec:est.sparsity}

In a Bayesian framework for \eqref{eq:linearmodel} the model $M=(m_1,\ldots,m_p)$ is a vector of variable inclusion indicators $m_i=I(\beta_i \neq 0)$. Let $|M|=\sum_{j=1}^p m_j$. Consider a joint prior on parameters and models 
\begin{equation}\label{eq:bayesianframework}
    p(\boldsymbol{\beta},M \mid \boldsymbol{\theta})\;=\;p( \boldsymbol{\beta} \mid  M) p(M \mid \boldsymbol{\theta})
\end{equation}
where $p( \boldsymbol{\beta} \mid  M)$ is a prior on regression coefficients given the model, and 
$p(M \mid \boldsymbol{\theta})$ the prior probability of model $M$. The latter depends on hyperparameters $\bs \theta$ giving the prior inclusion probabilities in each block.
Specifically, assume that variable inclusions are independent a priori,
with constant inclusion probability $\theta^{(j)}$ within each block $j$. Then
\begin{eqnarray} \label{eq:block_modelprior}
p( M \mid \boldsymbol{\theta}) \propto
        \prod_{i=1}^p \operatorname{Bern}\left(m_i ; \theta_i\right) I(M \in \M),
\end{eqnarray}
where $\forall \, i\in B_j, \: \theta_i=\theta^{(j)}$ and $\bs \theta=(\theta^{(1)},\ldots,\theta^{(b)})$.

Posterior model probabilities are $p(M \mid \bs y, \bs \theta) \propto p(\bs y \mid M) p(M \mid \bs \theta)$, where $p(\bs y \mid M)$ is the so-called marginal likelihood of model $M$. The BIC approximation \citep{BIC} to $p( \bs y  \mid M)$ gives
%\begin{eqnarray*}
% p( \bs y  \mid M)\;\approx\; p(\bs y \mid \tilde{\bs \beta}^{(M)})n^{-|M|/2}
%\end{eqnarray*}
%Hence
%It follows that the posterior model probabilities satisfy
\begin{eqnarray*}
\ln p(M \mid\bs y, \bs \theta) \;\approx\; \ln p(\bs y \mid \tilde{\bs \beta}^{(M)}) - \frac{|M|}{2}\ln(n) + \ln p(M\mid\bs \theta) + c_M,
\end{eqnarray*}
for a wide family of priors $p(\bs \beta \mid M)$, where $\tilde{\bs \beta}^{(M)}$ is the MLE under model $M$, and $c_M$ a constant that may depend on $M$.
In particular, when $p(\bs \beta \mid M)$ is Zellner's unit information prior \citep{zellner1986assessing}, then $c_M$ does not depend on $M$ and the approximation is exact by replacing $\tilde{\bs \beta}^{(M)}= (\bs X_M^\top \bs X_M + n^{-1} I)^{-1} \bs X_M^\top \bs y$.
Neglecting the constant $c_M$, simple algebra shows that 
the block $\ell_0$ penalty selector discussed in Section \ref{sec:properties_legression} approximately maximizes $p(M \mid \bs y, \bs \theta)$. %for some $p(M \mid \bs \theta)$. 
Specifically, take the normalized criterion $NC(M)$ in \eqref{eq:CM} to be proportional to
\begin{align}
p(\bs y \mid \tilde{\bs \beta}^{(M)})  
\prod_{j=1}^b \left( n^{\frac{1}{2}}(1/\theta^{(j)} - 1) \right)^{-|M_j|},
    \nonumber
\end{align}
which corresponds to taking the block penalties  
%Then $\ln p(M \mid \bs \theta)$ defines the block penalties studied in Section \ref{sec:properties_legression}
\begin{equation}\label{eq:blockpenpriorinclu}
    \kappa_j=\tfrac{1}{2}\ln(n)+\ln\big(1/{\theta^{(j)}} - 1\big).
\end{equation}

In summary, one may think of $NC(M) \approx P(M \mid \bs y,\bs \theta)$, where the $\kappa_j$'s are a suitable function of $\theta^{(j)}$. %, the normalized score in Section \ref{sec:properties_legression}.
This connection motivates estimating the number of truly active variables in block $j$ by %for any block $\ell_0$ procedure:
\begin{equation}\label{eq:shat}
    \hat{s}_j \,:=\, \sum_{i\in B_j}\,\,\sum_{M\in \M | i \in M}NC(M)
%\end{equation}
%The normalized score $NC(M)$ can be seen as the posterior probability for model $M$ in Bayesian framework \eqref{eq:bayesianframework}. 
%To be clear, replacing $NC(M)$ by $P(M \mid \bs y,\bs \theta)$ in \eqref{eq:shat} gives
%\begin{equation*}
\approx    \sum_{i\in B_j}\,\,\sum_{M\in \M | i \in M}P(M\mid \bs y,\bs \theta) = \sum_{i\in B_j}P(\beta_i \neq 0 \mid \bs y,\bs \theta),
\end{equation}
where the right-hand side is the posterior mean $E(s_j \mid \bs y,\bs \theta)$.
%Hence, $\hat{s}_j$ is the posterior mean of $s_j$ in a setting where one uses the BIC approximation to the marginal likelihood.
One may then set prior inclusion probabilities $\hat{\theta}^{(j)}= \hat{s}_j/p_j$, the estimated proportion of truly active variables in block $j$.
%From a Bayesian perspective, $\hat{s}_j$ is then %the sum of posterior inclusion probabilities and 
%the posterior mean of $s_j$. 
Section S0 of the Supplement discusses how $\hat{\theta}^{(j)}=\hat{s}_j/p_j$ can also be motivated as an approximation to an empirical Bayes estimator maximizing the marginal likelihood of $\bs y$ given $\bs \theta$.

We now show that, besides being well-founded from a Bayesian perspective, $\hat{\theta}^{(j)}=\hat{s}_j/p_j$ %and, as shown below, also 
has attractive frequentist properties under mild assumptions. In \eqref{eq:defslsi} we defined $S^{S}_j(\kappa)$ to be the set of small signals in block $j$ and  $S^{L}_j(\kappa)$ the set of larger signals. In Theorem~\ref{thm:convtoT} we also derived a convergence result on pseudo posterior probabilities that yields the following asymptotic bounds on the frequentist expectation of $\hat{s}_j/p_j$.
\begin{prop}
    \label{prop:shatconsistence}
Assume \hyperlink{lab:A1}{A1}, \hyperlink{lab:A6}{A6}, $|S^{I}(\kappa)|=O(1)$ and $|S^{S}_j(\kappa)|=O(p_j-s_j)$ for every $j=1,\ldots,b$. Then
$$ \,\, \frac{|S^{L}_j(\kappa)|}{p_j}\leq \lim_{n\to\infty}\E\bigg(\frac{\hat{s}_j}{p_j}\bigg)\leq \frac{s_j-|S^{S}_j(\kappa)|}{p_j}\qquad\text{for all }j=1,\ldots,b$$
\end{prop}
An immediate consequence is that, if the betamin condition in Assumption~\hyperlink{lab:A7}{A7} also holds, then $\hat{s}_j/p_j \stackrel{L_1}{\longrightarrow} s_j/p_j$
because $S_j=S_j^{L}(\kappa)$ and $S_j^{S}(\kappa)=\emptyset$ for all $j$. The proposition also shows that, when Assumption~\hyperlink{lab:A7}{A7} is not met, $\hat{s}_j/p_j$ is asymptotically downward biased by at least $|S^{S}_j(\kappa)|/p_j$, but it is guaranteed to be larger than the proportion of signals satisfying Assumption~\hyperlink{lab:A7}{A7}. 

\subsection{Data-based block selection method} \label{ssec:databased_blockselection}

The oracle block penalties of Section~\ref{sec:properties_legression} have varying strength depending on the unknown $\beta^*_{\min,j}$ and number of active signals in each block $s_j$, and we just saw that the latter can be reliably estimated. We propose a two-step procedure. First, we use a standard (non-block-based) penalty $\kappa^{\circ}$ and estimate the number of active signals $s_j$ in each block $j$ with $\hat{s}_j$. Second, we use the estimated $s_j$ to set block penalties. The variable selection consistency of our procedure follows from the results of Section~\ref{sec:properties_legression} and holds equally for fixed and diverging $p_j-s_j$ and $s_j$. While, in Section~\ref{sec:properties_legression}, the gains in conditions for consistency were driven by how $p_j-s_j$ compared to $p-s$, here it is driven by how $p_j-|S^L_j (\kappa^{\circ})|$ compares to $p-|S^L (\kappa^{\circ})|$. This occurs because  recovery of small signals in Step 1 is not guaranteed.

The procedure has two variants. In the first variant we directly use the approximate empirical Bayes approach where $\theta^{(j)}= \hat{s}_j/p_j$ %of Lemma~\ref{lem:empiricalbayesinterpret} 
to set block penalties in Step 2. 
The second variant is motivated by Theorem~\ref{theo:suffcondlinearmodel}. Considering jointly \eqref{eq:blockpenpriorinclu} and Assumption~\hyperlink{lab:A6}{A6}, a natural choice is setting $\theta^{(j)}=(p_j-s_j+1)^{-1}$. This choice yields block penalties that are sufficiently large for Theorem~\ref{theo:suffcondlinearmodel} to hold, where $f_j$ in Assumption~\hyperlink{lab:A6}{A6} takes value $\ln(n)/2$ for all $j$. The second variant approximates this choice.\\

\textbf{Algorithm 1}
\begin{enumerate}[label=(\roman*)]
    \item Set $\kappa_j=\kappa^{\circ}=\ln(p)+\frac{1}{2}\ln(n)$ 
    for $j=1,\ldots,b$. Compute $\hat{s}_j/p_j$ in \eqref{eq:shat} for $j=1,\ldots,b$.

    \item Obtain $\hat{S}^{EB,b}$ solving \eqref{eq:Mhat} with $\kappa_j^{EB}= \ln(p_j/\hat{s}_j -1)+ \tfrac{1}{2}\ln(n)$.
    Alternatively, obtain $\hat{S}^{A,b}$ solving \eqref{eq:Mhat} with
    $\kappa_j^{A} = \ln(p_j-\hat{s}_j)+ \tfrac{1}{2}\ln(n)$.
\end{enumerate}

We refer to $\hat{S}^{EB,b}$ as the block empirical Bayes selector, and to $\hat{S}^{A,b}$ as the block adaptive selector.
Step 1 approximates posterior model probabilities under equal prior inclusion probabilities $\theta^{(j)}=(p+1)^{-1}$ in \eqref{eq:block_modelprior} across blocks, and estimates the proportion of truly active coefficients with $\hat{s}_j/p_j$. Step 1 can then be approximated by any Bayesian computational method for posterior model probabilities. In Section~\ref{sec:simulation} we use an MCMC algorithm. Step 2 selects a model using the block penalties $\kappa_j^{EB}$ induced by $\theta^{(j)}=\hat{s}_j/p_j$,
or alternatively setting $\kappa_j^{A}= \kappa_j^{EB} + \ln(\hat{s}_j)$. Any fast computational method for the exact or approximate resolution of the $\ell_0$ problem may be used for Step 2. 
Under mild assumptions both $\kappa_j^{EB}$ and $\kappa_j^{A}$ lead to variable selection consistency.
The main difference is that in non-sparse settings where $s$ grows faster than $\sqrt{n}$, the penalty $\kappa_j^{EB}$ can be insufficient and there $\kappa_j^{A}$ might be preferred.

Theorem~\ref{theo:empbayesselconsist} shows
the consistency of $\hat{S}^{EB,b}$ under
assumptions:\\[5pt]
\hypertarget{lab:A8}{(A8)}\quad $s=o(\sqrt{n})$\\[5pt]
\hypertarget{lab:A9}{(A9)}\quad For each block $j$, there exists $a_j\to \infty$ such that for every sufficiently large $n$, 
\begin{equation*}\label{cond:suffIIreg2}
        \sqrt{\frac{ (1-\psi)n\rho(\bs X)}{6}}{\beta_{\min,j}^*} \,- \,\sqrt{\ln\bigg(\frac{p_j}{|S^{L}_j(\kappa^{\circ})|}-1\bigg)+\frac12\ln(n)} \;=\; \sqrt{\ln(s_j)} + a_j .
    \end{equation*}
where $\psi=\tfrac12\Big(1+\max_j \ln(p_j-s_j)/\big(\ln(p_j/s_j-1)+\ln(n)/2\big)\Big)$.
%Using Theorem~\ref{thm:convtoT} and Proposition~\ref{prop:shatconsistence}, we show that under \hyperlink{lab:A8}{A8} and \hyperlink{lab:A9}{A9}, the $\kappa^{EB}_j$'s satisfy \hyperlink{lab:A6}{A6} and \hyperlink{lab:A7}{A7} with probability going to 1 and conclude on the selection consistency  of $\hat{S}^{EB,b}$.

\begin{thm}\label{theo:empbayesselconsist}
    Assume \hyperlink{lab:A1}{A1}, \hyperlink{lab:A8}{A8}, \hyperlink{lab:A9}{A9}, and $|S^{S}_j(\kappa^{\circ})|=O(p_j-s_j)$ for every $j=1,\ldots,b$. Then
$\lim_{n \to \infty} P(\hat{S}^{EB,b} = S)= 1$.
\end{thm}

Theorem~\ref{theo:informedl0consist} below shows
the consistency of $\hat{S}^{A,b}$. It no longer requires Assumption~\hyperlink{lab:A8}{A8}, and Assumption~\hyperlink{lab:A9}{A9} is replaced by the more stringent Assumption A10: \\[5pt]
\hypertarget{lab:A10}{(A10)}\quad For each block $j$, there exists $c_j\to \infty$ such that for  all sufficiently large $n$,
\begin{equation*}
        \sqrt{\frac{ (1-\xi)n\rho(\bs X)}{2}}\beta_{\min,j}^* \,- \,\sqrt{\ln\big(p_j-|S^{L}_j(\kappa^{\circ})|\big)+\frac12\ln(n)} \;=\; \sqrt{\ln(s_j)} + c_j .
    \end{equation*}
where  $\xi=\tfrac12(1+\max_j \ln(p_j-s_j)/(\ln(p_j-s_j)+0.5\ln(n)))$.

\begin{thm}\label{theo:informedl0consist}
    Assume \hyperlink{lab:A1}{A1}, \hyperlink{lab:A10}{A10}, and $|S^{S}_j(\kappa)|=O(p_j-s_j)$ for every $j=1,\ldots,b$, then \\
$\lim_{n \to \infty} P(\hat{S}^{A,b} = S) = 1$.
\end{thm}

\subsection{Benefits of data-based block selection}

By Theorem~\ref{theo:empbayesselconsist}, a standard (non-block-based) empirical Bayes $\hat S^{EB}$ selector that sets in Step 2 a single common penalty $\kappa^{EB}=\ln(p/\hat{s} -1)+ \tfrac{1}{2}\ln(n)$ is variable selection consistent under the betamin assumption:
\begin{equation}\label{eq:betaminEB}
        \sqrt{\frac{ (1-\psi')n\rho(\bs X)}{6}}{\beta_{\min}^*} \,- \,\sqrt{\ln\bigg(\frac{p}{|S^{L}(\kappa^{\circ})|}-1\bigg)+\frac12\ln(n)} \;=\; \sqrt{\ln(s)} + a'_j .
    \end{equation}
where $\psi'=\tfrac12\Big(1+ \ln(p-s)/\big(\ln(p/s-1)+\ln(n)/2\big)\Big)$ and $a'_j \to \infty$. We have
\begin{eqnarray*}
&\sqrt{\ln\bigg(\frac{p}{|S^{L}(\kappa^{\circ})|}-1\bigg)+\frac12\ln(n)} + \sqrt{\ln(s)} \geq \sqrt{\ln\bigg(\frac{p_j}{|S^{L}_j(\kappa^{\circ})|}-1\bigg)+\frac12\ln(n)} + \sqrt{\ln(s_j)}
\end{eqnarray*}
and also $\psi > \psi'$. Sufficient conditions for consistency and the smallest signal recoverable are then milder for $\hat S^{EB,b}$ than for $\hat S^{EB}$. %The smaller $p_j/|S^{L}_j(\kappa^{\circ})|$ is compared to $p/|S^{L}(\kappa^{\circ})|$, the milder is \hyperlink{lab:A9}{A9}  compared to \eqref{eq:betaminEB}. % in block $B_j$. 
%$\hat S^{EB,b}$ is then most beneficial when the external information identifies blocks with low proportion of inactive or small signals.

Similarly, a standard (non-block-based) adaptive selector $\hat S^{A}$ setting a common penalty $\kappa^{A}=\ln(p-\hat{s})+ \tfrac{1}{2}\ln(n)$ in Step 2 is consistent under stricter assumptions than $\hat S^{A,b}$. The smallest signal recoverable is also smaller with $\hat S^{A,b}$. %In that case, the gains are driven by the ability of external information to identify blocks with low \emph{numbers} of inactive or small signals.

%Our lack of convergence rates for $\hat S^{EB,b}$ and $\hat S^{A,b}$ precludes a detailed discussion of their benefits in this aspect. However, we remark that the $\kappa^{EB}_j$ and $\kappa^{A}_j$ are constantly lower than the single $ \kappa^{EB}$ and $\kappa^{A}$ respectively. This is in contrast to the oracle block penalties $\kappa^*_j$ that are larger in blocks with large signals. As a result, $\hat S^{EB,b}$ and $\hat S^{A,b}$ have higher probability of false positives and lower probability of false negatives than $\hat S^{EB}$ and $\hat S^{A}$ respectively.  $\hat S^{EB,b}$ and $\hat S^{A,b}$ underperform in settings with large signals and large number of inactive variables and outperform in settings with small signals. We show next with empirical results that the loss in probability of support recovery with $\hat S^{EB,b}$ and $\hat S^{A,b}$ is minor while the benefit potentially large.  

%\paul{Theorem on rates in Section~\ref{sec:properties_legression} hold when \hyperlink{lab:A6}{A6} and \hyperlink{lab:A7}{A7} hold. We can only prove \hyperlink{lab:A6}{A6} and \hyperlink{lab:A7}{A7} hold with probability going to 1. To have a rate for the EB we would need to have a rate at which \hyperlink{lab:A6}{A6} and \hyperlink{lab:A7}{A7} hold}

\section{Numerical illustrations}\label{sec:simulation}

We illustrate the performance of Algorithm 1 on simulated data. We run our proposed method in linear regression under the asymptotic regimes and block sparsity assumptions of Examples 1, 3 and 4 in Table~\ref{tab:examples}. We also consider an additional setting, Example 5, that highlights differences between $\hat S^{EB,b}$ and $\hat S^{A,b}$. In that example, we set $p=n/2$, $s=3\ln(n)$, and two blocks such that $p_1-s_1=(n-\sqrt{n})/2$, $p_2-s_2=\sqrt{n}/2$ and $s_1 = s_2=3\ln(n)/2$. In Step 1 of Algorithm 1 we take $\kappa^{\circ}=\tfrac12\ln(n)+\ln(p)$. 
To search over models, we rely on 
%We do not perform an exhaustive enumeration and scoring all of the models. We rely instead on 
the connection between $\ell_0$ penalties and Bayesian variable selection and %search over the model space through 
use the MCMC algorithm in function \texttt{bestIC} in R package \texttt{mombf}. Each visited model is 
%The models explored are 
scored with the BIC approximation \eqref{eq:blockpenpriorinclu}. In Step 2 of Algorithm 1, we obtain $\hat S^{EB,b}$, $\hat S^{A,b}$, $\hat S^{EB}$, $\hat S^{A}$ by scoring all the models visited by the MCMC in Step 1.

We simulate data with $n \in \{20,700\}$ and Gaussian covariates with unit variance and all pairwise correlations equal to 0.5.
%under Gaussian designs  with $n$ ranging from 20 to 700 and population covariance matrix:
%$$\Sigma=\begin{pmatrix}
%1\,     &\, 0.5\,  &\ldots & \,0.5\\
%0.5\,   & \,1\,    &\ldots & \,0.5\\
%\vdots  &          &       & \vdots\\
%0.5\,   &\ldots    &\,1 \, & \,0.5 \\ 
%0.5\,   &\ldots    &\,0.5\,& \,1 
%\end{pmatrix}$$
Table~\ref{tab:betaminsim} summarizes the $(\beta_{\min,1}^*,\beta_{\min,2}^*)$ %betamin assumptions 
used in our simulations. %The $\beta_{\min}^*=\beta_{\min,2}^*$ are chosen so the probability of recovery with block and standard selectors approximately spans the interval $(0,1)$ for the range of $n$ considered. 
Example 1, 3 and 5 are discriminative settings and we set $\beta_{\min,1}^*>\beta_{\min,2}^*=\beta_{\min}^*$. In Example 4, we set $\beta_{\min,1}^*=\beta_{\min,2}^*$ to represent a setting with non-discriminative blocks. Other truly active signals are drawn from the uniform distribution with support $[1,3]$.  %We use MCMC search and BIC approximation too.
\begin{table*}[ht]
\caption{Smallest active signals in simulations}
\label{tab:betaminsim}
\begin{tabular}{@{}cccccc@{}}
\hline \\[-1em]
Example & $1$ & $3$ & $4$ & $5$ \\[3pt]   
\hline\\[-1em]
$\beta_{\min,1}^*$  & 0.8  & 0.8  & 0.33  & 0.8  \\[3pt]
$\beta_{\min,2}^*$  & 0.33 & 0.33 & 0.33  & 0.2 \\
\hline
\end{tabular}
\end{table*}

\begin{figure}[ht]
\begin{minipage}{1\columnwidth}
\begin{tabular}{M{63mm}M{63mm}}
        Example 1& Example 3\\
         \includegraphics[scale=1]{images/p.probrecov.ex1.slogn.pdf} & \includegraphics[scale=1]{images/p.probrecov.ex3.slogn.pdf} \\ 
         Example 4& Example 5\\
         \includegraphics[scale=1]{images/p.probrecov.ex4.slogn.pdf}  & 
         \includegraphics[scale=1]{images/p.probrecov.ex5.slogn.pdf} 
\end{tabular}
\end{minipage}
\caption{Probability of correct selection with $\hat S^{EB,b}$ (solid black), $\hat S^{A,b}$ (dashed black), $\hat S^{EB}$ (solid grey), $\hat S^{A}$ (dashed grey), and the EBIC penalty (dotted grey) in Example 1 (top left), 2 (top right), 4 (bottom left) and 5 (bottom right)}
\label{fig:linregslogn}
\end{figure}

% \begin{figure}[ht]
% \begin{minipage}{1\columnwidth}
% \begin{tabular}{M{63mm}M{63mm}}
%         Example 1 - GSM & Example 3 - GSM \\
%          \includegraphics[scale=1]{images/p.probrecov.ex1.slogn.gsm.pdf} & \includegraphics[scale=1]{images/p.probrecov.ex3.slogn.gsm.pdf} \\ 
%          Example 4 - GSM & Example 5 - GSM  \\
%          \includegraphics[scale=1]{images/p.probrecov.ex4.slogn.gsm.pdf}  & 
%          \includegraphics[scale=1]{images/p.probrecov.ex5.slogn.gsm.pdf} \\
%          Example 5 larger penalty - GSM  & empty 
% \end{tabular}
% \end{minipage}
% \caption{Probability of correct selection with $\hat S^{EB,b}$, $\hat S^{A,b}$, $\hat S^{EB}$, $\hat S^{A}$, and the EBIC penalty in Example 1 (top left), 2 (top right), 4 (bottom left) and 5 (bottom right)}
% \label{fig:gsmslogn}
% \end{figure}

Figure~\ref{fig:linregslogn} plots the empirical probabilities of correct recovery for $\hat S^{EB,b}$, $\hat S^{A,b}$, $\hat S^{EB}$, $\hat S^{A}$ and the EBIC penalty, for Examples 1, 2, 4 and 5. The probabilities are computed over 100 simulations for each $n$.
In Examples 1 and 3 where blocks are discriminative,
%In the examples with discriminative blocks, 
$\hat S^{EB,b}$ and $\hat S^{A,b}$ outperform $\hat S^{EB}$, $\hat S^{A}$ and the EBIC, particularly for small $n$. % in the harder conditions when $n$ is small. 
%When $n$ becomes large, $\hat S^{EB,b}$ and $\hat S^{A,b}$ lose their edge over $\hat S^{EB}$ and $\hat S^{A}$ respectively. This is expected because the penalties in $\hat S^{EB,b}$ and $\hat S^{A,b}$ are smaller than in $\hat S^{EB}$ and $\hat S^{A}$. When $n$ is large, false negatives are less of an issue and there is no more benefit in smaller penalties. 
%Importantly, the underperformance of $\hat S^{EB,b}$ and $\hat S^{A,b}$ is null or minor: the probabilities of the block selectors and their standard counterparts converge to same levels in our examples. $\hat S^{EB,b}$ and $\hat S^{A,b}$ then control the probability of false positives as well as $\hat S^{EB}$ and $\hat S^{A}$ with smaller penalties: they achieve a better balancing of the probabilities of false negatives and false positives. 
In the nondiscriminative block setting, Example 4, $\hat S^{EB,b}$ and $\hat S^{A,b}$ perform very similarly to  $\hat S^{EB}$ and $\hat S^{A}$ respectively. In Examples 1, 3 and 4, the adaptive block selector $\hat S^{A,b}$ outperforms the empirical Bayes selector $\hat S^{EB,b}$ for large $n$. This occurs because $\hat{S}^{A,b}$ uses larger penalties
%is explained by the larger strength of $\hat S^{A,b}$ penalties 
that control better false positives, whereas false negatives are essentially not an issue anymore for such large $n$. In contrast, Example 5 is a setting where smaller penalties are advantageous, because the number of of inactive variables and $\beta_{\min}^*$ are small. In that case,  $\hat S^{EB,b}$ outperforms $\hat S^{A,b}$.

\section{Discussion}\label{sec:discussion}

We studied %provided an understanding of 
how incorporating external information as possible with data integration and transfer learning can facilitate model selection in the sequence model and high-dimensional linear regression. We studied the case where external information partitions
variables %external information helps partition variables 
into blocks %with potentially distinct characteristics 
and introduced corresponding block-based %externally-informed 
$\ell_0$ selectors. % that let the penalty depend on those blocks. 
We showed that an oracle externally-informed selector converges faster and under milder conditions than the standard $\ell_0$ oracle. In particular, it softens the stringent conditions on signal strength. We also provided concrete data analysis methods that incorporate external information to improve variable selection properties without requiring oracle knowledge. Efficient computation is possible for those methods via standard MCMC technique.

A question %that immediately arise 
for future work is how much the assumption of fixed number of blocks $b$ can be relaxed. Our current proof strategies %consistency results 
are robust to moderate increases in the number of blocks but do not work when $b=p$ for example. Also, our setting is motivated by situations where one has a discrete meta-covariate that allows dividing parameters into blocks, e.g. whether a variable refers to patient history or genomic biomarkers. Hence, another natural extension is to consider continuous meta-covariates, e.g. allow the prior inclusion probability of a covariate on its estimated effect in a related disease.

Another interesting research direction is understanding the benefit of external information for parameter estimation and prediction error.
For example, it is possible to obtain estimation error bounds for the sequence model, but the results depend on the chosen estimator (e.g. $\ell_0$, $\ell_1$ or $\ell_2$) and ensuring their tightness requires separate work elsewhere. By focusing on model selection, we obtained results that apply to essentially all penalties / Bayesian methods in the sequence model.


% We also observe that our results bear similarities with conclusions from on related problems. In Section \ref{sec:examples}, we discussed how the benefits of externally-informed selectors depend on how discriminative is the external information. \cite{scarlett2012compressed} similarly showed that support recovery in compressed sensing becomes easier as the entropy of the distribution of inclusion probabilities across blocks shrinks. In hypothesis testing, \cite{genovese2006} demonstrated that the power of their informed weighted testing procedure increases when assignment of weights is positively associated with the null hypotheses being false. We expect analogous results to be achievable in other related problems like structural learning in graphical models or variable selection over partially ordered sets as introduced in \cite{taeb}, which are further interesting future research avenues. 


%%%%%%%%%%%%%%%%%%%%%%%%%%%%%%%%%%%%%%%%%%%%%%
%% Single Appendix:                         %%
%%%%%%%%%%%%%%%%%%%%%%%%%%%%%%%%%%%%%%%%%%%%%%
%\begin{appendix}
%\section*{???}%% if no title is needed, leave empty \section*{}.
%\end{appendix}
%%%%%%%%%%%%%%%%%%%%%%%%%%%%%%%%%%%%%%%%%%%%%%
%% Multiple Appendixes:                     %%
%%%%%%%%%%%%%%%%%%%%%%%%%%%%%%%%%%%%%%%%%%%%%%
%\begin{appendix}
%\section{???}
%
%\section{???}
%
%\end{appendix}

%%%%%%%%%%%%%%%%%%%%%%%%%%%%%%%%%%%%%%%%%%%%%%
%% Support information, if any,             %%
%% should be provided in the                %%
%% Acknowledgements section.                %%
%%%%%%%%%%%%%%%%%%%%%%%%%%%%%%%%%%%%%%%%%%%%%%
\begin{acks}[Acknowledgments]
%The authors would like to thank the anonymous referees, an XXXXX and the XXXXX for their constructive comments. 
Paul Rognon-Vael is also affiliated with the Department of Statistics and Operational Research, Universitat Politècnica de Catalunya. David Rossell is also affiliated to the Data Science Center, Barcelona School of Economics. Piotr Zwiernik is also affiliated with the Department
of Economics and Business, Universitat Pompeu Fabra and Data Science Center, Barcelona School of
Economics.
\end{acks}
%%%%%%%%%%%%%%%%%%%%%%%%%%%%%%%%%%%%%%%%%%%%%%
%% Funding information, if any,             %%
%% should be provided in the                %%
%% funding section.                         %%
%%%%%%%%%%%%%%%%%%%%%%%%%%%%%%%%%%%%%%%%%%%%%%
\begin{funding}
PRV acknowledges the support of Departament de Recerca i Universitats de la Generalitat de Catalunya and the European Social Fund. DR was funded by grant Consolidaci\'on investigadora CNS2022-135963 by the AEI, and PID2022-138268NB-I00 by MCIN/AEI/10.13039/501100011033 /FEDER.
%Europa Excelencia EUR2020-112096 from the AEI/10.13039/501100011033 and European Union NextGenerationEU/PRT.
PZ was supported by NSERC grant RGPIN-2023-03481.
\end{funding}

%%%%%%%%%%%%%%%%%%%%%%%%%%%%%%%%%%%%%%%%%%%%%%
%% Supplementary Material, including data   %%
%% sets and code, should be provided in     %%
%% {supplement} environment with title      %%
%% and short description. It cannot be      %%
%% available exclusively as external link.  %%
%% All Supplementary Material must be       %%
%% available to the reader on Project       %%
%% Euclid with the published article.       %%
%%%%%%%%%%%%%%%%%%%%%%%%%%%%%%%%%%%%%%%%%%%%%%
\begin{supplement}
\stitle{Supplementary material to "Improving variable selection properties by levaring external data"}
\sdescription{In the supplementary material, we provide additional motivation for our estimator of sparsity, auxiliary results, proofs, properties of $\hat{S}^b$ in the Gaussian sequence model with fixed number of active signals, properties of non linear block $\ell_0$ penalties in high-dimensional linear regression, and a discussion of the tightness of our conditions for variable selection consistency in linear regression.}
\end{supplement}

%%%%%%%%%%%%%%%%%%%%%%%%%%%%%%%%%%%%%%%%%%%%%%%%%%%%%%%%%%%%%
%%                  The Bibliography                       %%
%%                                                         %%
%%  imsart-???.bst  will be used to                        %%
%%  create a .BBL file for submission.                     %%
%%                                                         %%
%%  Note that the displayed Bibliography will not          %%
%%  necessarily be rendered by Latex exactly as specified  %%
%%  in the online Instructions for Authors.                %%
%%                                                         %%
%%  MR numbers will be added by VTeX.                      %%
%%                                                         %%
%%  Use \cite{...} to cite references in text.             %%
%%                                                         %%
%%%%%%%%%%%%%%%%%%%%%%%%%%%%%%%%%%%%%%%%%%%%%%%%%%%%%%%%%%%%%

%% if your bibliography is in bibtex format, uncomment commands:
\bibliographystyle{imsart-nameyear}  % Style BST file (imsart-number.bst or imsart-nameyear.bst)
\bibliography{bibliography.bib}       % Bibliography file (usually '*.bib')

%% or include bibliography directly:
% \begin{thebibliography}{}
% \bibitem{b1}
% \end{thebibliography}
\newpage
\hspace{0.2cm}\\
\begin{center}
    \Large{Supplementary material to "Improving variable selection properties by leveraging external data"}
\end{center}

\renewcommand*\thesection{S\arabic{section}}
\renewcommand*{\theHsection}{S\arabic{section}}
\setcounter{section}{-1}
\renewcommand{\theequation}{S\arabic{equation}}
\renewcommand{\thetable}{S\arabic{table}}
\hspace{0.2cm}\\
We provide the following additional material:\\
\noindent\ref{suppsec:ebayes}: Additional motivation for our estimator of sparsity\\
\ref{suppsec:auxiliary}: Auxiliary results\\
\ref{suppsec:proofsecgsm}: Proofs of Section~\ref{sec:properties_orth}\\
\ref{suppsec:proofseclinreg}: Proofs of Section~\ref{sec:properties_legression}\\
\ref{suppsec:proofsempbayes}: Proofs of Section~\ref{sec:informed_l0}\\
\ref{suppsec:appendixGSMsfinite}: Gaussian sequence model with fixed number of active signals.\\
\ref{suppsec:appendixnonlinear}: Non linear block $\ell_0$ penalties in high-dimensional linear regression.\\
\ref{suppsec:tightness}: Tightness of conditions for variable selection consistency in linear regression.

\section{Additional motivation for our estimator of sparsity}\label{suppsec:ebayes} In Section~\ref{ssec:est.sparsity}, we presented our estimator of the proportion of truly active variables in block $j$ as the posterior mean $E(s_j \mid \bs y,\bs \theta)$. A related standard empirical Bayes strategy is to set $\hat{\theta}^{(j)}$ by %The estimator of the proportion of truly active variables $\hat{\theta}^{(j)}= \hat{s}_j/p_j$ can also be justified as an empirical Bayes estimator %procedure. 
%The empirical Bayes estimator of $\bs \theta$ (cf \eqref{eq:block_modelprior}) 
maximizing the marginal likelihood of $\bs y$ given $\bs \theta$, i.e.
\begin{eqnarray}
\hat{\bs \theta}
%\hat{\bs \theta}^{EB}
\,:=\, \arg\max_{\bs \theta} p(\bs y \mid \bs \theta)=
\arg\max_{\bs \theta} \int p(\bs y \mid \bs \beta, M) d P(\bs \beta, M \mid \bs \theta).
\nonumber
\end{eqnarray}
% In the Bayesian framework \eqref{eq:block_modelprior}, $\hat{s}_j/p_j$ corresponds to the empirical Bayes estimate of $\bs \theta$.
\begin{lem}\label{lem:empiricalbayesinterpret} For $p(M \mid \bs \theta)$ in \eqref{eq:block_modelprior}, the empirical Bayes estimator satisfies 
$$\hat{\theta}^{(j)} = \frac{1}{p_j}\sum_{i\in B_j}P(\beta_i \neq 0 \mid \bs y,\hat{\bs \theta}).$$
\end{lem}
The posterior mean estimator $\hat{\theta}^{(j)}=\hat{s}_j/p_j$ can be seen as an approximation to the fixed-point equation in Lemma \ref{lem:empiricalbayesinterpret}, where one replaces $\hat{\bs \theta}$ in the right-hand side by an initial guess (implicitly defined in Section \ref{ssec:databased_blockselection}). 

\subsection{Proof of Lemma~\ref{lem:empiricalbayesinterpret}}

Denote the marginal likelihood of $\bs y$ given $\bs \theta$ by $H(\bs \theta)=p(\bs y \mid \bs \theta)$. For every $j=1,\ldots,b$, the partial derivative of its logarithm with respect to $\theta^{(j)}$ is
    \begin{equation}\label{eq:logderivative}
        \frac{\partial \ln H(\bs \theta)}{\partial \theta^{(j)}}=\frac{\partial H(\bs \theta)}{\partial \theta^{(j)}}\,H(\boldsymbol{\theta})^{-1}.
    \end{equation}
    Observe that:
    \begin{equation}\label{eq:partialmarglikelihood}
    H(\bs \theta)=\sum_{M \in \M}p(\bs y \mid M,\bs\theta)p(M \mid \bs \theta)
    \quad \text{and}\quad 
    \frac{\partial H(\bs \theta)}{\partial \theta^{(j)}}=\sum_{M\in \M}p(\bs y \mid M,\bs\theta)\frac{\partial p(M \mid \bs \theta)}{\partial \theta^{(j)}}
    \end{equation}
    Recall that each model is defined as $M=(m_1,\ldots,m_p)$ where $m_i = I(\beta_i \neq 0)$ indicates whether variable $j$ is included under $M$, and that our choice of model prior in \eqref{eq:block_modelprior} is 
    \begin{align}
        p( M \mid \boldsymbol{\theta})=
        \prod_{j=1}^b \big(\theta^{(j)}\big)^{\sum_{i \in B_j} m_i}
        \big(1 - \theta^{(j)}\big)^{p_j - \sum_{i \in B_j} m_i}.
        \nonumber
    \end{align}
    Hence, simple algebra shows that for every $M\in \mathcal M$
    \begin{equation}\label{eq:partialprobmodel}
    \frac{\partial p(M \mid \bs \theta)}{\partial \theta^{(j)}}= p(M \mid \bs \theta) \bigg(\frac{\sum_{i\in B_j}m_i}{\theta^{(j)}}-\frac{p_j-\sum_{i\in B_j}m_i}{1-\theta^{(j)}} \bigg).       
    \end{equation}
    
    Replacing \eqref{eq:partialprobmodel} into \eqref{eq:partialmarglikelihood}, and using that for any function $f$ $$\sum_{M\in \M}\sum_{i\in B_j}m_if(M)=\sum_{i\in B_j}\sum_{M\in \M:m_i=1}f(M),$$ we get
    \begin{align}\label{eq:partialmarglikelihood2}
        \frac{\partial H(\bs \theta)}{\partial \theta^{(j)}}&\;=\;\frac{1}{\theta^{(j)}}\sum_{i \in B_j}\sum_{M\in \M: m_i=1}p(\bs y \mid M,\bs\theta)p(M \mid \bs \theta) \\&\quad- \frac{1}{1-\theta^{(j)}}\bigg[p_j H(\bs \theta)-\sum_{i \in B_j}\sum_{M\in \M: m_i=1}p(\bs y \mid M,\bs\theta)p(M \mid \bs \theta). \nonumber
        \bigg]
    \end{align}
    Note that
    \begin{align*}
        \frac{\sum_{M\in \M: m_i=1}p(\bs y \mid M,\bs\theta)p(M \mid \bs \theta)}{H(\bs \theta)} = \frac{\sum_{M\in \M: m_i=1}p(\bs y , M\mid \bs \theta)}{p(\bs y \mid \bs \theta)}= P(m_i = 1\mid \bs y,\bs \theta)
    \end{align*}
    By \eqref{eq:logderivative}, we then get the following expression of the partial derivative, for every $j=1,\ldots,b$,
    $$\frac{\partial \ln H(\bs \theta)}{\partial \theta^{(j)}}=\frac{1}{\theta^{(j)}}\sum_{i \in B_j}P(m_i = 1\mid \bs y,\bs \theta) - \frac{1}{1-\theta^{(j)}}\bigg[p_j-\sum_{i \in B_j}P(m_i = 1\mid \bs y,\bs \theta)
        \bigg]$$
    Setting the partial derivatives to 0 and solving for $\theta^{(j)}$ gives the desired result.

%\renewcommand*{\theHchapter}{\thepart.\arabic{chapter}}
\section{Auxiliary results}\label{suppsec:auxiliary}

In this section we collect some technical results.

\subsection{Solutions to penalized likelihood problems}

\begin{lem}\label{lem:l1isthreshold}
	In the sequence model \eqref{eq:GSM}, selecting the non-zero entries  of the solution to 
\begin{equation}\label{eq:genl1}
	\underset{\bs \beta \in \R^p}{\rm minimize} \;\;\tfrac12\|\bs y-\sqrt{n}\bs\beta\|^2+\sum_{j=1}^b \lambda_j \sum_{i \in B_j}|\beta_{i}|
\end{equation}
for non-negative $\lambda_1,\ldots,\lambda_b$  is equivalent to taking the selector $\hat S^{b}$ in \eqref{prop:isthresh} with $\bs \tau=(\lambda_1/n,\ldots,\lambda_b/n)$. Let $\bs\beta^\circ$ be either the MLE or the LASSO estimator with penalty $\lambda^\circ$. Selecting the non-zero entries of the solution to the problem
\begin{equation}\label{eq:genl1ada}
	\underset{\bs \beta \in \R^p}{\rm minimize} \;\;\tfrac12\|\bs y-\sqrt{n}\bs\beta\|^2+\sum_{j=1}^b \lambda_j \sum_{i \in B_j}\tfrac{|\beta_{i}|}{|\beta_i^\circ|}.
\end{equation}
is equivalent to taking the selector $\hat S^{b}$ with $\bs \tau=(\sqrt{\lambda_1/n},\ldots,\sqrt{\lambda_b/n})$ if $\bs\beta^\circ$ is the MLE and with $\bs \tau=(\lambda^\circ/2n+\sqrt{(\lambda^\circ/2n)^2+\lambda_1/n},\ldots,\lambda^\circ/2n+\sqrt{(\lambda^\circ/2n)^2+\lambda_b/n})$ if $\bs\beta^\circ$ is the LASSO estimator.
\end{lem}

\begin{proof}
    Denote by $\hat{\bs \beta}$ the minimizer of the problem \eqref{eq:genl1}. The MLE under the full model is $\tilde{\bs\beta}=\tfrac{1}{\sqrt{n}} \bs y$ and the optimized function in \eqref{eq:genl1} can be then rewritten as
	$$
	\tfrac{1}{2}\bs y^{\top}\bs y +\tfrac{n}{2}\sum_{j=1}^b\sum_{i\in B_j}\left(\beta_i^2-2\tilde\beta_i\beta_i+2\tfrac{\lambda_j}{n}|\beta_i|\right)
	$$
	which we optimize with respect to each $\beta_i$ separately. For any $i\in B_j$ we have $\hat\beta_i=0$ if and only if $|\tilde \beta_i|\leq \lambda_j/n$. Similarly  for \eqref{eq:genl1ada}, denoting $\hat{\bs \beta}$ the minimizer of the problem, for any $i\in B_j$ we have $\hat\beta_i=0$ if and only if $|\tilde \beta_i|\leq \lambda_j/(|\beta^{\circ}_i|n)$. If $\bs \beta^{\circ}= \tilde{\bs\beta}$, then for any $i\in B_j$ $\hat{\beta}_i = 0$ if and only if $|\tilde \beta_i|\leq \sqrt{\lambda_j/n}$. If $\bs\beta^\circ$ is a LASSO estimate with penalization $\lambda^\circ$, then for any $i\in B_j$ 
 $\hat{\beta}_i = 0$ if and only if $\tilde{\beta}^2_i + \operatorname{sign}(\lambda^{\circ}/n-\tilde{\beta}_i)\lambda^\circ/(n\tilde{\beta}_i)-\lambda_j/n \leq 0$. Equivalently, for any $i\in B_j$ 
 $\hat{\beta}_i = 0$ if and only if $|\tilde{\beta}_ i| \leq  \tfrac{1}{2}\left(\lambda^\circ/n+\sqrt{(\lambda^\circ/n)^2+4 \lambda_j/n}\right)$.
\end{proof}

\subsection{Tail bounds}

\begin{lem}\label{lemma:3tailbounds}\textbf{Tail bounds on the maximum and minimum of folded Gaussians} 
\begin{enumerate}[label=(\roman*)]
        \item If $\bs y\sim N_p(0,n^{-1}I_p)$, $p>1$, and $a \geq \sqrt{2\ln(p)/n}$, $$P\Big(\underset{i\in\{1,\ldots,p\}}{\max}|y_i|> a \Big) \leq \frac{e^{-\frac{n}{2}\left(a^2-\frac{2\ln (p)}{n}\right)}}{\sqrt{\pi\ln (p)}}.$$
    \item If $\bs y\sim N_s\left(\bs\mu,\sigma^2I_s\right)$, $s>1$, and $a \leq \underset{i\in\{1,\ldots,s\}}{\min}|\mu_i|$,
        $$P\Big(\underset{i\in\{1,\ldots,s\}}{\min}|y_i| > a\Big) \geq P\Big(\underset{i\in\{1,\ldots,s\}}{\min}|\mu_i|-\underset{i\in\{1,\ldots,s\}}{\max}|y_i-\mu_i| > a \Big).$$
    \item If $y_i\sim N\left(\mu,\sigma^2\right)$ for $i=1,\ldots,s$ are independent, and $0<a<|\mu|$, then
        $$P\Big(\underset{i\in\{1,\ldots,s\}}{\min}|y_i| > a \Big) \leq \exp\left\{-\frac{s}{2}\frac{e^{-2\sigma^{-2}(|\mu|-a)^2/\pi}-e^{-\sigma^{-2}(a+|\mu|)^2/2}}{\left(1-e^{-2\sigma^{-2}(|\mu|-a)^2/\pi}\right)^{\frac{1}{2}}+\left(1-e^{-\sigma^{-2}(a+|\mu|)^2/2}\right)^{\frac{1}{2}}}\right\}.$$
\end{enumerate}
\end{lem}

\begin{proof}

\textbf{Part (i)} \,\, By the union bound and the identical distribution of the $y_i$'s,
\begin{center}
    $P\Big(\underset{i\in\{1,\ldots,p\}}{\max}|y_i|> a \Big)     \leq \sum_{i=1}^{p}{ P\left(|\sqrt{n}y_i|>\sqrt{n}a \right)}
    = p\; P\left(|z|>\sqrt{n}a\right)$
\end{center}
    where $z\sim N(0,1)$. 
    By symmetry and using the standard tail bound for standard normal, $P(z \geq \delta) \leq (2\pi)^{-1/2} \delta^{-1} e^{-\delta^{2} / 2}$ for $\delta \geq 0$, we obtain that 
\begin{center}
  $P\Big(\underset{i\in\{1,\ldots,p\}}{\max}|y_i|> a \Big) 
     \leq  p\; 2 \, P\left(z>\sqrt{n}a\right) \leq \frac{1}{\sqrt{\pi}}\frac{\sqrt{2}}{\sqrt{n}a}p e^{-\frac{n}{2} a^{2}}.$
\end{center}
Since $a \geq \sqrt{2\ln(p)/n}$, we have that $\sqrt{2}/\sqrt{n}a\leq 1/\sqrt{\ln(p)}$. Taking $a^2=a^2-\frac{2\ln(p)}{n}+\frac{2\ln(p)}{n}$, we get \begin{center}
$ P\Big(\underset{i\in\{1,\ldots,p\}}{\max}|y_i|> a \Big) \leq \frac{1}{\sqrt{\pi\ln(p)}} e^{-\frac{n}{2} \left(a^{2}-\frac{2\ln(p)}{n}\right)}.$ \\\end{center}

\textbf{Part (ii)} Consider the events $A:=\{ \underset{i \in \{1,\ldots,s\}}{\min}|\mu_i|- \underset{i \in \{1,\ldots,s\}}{\max}|y_i-\mu_i| > a\}$ and $B:=\{\underset{i \in \{1,\ldots,s\}}{\min}|y_i| > a\}$. We first show that $A$ implies $B$. $A$ implies that, for all $i$, $(\underset{i \in \{1,\ldots,s\}}{\min}|\mu_i|) - |y_i-\mu_i| > a $.  By the triangle inequality, we have that $ -|y_i-\mu_i|\leq|y_i|-|\mu_i|$ and therefore that $$ \text{for all } i, \;\;\;(\min_{i \in \{1,\ldots,s\}}|\mu_i|) - |y_i-\mu_i| \leq  (\min_{i \in \{1,\ldots,s\}}|\mu_i|) +|y_i| - |\mu_i| < |y_i|.$$ 
Then, $A$ 
implies that, for all $i$, $|y_i| > a $, that is $B$. By the monotonicity of the probability $P$
\begin{center}$ P\Big(\underset{i \in \{1,\ldots,s\}}{\min}|\mu_i|-\underset{i \in \{1,\ldots,s\}}{\max}|y_i-\mu_i| > a\Big) \leq P\Big(\underset{i \in \{1,\ldots,s\}}{\min}|y_i| > a\Big) $. \end{center}


\textbf{Part (iii)}\,\, 
Since the $y_i$'s are independent and identically distributed we have $P(\min_{i\in\{1,\ldots,s\}}|y_i|>a)=P(|y_i|>a)^s= \exp\left(s\ln{P(|y_i|>a)}\right)$ for any $i\in\{1.\ldots,s\}$. Using that $\ln\left(1+ x \right)< x$ for $x\in(-1,0)$, we have
\begin{equation}\label{eq:grpoih}
    P(\min_{i\in\{1,\ldots,s\}}|y_i|>a)\leq \exp\left(s(P(|y_i|>a)-1)\right).
\end{equation}
To get a bound on $P(\min_{i\in\{1,\ldots,s\}}|y_i|>a)$ it is then enough to bound $P(|y_i|>a)$ for any $i$. We have
$$P(|y_i|>a) = P(y_i>a) + P(y_i<-a) = P(z>\frac{a-\mu}{\sigma}) + P(z<-\frac{a+\mu}{\sigma})$$
where $z\sim N(0,1)$. If $\mu>0$, $P(z>\frac{a-\mu}{\sigma})  =  P(z>\frac{a-|\mu|}{\sigma})= P(z<\frac{|\mu|-a}{\sigma})$ by symmetry of the standard Gaussian, and $P(z<-\frac{a+\mu}{\sigma})=P(z<-\frac{a+|\mu|}{\sigma})$, then $P(|y_i|>a) =  P(z<\frac{|\mu|-a}{\sigma}) + P(z<-\frac{a+|\mu|}{\sigma})$.
If $\mu<0$, then $P(z>\frac{a-\mu}{\sigma})=P(z>\frac{a+|\mu|}{\sigma})=P(z<-\frac{a+|\mu|}{\sigma})$, by symmetry of the standard Gaussian, and $P(z<-\frac{a+\mu}{\sigma})=P(z<\frac{|\mu|-a}{\sigma})$. Then for any $\mu$, $P(|y_i|>a)  = P(z<\frac{|\mu|-a}{\sigma}) + P(z<-\frac{a+|\mu|}{\sigma}) $. For any $a<|\mu|$ we further have
\begin{align}
    P(|y_i|>a) &= P\Big(z<\frac{|\mu|-a}{\sigma}\Big) + 1- P\Big(z<\frac{a+|\mu|}{\sigma}\Big) \nonumber \\
    &= P(z<0)
    +P\Big(0<z<\frac{|\mu|-a}{\sigma}\Big)+ 1 - P(z<0) - P\Big(0<z<\frac{a+|\mu|}{\sigma}\Big) \nonumber \\
    &= 1+P\Big(0<z<\frac{|\mu|-a}{\sigma}\Big)-P\Big(0<z<\frac{a+|\mu|}{\sigma}\Big)\label{eq:random}.
\end{align}
From \cite{chu}, for any $\delta>0$, $\frac{1}{2}\left(1-e^{-\delta^2/2}\right)^{\frac{1}{2}}\leq P(0 <z < \delta)\leq \frac{1}{2}\left(1-e^{-2\delta^2/\pi}\right)^{\frac{1}{2}}$. Thus, for any $\gamma,\delta>0$, we have that
$$ P(0<z \leq \gamma)-P(0<z \leq \delta)\leq
\frac{1}{2}\left(1-e^{-2\gamma^2/\pi}\right)^{\frac{1}{2}}-\frac{1}{2}\left(1-e^{-\delta^2/2}\right)^{\frac{1}{2}}.$$
Applying the above to \eqref{eq:random}, we get
$$P(|y_i|>a) \leq 
1+ \frac{1}{2}\left(\left(1-e^{-2(|\mu|-a)^2/\pi\sigma^2}\right)^{\frac{1}{2}}-\left(1-e^{-(a+|\mu|)^2/2\sigma^2}\right)^{\frac{1}{2}}\right).\\$$

Using that $x^{\frac{1}{2}}-y^{\frac{1}{2}}=\frac{x-y}{x^{\frac{1}{2}}+y^{\frac{1}{2}}}$ for $x,y>0$, we get
 \begin{equation}\label{eq:pafvih}
    P(|y_i|>a) \leq 1- \frac{e^{-2(|\mu|-a)^2/\pi\sigma^2}-e^{-(a+|\mu|)^2/2\sigma^2}}{2\left(\left(1-e^{-2(|\mu|-a)^2/\pi\sigma^2}\right)^{\frac{1}{2}}+\left(1-e^{-(a+|\mu|)^2/2\sigma^2}\right)^{\frac{1}{2}}\right)}. 
 \end{equation}
Finally, inputting in the bound in \eqref{eq:pafvih} in \eqref{eq:grpoih} gives the desired inequality.
\end{proof}

\begin{lem}\label{lem:boundrhogen}
 For any $T\subseteq S$ and $M\in\M$ such that $T\not\subseteq M$, let $Q_T=M \cup T$. The non-centrality parameter defined in \eqref{eq:muQM} satisfies: 
\begin{equation}
    \mu_{Q_TM}\; \geq \;n\,\rho(\bs X)\,\sum_{j=1}^b|T_j\setminus M_j|\,\min_{i\in T_j\setminus M_j}{\beta_{i}^*}^2.
\end{equation}
\end{lem}

\begin{proof} 
The non-centrality parameter $\mu_{Q_TM}$, as defined in \eqref{eq:muQM}, satisfies
\begin{eqnarray*}
    \mu_{Q_TM}&=&\|\big(I_n-P_M\big) \bs X_{Q_T\setminus M} \bs \beta^*_{Q_T\setminus M}\|^2\\
    &=&\|\big(I_n-P_M\big) \bs X_{T\setminus M} \bs \beta^*_{T\setminus M}\|^2\\
    &= & n {\beta^*_{T\setminus M}}^\top \left(\tfrac1n \bs X_{T\setminus M}^\top(I_n-P_M)\bs X_{T\setminus M}\right)\beta^*_{T\setminus M}\\
    &\geq & n\lambda_{\min}\left(\tfrac1n \bs X_{T\setminus M}^\top(I_n-P_M)\bs X_{T\setminus M}\right)\|\beta^*_{T\setminus M}\|^2,
\end{eqnarray*}
where the second equality follows from observing that $Q_T\setminus M=T\setminus M$. Since $T$ is a subset of $S$, by reordering columns, $X_{S\setminus M}=[X_{T\setminus M},X_{S\setminus (M\cup T)}]$ and therefore $\tfrac1n \bs X_{T\setminus M}^\top(I_n-P_M)\bs X_{T\setminus M}$ is a principal submatrix of $\tfrac1n \bs X_{S\setminus M}^\top(I_n-P_M)\bs X_{S\setminus M}$. Hence, Cauchy's interlacing theorem gives that $\lambda_{\min}\left(\tfrac1n \bs X_{T\setminus M}^\top(I_n-P_M)\bs X_{T\setminus M}\right) \geq \lambda_{\min}\left(\tfrac1n \bs X_{S\setminus M}^\top(I_n-P_M)\bs X_{S\setminus M}\right)$. Finally, by definition of $\rho(\bs X)$ in \eqref{eq:rhoX} we have that
$\lambda_{\min}\left(\tfrac1n \bs X_{S\setminus M}^\top(I_n-P_M)\bs X_{S\setminus M}\right) \geq \rho(\bs X)$, and further noting that $\|\beta^*_{T\setminus M}\|^2\geq \sum_{j=1}^b|T_j\setminus M_j|\,\min_{i\in T_j\setminus M_j}{\beta_{i}^*}^2$ gives the desired result.
\end{proof}

\begin{lem}\label{lemma:s1s3}
Let $W\sim\chi_{\nu}^2(\mu)$ with $\mu \geq 0$, then for any $w>\mu+\nu$
$$P(W > w) \leq e^{-\big(\frac{w+\mu}{2}  - \sqrt{2w( 2 \mu+ \nu)-2 \mu \nu - \nu^2 } \big)}. $$
Moreover, assume $w$, $\nu$ and $\mu$ are functions of $n$ such that $w$ is increasing, $\nu=o(w)$, and $\mu=o(w)$. Then, for any $\phi\in(0,1)$ and $n$ large enough
$$P(W > w)  \leq  e^{-\phi\frac{w}{2}}.$$
\end{lem}

\begin{proof}
By \cite{birge}, Lemma 8.1 we have that for any $x>0$
$$P(W> (\nu+\mu)+2 \sqrt{(\nu+2 \mu) x}+2 x) \leq e^{-x}.$$
The function $f:x\mapsto (\nu+\mu)+2 \sqrt{(\nu+2 \mu) x}+2 x$ is  one-to-one between $\R^+$ and $( \nu+
\mu,\infty)$. 
Hence, we have that for any $w>\mu+\nu$,
$$P(W > w) \leq e^{-f^{-1}(w)} \:=\; e^{-\big(\frac{w+\mu}{2}  - \sqrt{2w(2 \mu+\nu)-2 \mu \nu - \nu^2 }\big) } .$$
Observe that
$$\frac{w+\mu}{2}  - \sqrt{2w( 2 \mu+ \nu)-2 \mu \nu - \nu^2 }= \frac{w}{2}\bigg(1+\frac{\mu}{w} -\sqrt{\frac{8(2 \mu+\nu)}{w}\bigg(1-\frac{2\nu\mu-\nu^2}{2(2w\mu+w\nu)} \bigg)}\bigg).$$
 Since $\nu=o(w)$ and $\mu=o(w)$ by assumption, we have $\frac{w+\mu}{2}  - \sqrt{2w( 2 \mu+ \nu)-2 \mu \nu - \nu^2 }= \frac{w}{2}(1+o(1))$. Therefore, for any $\phi\in(0,1)$ and every $n$ large enough.
$$P(W > w) \leq  e^{-\phi\frac{w}{2}}.$$
\end{proof}


\begin{lem}\label{lemma:s2} Let $W \sim \chi_\nu^2(\mu)$ with $\mu>0$. For any $w<\mu$, 
$$
P(W<w) \leq \frac{e^{-\frac{1}{2}(\sqrt{\mu}-\sqrt{w})^2}}{(\mu / w)^{\nu / 4}}.$$
\end{lem}

\begin{proof}
The result follows directly from \cite{rossell2022concentration}, Lemma S2.
\end{proof}

\begin{lem}\label{lemma:s18} 
    Let $W\sim \chi_\nu^2(\mu)$ with $\mu \geq 0$. Assume that $g$, $\nu$ and $\mu$ are functions of $n$ such that $g$ is positive and increasing,  $\nu=o(\ln(g))$, and $\mu=o(\ln(g))$. Let $\bar{u},\underline{u}$ in $(0,1)$ such that $1>\bar{u}>\underline{u}\geq\left(1+g^{\phi}e^{-(\nu+\mu)/2}\right)^{-1}$ where $\phi\in(0,1)$, then for every $n$ large enough, we have 
$$
\int_{\underline{u}}^{\bar{u}} P\left(W>2\ln \left(\frac{g}{1 / u-1}\right)\right) d u \leq  \frac{1}{g^{\phi}}\Big(\bar{u}-\underline{u}+\ln\Big(\frac{\bar{u}}{\underline{u}}\Big)\Big).
$$
\end{lem}

\begin{proof}
For any $u\in [\underline{u},\bar{u}]$, we have $2\ln \left(\frac{g}{1 / u-1}\right) \geq 2\ln \left(\frac{g}{1 / \underline{u}-1}\right)$. Since $\underline{u}\geq (1+g^{\phi}e^{-(\nu+\mu)/2})^{-1}$ by assumption  we also have that, for any $u\in (\underline{u},\bar{u})$,
$$2\ln \left(\frac{g}{1 / u-1}\right) \geq 2(1-\phi)\ln(g)+\nu+\mu.$$
It follows,
$$\frac{\nu}{2\ln \left(\frac{g}{1 / u-1}\right)}\leq \frac{\nu}{2(1-\phi)\ln(g)+\nu+\mu}
\quad\text{and}\quad
\frac{\mu}{2\ln \left(\frac{g}{1 / u-1}\right)}\leq\frac{\mu}{ 2(1-\phi)\ln(g)+\nu+\mu}.$$
By assumption $\nu=o\left(\ln (g)\right)$ and $\mu=o\left(\ln (g)\right)$, then for any $u\in (\underline{u},\bar{u})$, $\nu=o\left(2\ln \left(\frac{g}{1 / u-1}\right)\right)$ and $\mu=o\left(2\ln \left(\frac{g}{1 / u-1}\right)\right)$. By Lemma~\ref{lemma:s1s3}, for any $\phi\in(0,1)$ and every $n$ large enough,
\begin{equation}\label{eq:randomm}
    \int_{\underline{u}}^{\bar{u}} P\left(W>2\ln \left(\frac{g}{1 / u-1}\right)\right) d u< \frac{1}{g^{\phi}} \int_{\underline{u}}^{\bar{u}}(1 / u-1)^{\phi} d u .
\end{equation}
Applying the change of variables $v=1 / u-1$ to the integral on the right-hand side above gives
\begin{equation}\label{eq:rpoig}
    \int_{\underline{u}}^{\bar{u}}(1 / u-1)^{\phi} d u=\int_{1 / \bar{u}-1}^{1 / \underline{u}-1} \frac{v^{\phi}}{(v+1)^2} d v.
\end{equation}
Rewrite $v^{\phi}=(v-1 + 1)^{\phi}$. Since $\bar{u}<1$, we have that for any $v>1 / \bar{u}-1$, $v-1>-1$. Note that for any $x \geq-1$ and $r\in[0,1]$ $(1+x)^r \leq 1+r x$. Then, for any $v>1 / \bar{u}-1$, $v^{\phi}=(v-1 + 1)^{\phi}\leq 1 + \phi (v-1)\leq 1 + \phi(v+1)$. Applying this last inequality to the right-hand side in \eqref{eq:rpoig} gives
\begin{equation}\label{eq:random2}
\int_{\underline{u}}^{\bar{u}}(1 / u-1)^{\phi} d u<\int_{1 / \bar{u}-1}^{1 / \underline{u}-1} \frac{1}{(v+1)^2} +\frac{\phi}{v+1} d v = \bar{u}-\underline{u}+\phi\ln\Big(\frac{\bar{u}}{\underline{u}}\Big).
\end{equation}
The result follows inputing the bound from \eqref{eq:random2} in \eqref{eq:randomm} and using that $\phi<1$ and $\ln(\bar{u}/\underline{u}) \geq 0$ ($\underline{u} \leq \bar{u}$).
\end{proof}

\begin{lem}\label{lemma:s20}
 Let $W\sim \chi_\nu^2(\mu)$ with $\mu \geq 0$. Assume that $g$, $\nu$ and $\mu$ are functions of $n$ such that $g$ is positive and increasing,  $\nu=o(\ln(g))$, and $\mu=o(\ln(g))$. Then for any $\alpha\in(0,1)$ and every large enough $n$, we have
$$
\int_0^1 P\left(W>2 \ln \left(\frac{g}{1 / u-1}\right)\right) d u = o\big(  g^{-\alpha}\big).$$
\end{lem}

\begin{proof}
Since a probability is bounded by $1$, for any $a\in(0,1)$,
\begin{equation}\label{eq:basebound}
    \int_0^1 P\left(W>2 \ln \left(\frac{g}{1 / u-1}\right)\right) d u \;\leq\; 2 a+\int_{a}^{1-a} P\left(W>2\ln\left(\frac{g}{1 / u-1}\right)\right) d u.
\end{equation}

Take $a=\left(1+g^{\phi}e^{-(\nu+\mu)/ 2}\right)^{-1}$ for some $\phi\in(\alpha,1)$. By Lemma~\ref{lemma:s18} with $\underline{u}=a$ and $\bar{u}=1-a$, we have that
$$
\int_0^1 P\left(W>2 \ln \left(\frac{g}{1 / u-1}\right)\right) d u \;\leq\; \frac{2}{1+g^{\phi}e^{-(\nu+\mu)/2}}+ 
\frac{1-2a+\ln(g^{\phi}e^{-(\nu+\mu)/2})}{g^{\phi}}. $$
Since $1+g^{\phi}e^{-(\nu+\mu)/ 2}>g^{\phi}e^{-(\nu+\mu)/ 2}$ and $1-2a-\frac{\nu+\mu}{2}\leq \ln(g^\phi)$ for every $n$ large enough, we have
$$
\int_0^1 P\left(W>2 \ln \left(\frac{g}{1 / u-1}\right)\right) d u \;\leq\; g^{-\phi}\big(2e^{(\nu+\mu)/ 2}+2\ln(g^{\phi})\big) $$
We have that
\begin{equation}\label{eq:random24}
    \frac{g^{-\phi}2e^{(\nu+\mu)/ 2}}{g^{-\alpha}} 
= e^{-(\phi-\alpha) \ln(g) + \ln(2) + (\nu + \mu)/2}
    \;=\;e^{-(\phi-\alpha)\ln(g)\big(1-\frac{\nu+\mu}{2(\phi-\alpha)\ln(g)}- \frac{\ln(2)}{(\phi-\alpha)\ln(g)}\big)}
\end{equation}
and similarly that
\begin{equation}\label{eq:random25}
    \frac{g^{-\phi}2\ln(g^{\phi})}{g^{-\alpha}} \;=\;e^{-(\phi-\alpha)\ln(g)\big(1-\frac{\ln(\phi\ln(g))}{(\phi-\alpha)\ln(g)}- \frac{\ln(2)}{(\phi-\alpha)\ln(g)}\big)}.
\end{equation}
Since $\alpha<\phi$ as stated above, and by assumption $g$ is increasing, $\nu=o(\ln(g))$ and $\mu=o(\ln(g))$, both expressions in \eqref{eq:random24} and \eqref{eq:random25} vanish as $n$ grows. Hence, 
$$\int_0^1 P\left(W>2 \ln \left(\frac{g}{1 / u-1}\right)\right) d u \;=\; o(g^{-\alpha}).$$
\end{proof}

\subsection{A general necessary condition on signal strength in the Gaussian sequence model}
Lemma~\ref{lem:gennecIIblockthresh} gives a necessary condition on signal strength for support recovery with $\hat{S}^b$ that applies independently on whether the $s_j$ are fixed or diverging. It is analogous to a necessary condition for recovery with $\hat{S}$ shown in \cite{abraham2023sharp}.
\begin{lem}\label{lem:gennecIIblockthresh} In the
sequence model \eqref{eq:GSM}, assume \hyperlink{lab:A1}{A1} and \hyperlink{lab:A2}{A2}.
%\begin{enumerate}[label=(\roman*)]
%    \item 
Suppose that $\tau_j<\beta^*_{\min,j}$ satisfies $\lim_{n\to \infty} \sqrt{n}\tau_j/\sqrt{2\ln(p_j-s_j)}\;\geq\;1$ for some $j \in \{1,\ldots,b\}$. If \begin{equation}\label{cond:necIIallblockthresh2}
    \lim_{n\to \infty} \sqrt{n}(\beta_{\min,j}^*-\tau_j)\; <\; \infty
\end{equation}
then $\lim_{n \rightarrow \infty} P(\hat S^{b}\supseteq S  ) < 1$.
% \item Suppose there exists $j\in \{1,\ldots,b\}$ such that  $\beta^*_i=\beta_{\min,j}^*$ for all $i \in S_j$, $\tau_j<\beta^*_{\min,j}$ satisfies \eqref{cond:necIallblockthresh}, and  $s_j/p_j<1$. If Assumption~\hyperlink{lab:A3}{A3} holds and, in addition, 
% \begin{equation}\label{cond:necIIallblockthresh}
% \lim_{n\to \infty} 	\frac{\beta_{\min,j}^*-\tau_j}{\sqrt{\frac{\pi}{2}\frac{\ln(s_j)}{n}}} \;\leq\;  1\end{equation}
% then $\lim_{n \rightarrow \infty} P(\hat S^{b}\supseteq S ) < 1$.
%\end{enumerate}
\end{lem}

\begin{proof}
By independence we have that
$$P(\hat{S}^{b}\supseteq S)=\prod_{j=1}^b P(\min_{i \in S_j}|y_i/\sqrt{n}| > \tau_j).$$ 
where
$$
P(\min_{i \in S_j}|y_i/\sqrt{n}| > \tau_j)=\prod_{i \in B_j} P(|y_i/\sqrt{n}| > \tau_j). 
$$
Take any $j$, satisfying \eqref{cond:necIIallblockthresh2}. 
Denote by $i_{\circ}\in B_j$ an entry such that $|\beta_{i_{\circ}}^*|=\beta^*_{\min,\,j}$. In the proof of Lemma~\ref{lemma:3tailbounds} (iii), we show that if $y\sim N(\mu, \sigma^2)$, then for any $a<|\mu|$ , $P(|y|>a)  = P(z>\frac{a-|\mu|}{\sigma}) + P(z<-\frac{a+|\mu|}{\sigma})$ where $z\sim N(0,1)$. Using the latter, we have
$$P(|y_{i_{\circ}}/\sqrt{n}|>\tau_j) 
    = P\left(z>\sqrt{n}\left(\tau_j-\beta_{\min,\,j}^*\right)\right) + P\left(z<-\sqrt{n}\left(\tau_j+\beta_{\min,\,j}^*\right)\right)$$ where $z\sim N(0,1)$. Since $\tau_j<\beta^*_{\min,j}$ satisfies $\lim_{n\to \infty} \sqrt{n}\tau_j/\sqrt{2\ln(p_j-s_j)}\;\geq\;1$, $\sqrt{n}\left(\tau_j+\beta_{\min,\,j}^*\right)\to\infty$ and we have that $P(z<-\sqrt{n}\big(\tau_j+\beta_{\min,\,j}^*\big)) \to 0$.
    
    Further, by \eqref{cond:necIIallblockthresh2}  we have that $\underset{n \to \infty}{\lim}P\left(z>\sqrt{n}\left(\tau_j-\beta_{\min,\,j}^*\right)\right)<1$ and hence we get $\underset{n \to \infty}{\lim}P\left(\hat{S}^{b} \supseteq S \right)<1$, as we wished to prove.
\end{proof}

\subsection{Bounds related to model normalized scores}

Let $NC(M)$ be the normalized score for model $M$ defined in \eqref{eq:CM}, $\M$ the set of models under consideration, and $\mu_{Q M}$ be the noncentrality parameter for any two nested models $Q \supseteq M$ defined in Lemma~\ref{lemma:s7}. Lemma~\ref{lem:L0toL1convergencegen} generalizes Lemma~\ref{lem:L0toL1convergence} and shows that the probability of not selecting a set of models is bounded above by the expected sum of the normalized scores of the models outside the set. Lemma~\ref{lem:boundexpnc} provides a bound on the expected normalized score of any $M \in \M$ based on the pairwise comparison $C(T)-C(M)$ (\emph{cf} \eqref{eq:CM}) for some $T\subseteq S$. It essentially shows that, if the block penalties diverge and the signals in $M\setminus T$ are small, $NC(M)$ is small.  Lemma~\ref{lem:upperboundnoncentralgen} gives, for any $T\subseteq S$ and $M\in \M$, upper and lower bounds on $\mu_{Q_TT}$ where $Q_T=M\cup T$.

\begin{lem}\label{lem:L0toL1convergencegen} For $\hat S^b$ as in \eqref{eq:Mhat} and any $k$ models $M_1,\ldots,M_k$
$$P(\hat{S}^b\not\in \{M_1,\ldots,M_k\})\;\leq\; (k+1) \sum_{M \in \M\setminus \{M_1,\ldots,M_k\}}\E\left( NC(M)\right).$$
\end{lem}

\begin{proof}
Suppose that $NC(M_1)+\ldots+NC(M_k)>\frac{k}{k+1}$ then for any $M\not\in \{M_1,\ldots,M_k\}$, we have that
$$NC(M)= 1-\sum_{M' \neq M} NC(M') < 1 - \sum_{M' \in \{M_1,\ldots,M_k\}} NC(M') < \frac{1}{k+1}.$$ 
In addition, if $NC(M_1)+\ldots+NC(M_k)>\frac{k}{k+1}$ then necessarily $\max_{l=1,\ldots,k}NC(M_l)> \frac{1}{k+1} > NC(M) $ for any $M \not\in \{M_1,\ldots,M_k\}$, and therefore $\hat{S}^{b}\in \{M_1,\ldots,M_k\}$. Consequently, 
$$
P(\hat{S}^b\not\in \{M_1,\ldots,M_k\})\; \leq\; P\Big(NC(M_1)+\ldots+NC(M_k)\leq\frac{k}{k+1}\Big).$$
Moreover, we have
$$P\Big(NC(M_1)+\ldots+NC(M_k)\leq\frac{k}{k+1}\Big)\;=\;P\Big(\sum_{M \in \M\setminus \{M_1,\ldots,M_k\}} NC(M)\geq \frac{1}{k+1}\Big).$$
The result follows from the Markov's inequality applied to the right-hand side above.
\end{proof}

\begin{lem}\label{lem:boundexpnc} For any $T\subseteq S$ and $M\in \M\setminus\{T\}$, denote $Q_T= M\cup T$ and $A_T:=\gamma\Delta_{MT}+\tfrac{1-\gamma}{6}\mu_{Q_TM}$ (\emph{cf} \eqref{eq:DeltaM} and \eqref{eq:muQM}). Suppose that, for some $\gamma \in (1/2,1]$, it holds that $A_T>0$, $|M\setminus T|=o(A_T)$, and $\mu_{Q_T T}=o(A_T)$. For any $\psi\in(0,1)$ and every $n$ large enough,
$$\E\left(NC(M)\right)  \;\;\leq\;\; e^{-\psi A_T}.$$
\end{lem}

\begin{proof}
Since $M\in \mathcal M\setminus \{T\}$, by Lemma~\ref{lem:L0toL1convergence} (ii),  $NC(M)\leq (1+e^{C(T)-C(M)})^{-1} \in [0,1]$.  
The first step of the proof is to use that for any random variable $Z \geq 0$ we have $\E(Z)=\int_0^\infty P(Z>u) du$, so that  %By the standard formula for expectation of a nonnegative random variable in terms of its survival function, we get that
\begin{eqnarray*}
  \E\left(NC(M)\right)  &\leq& \int_0^1 P\left((1+e^{C(T)-C(M)})^{-1}\geq u\right){\rm d} u\\
  &=&\int_0^1 P\left(C(T)-C(M) \leq \ln{(\tfrac1u-1)}\right){\rm d}u\\
  &=&\int_0^1 P\Big(-\tfrac12 L_{TM} \geq \Delta_{MT}-\ln{(\tfrac1u-1)}\Big) {\rm d}u.
\end{eqnarray*}

The second step of the proof is to use the union bound to upper bound the probability in the integrand above.
Let $Q_T=M \cup T$ and recall that $L_{TM}=L_{Q_TM}-L_{Q_TT}$. For any $\gamma$ 
$$
-\tfrac12 L_{TM} -(\Delta_{MT}-\ln{(\tfrac1u-1)})\;=\;\left(\tfrac12 L_{Q_TT}-\ln{\Bigl(\frac{e^{\gamma\Delta_{MT}}}{\tfrac1u-1}\Bigl)}\right)-\Bigl(\tfrac12 L_{Q_TM}+(1-\gamma)\Delta_{MT}\Bigl).
$$
Observe that for any random variables $U$, $V$, and any $\epsilon,\gamma' \geq 0$, the event $
\{U-V \geq 0\}$ implies $\{U \geq \gamma' \epsilon\}\cup
\{V  < \gamma' \epsilon\}$. Let $U=\tfrac12 L_{Q_TT}-\ln{\Bigl(\frac{e^{\gamma\Delta_{MT}}}{1/u-1}\Bigl)}$ and $V=\tfrac12 L_{Q_TM}+(1-\gamma)\Delta_{MT}$. Take $\epsilon=\mu_{Q_TM}$ and $\gamma'=\tfrac16({1-\gamma})$, and observe that
$A_T:=\gamma\Delta_{MT}+\tfrac{1-\gamma}{6}\mu_{Q_TM}= \gamma\Delta_{MT}+\gamma'\mu_{Q_TM}$. We then have
\begin{align*}
\{U\geq \gamma'\epsilon\}\;&=\;\{\tfrac12 L_{Q_TT}\geq \ln{\Bigl(\frac{e^{\gamma\Delta_{MT}}}{\tfrac1u-1}}\Bigl)+\gamma'\mu_{Q_TM}\}\;=\;\{\tfrac12 L_{Q_TT}\geq \ln{\Bigl(\frac{e^{A_T}}{\tfrac1u-1}}\Bigl)\}\\
\{V<\gamma'\epsilon\}\;&=\;\{\tfrac12 L_{Q_TM}<-(1-\gamma)\Delta_{MT}+\gamma'\mu_{Q_TM}\}\;=\;\{\tfrac12 L_{Q_TM}<\gamma'(\mu_{Q_TM}-6\Delta_{MT})\}.
\end{align*}
By the union bound we have that
\begin{equation}\label{eq:prooflemmas20nonspu}
    \E\left(NC(M)\right)  \leq \int_0^1P\Bigl(\tfrac12 L_{Q_TT}\geq \ln{\Bigl(\frac{e^{A_T}}{\tfrac1u-1}\Bigl)}\Bigl){\rm d}u+ P\Bigl(\tfrac12 L_{Q_TM}<\gamma'(\mu_{Q_TM}-6\Delta_{MT})\Bigl).
\end{equation}

The third and final step of the proof is to upper bound each of the terms in the right-hand side of \eqref{eq:prooflemmas20nonspu}. The intuition is that both $T$ and $M$ are nested within $Q_T$, and therefore $L_{Q_T T}$ and $L_{Q_T M}$ follow chi-squared distributions.
We first bound the first term. %in the right-hand side of \eqref{eq:prooflemmas20nonspu}. 
If $M \subset T$ then $Q_T=T$, $L_{Q_TT}=0$, and this term is zero. Suppose now that $M \not\subset T$. %$T$ is a strict subset of $Q_T$. 
Then, by Lemma~\ref{lemma:s7}, %$L_{Q_TT}$ has a $\chi^2$-distribution. More precisely, 
$L_{Q_TT} \sim \mathcal{\chi}_{|Q_T\setminus T|}^2(\mu_{Q_TT})$ with $|Q_T\setminus T|=|M\setminus T|$. By assumption, $A_T>0$, $|M\setminus T|=o(\ln(e^{A_T}))$ and $\mu_{Q_TT}=o(\ln(e^{A_T}))$, then  by Lemma~\ref{lemma:s20}, for $\alpha\in(\psi,1)$, and every $n$ large enough,
\begin{equation}\label{eq:boundLQS}
\int_0^1P\Bigl(L_{Q_TT}>2\ln{\Bigl(\frac{e^{A_T}}{1 / u-1}\Bigl)}\Bigl){\rm d}u \;<\; 
e^{-\alpha A_T}.
\end{equation}

We now bound the second term in \eqref{eq:prooflemmas20nonspu}. If $M \supset T$, then $Q_T=M$, $L_{Q_TM}=0$, $\mu_{Q_TM}=0$, and this term is zero. Alternatively, if $M \not\supset T$ %$M$ is a strict subset of $Q_T$ 
then, by Lemma~\ref{lemma:s7}, $L_{Q_TM} \sim \chi_{|Q_T\setminus M|}^2(\mu_{Q_TM})$ with $|Q_T\setminus M|=|T\setminus M|$.
Clearly, when $\mu_{Q_TM}\leq 6\Delta_{MT}$ this term is also zero, so suppose that $\mu_{Q_TM}> 6\Delta_{MT}$.  
We have  
$$P\left(L_{Q_TM}<2\gamma'(\mu_{Q_TM}-6\Delta_{MT})\right)\; \leq \; P\left(L_{Q_TM}<2\gamma'\mu_{Q_TM}\right)$$ 
and, by Lemma~\ref{lemma:s2} and using that $\gamma'\in (0,\tfrac{1}{12})$, we obtain that 
$$P\left(L_{Q_TM}<2\gamma'\mu_{Q_TM}\right) \;\;\leq\;\; (2\gamma')^{\frac{|T\setminus M|}{4}} e^{-\frac{1}{2}(1-\sqrt{2\gamma'})^2\mu_{Q_TM}}\;\;\leq\;\;  (\tfrac16)^{\frac{|T\setminus M|}{4}}e^{-\frac{1}{2}(1-\sqrt{2\gamma'})^2\mu_{Q_TM}}.$$
Since $\mu_{Q_TM}>6\Delta_{MT}$, we get
$$
A_T\;=\;\gamma'\mu_{Q_TM}+\gamma\Delta_{MT}\;<\;\frac{1-\gamma}{6}\mu_{Q_TM}+\frac{\gamma}{6}\mu_{Q_TM}\;=\;\frac16\mu_{Q_TM}\;\leq\;\frac12(1-\sqrt{2\gamma'})^2\mu_{Q_TM},
$$
where the last inequality follows from the fact that $\gamma'\in(0, \tfrac{1}{12})$. We then get
\begin{equation}\label{eq:boundLQM}
    P\left(L_{Q_TM}<2\gamma'(\mu_{Q_TM}-8\Delta_{MT})\right) \;\leq\; \left(\frac14\right)^{\frac{|T\setminus M|}{6}}e^{-A_T} \leq e^{-A_T}\leq e^{-\alpha A_T}
\end{equation}
Summing the bounds in \eqref{eq:boundLQS} and \eqref{eq:boundLQM} gives that for every $n$ large enough $\E\left(NC(M)\right) < 2e^{-\alpha A_T}$. Since $\psi<\alpha$, for every $n$ large enough, we have that $\E\left(NC(M)\right) < e^{-\psi A_T}$ as we wished to prove.
\end{proof}

\begin{lem}\label{lem:upperboundnoncentralgen}
    For any $T\subseteq S$ and $M\in \M$, let $Q_T=M\cup T$, and $\mu_{Q_T T}$ as defined in \eqref{eq:muQM}, then, 
    \begin{equation}\label{eq:upperbounnnnnd}
      \mu_{Q_TT} \leq n\,\bar{\lambda}\,\sum_{j=1}^b|(S_j\cap M_j )\setminus T_j|\,\max_{i\in (S_j\cap M_j )\setminus T_j}{\beta_{i}^*}^2 \quad \text{where}\quad \bar{\lambda}\,:=\,\lambda_{\max}\Big(\frac1n \bs X_{S}^\top\bs X_{S}\Big)  
    \end{equation}
    and
    \begin{equation}\label{eq:lowerbounnnnnd}
    \mu_{Q_TT} \geq n\,\lambda_{\min}\Big(\frac1n \bs X_{S}^\top\bs X_{S}\Big)\,\sum_{j=1}^b|(S_j\cap M_j )\setminus T_j|\,\min_{i\in (S_j\cap M_j )\setminus T_j}{\beta_{i}^*}^2
    \end{equation}
\end{lem}

\begin{proof}
Using the definition of $\mu_{Q_T T}$ in \eqref{eq:muQM}, we have that 
\begin{eqnarray}
    \mu_{Q_TT}&=&{\bs \beta^*_{Q_T\setminus T}}^{\top}\bs X_{Q_T\setminus T}^{\top}  \big(I_n-P_T\big) \bs X_{Q_T\setminus T} \bs \beta^*_{Q_T\setminus T}\nonumber\\
    &=&{\bs \beta^*_{M\setminus T}}^{\top}\bs X_{M\setminus T}^{\top}  \big(I_n-P_T\big) \bs X_{M\setminus T} \bs \beta^*_{M\setminus T}\nonumber\\
    &=&{\bs \beta^*_{(S\cap M)\setminus T}}^{\top}\bs X_{(S\cap M)\setminus T}^{\top}  \big(I_n-P_T\big) \bs X_{(S\cap M)\setminus T} \bs \beta^*_{(S\cap M)\setminus T}\label{eq:poihe}
\end{eqnarray}
where the second equality follows from $Q_T\setminus T =M\setminus T $ and the third equality from $\bs \beta^*_{M\setminus S}=0$.
We start by showing the upper bound in \eqref{eq:upperbounnnnnd}. Denote for any square matrix $A$, its largest eigenvalue 
$\lambda_{\max}(A)$. By \eqref{eq:poihe}, we have that
$$\mu_{Q_TT}\leq n\lambda_{\max}\Big(\frac{1}{n}\bs X_{(S\cap M)\setminus T}^{\top}  \big(I_n-P_T\big) \bs X_{(S\cap M)\setminus T}\Big) \|\bs \beta^*_{(S\cap M)\setminus T}\|^2.$$
Let $B:=\tfrac{1}{n}\bs X_{(S\cap M)\setminus T}^{\top}  \big(I_n-P_T\big) \bs X_{(S\cap M)\setminus T}$, $C:=\tfrac{1}{n}\bs X_{(S\cap M)\cup T}^{\top}\bs X_{(S\cap M)\cup T}$ and $D:=\tfrac{1}{n}\bs X_T^{\top} \bs X_T$. $D$ is a principal submatrix of $C$ and $B$ is the Schur complement of $D$ of $C$. The inverse $B^{-1}$ is then a principal submatrix of $C^{-1}$, and by Cauchy's interlacing theorem we have that   $\lambda_{\min}(B^{-1})\geq \lambda_{\min}(C^{-1})$ and then $\lambda_{\max}(B)\leq \lambda_{\max}(C)$. Since $T \subseteq S$ by assumption, we also have that $(S\cap M)\cup T \subseteq S$, then by interlacing again $\lambda_{\max}(C)\leq \lambda_{\max}\Big(\tfrac{1}{n}\bs X_S^{\top}\bs X_S\Big)=\bar{\lambda}$. The upper bound in \eqref{eq:upperbounnnnnd} follows from the latter inequality and also observing that $\|\bs \beta^*_{(S\cap M)\setminus T}\|_2^2\leq \sum_{j=1}^b |(S_j\cap M_j)\setminus T_j|\max_{i\in (S_j\cap M_j)\setminus T_j}{\beta_i^*}^2$.

We now derive the lower bound in \eqref{eq:lowerbounnnnnd}. By \eqref{eq:poihe}, we have that %a fist lower bound on $\mu_{Q_TT}$ is: 
$$\mu_{Q_TT}\geq n\lambda_{\min}\Big(\frac{1}{n}\bs X_{(S\cap M)\setminus T}^{\top}  \big(I_n-P_T\big) \bs X_{(S\cap M)\setminus T}\Big) \|\bs \beta^*_{(S\cap M)\setminus T}\|_2^2.$$

Recall that $B^{-1}$ is a principal submatrix of $C^{-1}$, hence by interlacing $\lambda_{\max}(B^{-1})\leq \lambda_{\max}(C^{-1})$ and $\lambda_{\min}(B)\geq \lambda_{\min}(C)$. Since $T \subseteq S$ by assumption, we have that $(S\cap M)\cup T \subseteq S$, and hence $\lambda_{\min}(C)\geq\lambda_{\min}(\frac{1}{n}\bs X_S^{\top} \bs X_S)$.
The bound in \eqref{eq:lowerbounnnnnd} is obtained by using the latter inequality and noting that also $\|\bs \beta^*_{(S\cap M)\setminus T}\|_2^2\geq \sum_{j=1}^b |(S_j\cap M_j)\setminus T_j|\min_{i\in (S_j\cap M_j)\setminus T_j}{\beta_i^*}^2.$
\end{proof}

\section{Proofs of Section~\ref{sec:properties_orth}}\label{suppsec:proofsecgsm}
 
 

\subsection{Proof of Proposition~\ref{prop:isthresh}}
First note that by discarding constant term 
 $$\arg\max_{M\in \M} \left\{\max_{\beta\in \mathcal{L}_M}\ell(\bs y; \bs \beta)-\sum_{j=1}^b\kappa_j |M_j| \right\}=\arg\min_{M\in \M} \left\{\min_{\beta\in \mathcal{L}_M}\bigg\{ \tfrac{1}{2}\|\bs y-\sqrt{n}\bs\beta\|^2\right\}+\sum_{j=1}^b\kappa_j |M_j| \bigg\}$$

 We also have $\min_{\beta\in \mathcal{L}_M} \tfrac{1}{2}\|\bs y-\sqrt{n}\bs\beta\|^2=\tfrac{1}{2}\|\bs y-\sqrt{n}\tilde{\bs \beta}_M\|^2=\tfrac{1}{2}\|\bs y\|^2-\tfrac{1}{2}\|\sqrt{n}\tilde{\bs \beta}_M\|^2$. The maximization in \eqref{eq:Mhat} can be then replaced with:
     \begin{equation}\label{eq:MhatGSM}
  \arg\max_{M\in \M} \bigg\{\tfrac{1}{2}\| \sqrt{n}\tilde{\bs \beta}_M\|^2-\sum_{j=1}^b\kappa_j |M_j|\bigg\}, \qquad M=M_1\cup \ldots\cup M_b.	
\end{equation}
Under the assumptions of \eqref{eq:GSM}, we can write 
	$$
	\tfrac{1}{2}\|\sqrt{n}\tilde{\bs \beta}_M\|^2-\sum_{j=1}^b\kappa_j |M_j|\;=\;\sum_{j=1}^b\sum_{i\in M_j}(\tfrac{n}{2}{\tilde{ \beta}_i}^2-\kappa_j).
	$$
	Then \eqref{eq:MhatGSM} can be maximized with respect to each $M_j$ by including $i\in \hat S_j$ whenever $n\tilde\beta_i^2\geq 2\kappa_j$.

\subsection{Proof of Proposition~\ref{prop:suffnecblockthresh}} 

\subsubsection{Part (i)} 
By the union bound and by Lemma~\ref{lemma:3tailbounds} (i),
\begin{equation}\label{eq:type1convrate}
P\left(\hat{S}^{b} \not\subseteq S\right)\;\leq\; \sum_{j=1}^b P\Big(\max_{i\in B_j\setminus S_j}|y_i/\sqrt{n}|>\tau_{j}\Big)\;\leq\;\sum_{j=1}^b\frac{e^{-\tfrac{n}{2}\left(\tau_j^2-\tfrac{2\ln(p_j-s_j)}{n}\right)}}{\sqrt{\pi\ln(p_j-s_j)}} 	.
 \end{equation}
 By Assumption~\hyperlink{lab:A4}{A4}, the numerator on the right is bounded above by $1$ for all sufficiently large $n$. It follows that 
	$$P\left(\hat{S}^{b} \not\subseteq S\right) \;\leq\; \sum_{j=1}^b (\pi\ln (p_j-s_j))^{-1/2} \to 0\qquad \mbox{as } n\to \infty.$$
\subsubsection{Part (ii)}  By independence, we have 
$$P(\hat{S}^{b} \subseteq S)\;=\;\prod_{j=1}^b P(\max_{i \in B_j\setminus S_j}|y_i/\sqrt{n}| \leq \tau_j).$$ 
Consider $j$ such that $\lim_{n\to \infty} \frac{\sqrt{n}\tau_j}{\sqrt{2\ln(p_j-s_j)}} \;<\;1$. In particular, there exists $c<1$ such that, for all sufficiently large $n$, we have that $\tau_j\leq c\sqrt{2\ln(p_j-s_j)/n}$. Then, for any such $n$,
\begin{align*}
    P\Big(\max_{i \in B_j\setminus S_j}|y_i/\sqrt{n}| \leq \tau_j\Big)&\;\leq\; P\Big(\max_{i \in B_j\setminus S_j}|y_i| \leq c \sqrt{2\ln(p_j-s_j)}\Big)\\&\;\leq\;P\Big(\tfrac{1}{\sqrt{2\ln(p_j-s_j)}}\max_{i \in B_j\setminus S_j}y_i \leq c \Big).
\end{align*}
On the other hand, results from extreme value theory, \cite{galambos} Example 4.4.1, show that 
 $$\frac{1}{\sqrt{2 \ln(p_j-s_j)}}\max_{i\in B_j\setminus S_j}y_i \;\overset{P}{\longrightarrow} 1$$
 and so 
 $$
 P\Big(\frac{1}{\sqrt{2\ln(p_j-s_j)}}\max_{i \in B_j\setminus S_j}y_i \leq c \Big)\;\to\;0,
 $$ 
 which implies that $P(\hat{S}^{b} \subseteq S)\to 0$. 

\subsubsection{Part (iii)}
	 By the union bound, 
	$$P\big(\hat{S}^{b} \not\supseteq S\big)\leq \sum_{j=1}^b  P\left(\min_{i\in S_j}|y_i/\sqrt{n}|\leq\tau_j\right).$$ 
	By Lemma~\ref{lemma:3tailbounds} (ii), for each $j$,	
	$$P\left(\min_{i\in S_j}|y_i/\sqrt{n}|\leq\tau_j\right)\;\leq\; P\left(\max_{i\in S_j}|y_i/\sqrt{n}-\beta_i| \geq\beta_{\min,j}^*-\tau_j \right).$$ By Lemma~\ref{lemma:3tailbounds} (i)  
	\begin{equation}\label{eq:type2convrate}
P\left(\hat{S}^{b}\not\supseteq S \right) \;\leq\;\sum_{j=1}^b\frac{e^{-\tfrac{n}{2}\left((\beta^*_{\min,j}-\tau_j)^2-\tfrac{2\ln(s_j)}{n}\right)}}{\sqrt{\pi\ln(s_j)}}.		
	\end{equation}
 By Assumption~\hyperlink{lab:A5}{A5}, the nominator on the right is bounded above by $1$ for all sufficiently large $n$. It follows that, byAssumption  \hyperlink{lab:A3}{A3} 
	$$P\left(\hat{S}^{b} \not\supseteq S \right) \;\leq\;  \sum_{j=1}^b (\pi\ln (s_j))^{-1/2} \to 0\qquad \mbox{as } n\to \infty.$$ 

\subsubsection{Part (iv)}	
Take any $j=1,\ldots,b$ satisfying $\lim_{n\to \infty} \sqrt{n}(\beta_{\min,j}^*-\tau_j)/\sqrt{(\pi/2)\ln(s_j)}\leq 1$. Consider first the case $\beta^*_{\min,j}\leq\tau_j$. Then $ \lim_{n\to \infty} \sqrt{n}(\beta_{\min,j}^*-\tau_j)\; <\; \infty$, and by Lemma~\ref{lem:gennecIIblockthresh}, $\lim_{n \rightarrow \infty} P(\hat S^{b}\supseteq S  ) < 1$. We now consider the case $\beta^*_{\min,j}>\tau_j$. By Lemma~\ref{lemma:3tailbounds} (iii):
$$P\left(\min_{i \in S_j}\Big|\frac{y_i}{\sqrt{n}}\Big| > \tau_j \right)\;\leq\;
\exp\left\{-\frac{s_j}{2}\frac{e^{-\tfrac{2n}{\pi}(\beta^*_{\min,j}-\tau_j)^2}-e^{-\tfrac{n}{2}(\tau_j+\beta^*_{\min,j})^2}}{\left(1-e^{-2n(\beta^*_{\min,j}-\tau_j)^2/\pi}\right)^{\frac{1}{2}}+\left(1-e^{-n(\tau_j+\beta^*_{\min,j})^2/2}\right)^{\frac{1}{2}}}\right\}. $$
Let $a_{n}=2\left(1-e^{-2n(\beta^*_{\min,j}-\tau_j)^2/\pi}\right)^{\frac{1}{2}}+2\left(1-e^{-n(\tau_j+\beta^*_{\min,j})^2/2}\right)^{\frac{1}{2}}$ and note that $a_n\in (0, 4]$.  Thus, to show that $P\left(\min_{i \in S_j}|\tfrac{y_i}{\sqrt{n}}| > \tau_j \right)$ is bounded away from $1$, it is enough to show that 
\begin{align*}
    %\lim_{n\to \infty} \frac{s_j}{a_n}({e^{-2n(\beta^*_{\min,j}-\tau_j)^2/\pi}-e^{-n(\tau_j+\beta^*_{\min,j})^2/2}})&\;\geq\; 
\lim_{n\to \infty} \frac{s_j}{4}\big({e^{-2n(\beta^*_{\min,j}-\tau_j)^2/\pi}-e^{-n(\tau_j+\beta^*_{\min,j})^2/2}}\big)
&\;>\;0.
\end{align*}

Since $\lim_{n\to \infty} \sqrt{n}(\beta_{\min,j}^*-\tau_j)/\sqrt{(\pi/2)\ln(s_j)}\leq 1$, we have that $\lim_{n\to\infty} {s_j} e^{-2n(\beta^*_{\min,j}-\tau_j)^2/\pi}\geq 1$. To conclude, we show that
$$
\lim_{n\to\infty} {s_j} e^{-n(\tau_j+\beta^*_{\min,j})^2/2}\;\leq\;c\;<\;1
$$ 
for some $c\in(0,1)$. Take $c$ such that $\tfrac{s_j}{p_j}\leq c$. Such $c$ exists by our assumption $\tfrac{s_j}{p_j}<1$.
We equivalently need
$\sqrt{n}(\tau_j+\beta^*_{\min,j})\geq\sqrt{{2\ln (s_j/c)}}$ for all $n$ sufficiently large. To show that, note that, by assumption $\tau_j< \beta^*_{\min,j}$, and $s_j/c\leq p_j$, which gives
$$
\frac{\tau_j+\beta^*_{\min,j}}{\sqrt{\tfrac{2\ln(s_j/c)}{n}}}\;\geq\;\frac{2\tau_j}{\sqrt{\tfrac{2\ln(p_j)}{n}}}\;=\;\frac{\sqrt{\tfrac{2\ln(p_j-s_j)}{n}}}{\sqrt{\tfrac{2\ln(p_j)}{n}}}\frac{2\tau_j}{\sqrt{\tfrac{2\ln(p_j-s_j)}{n}}}\;=\;\sqrt{\frac{\ln(p_j-s_j)}{\ln(p_j)}}\frac{2\tau_j}{\sqrt{\tfrac{2\ln(p_j-s_j)}{n}}}.
$$
Since by assumption $\tau_j$ satisfies $\lim_{n\to \infty} \sqrt{n}\tau_j/\sqrt{2\ln(p_j-s_j)}\;\geq\;1$, the second term on the right converges to something $\geq 2$, and it is enough to show that the first term converges to something $>1/2$. Using that $\tfrac{s_j}{p_j}<c$, we get
$$
\sqrt{\frac{\ln(p_j-s_j)}{\ln(p_j)}}
= \sqrt{\frac{\ln(p_j) + \ln(1-\frac{s_j}{p_j})}{\ln(p_j)}}
\;\geq\;\sqrt{\frac{\ln(p_j)+\ln(1-c)}{\ln(p_j)}}\;\to\;1,
$$
which concludes the proof.

\subsection{Proof of Lemma~\ref{lem:necbetaminblockthresh}}
    We have $$\frac{\beta_{\min,j}^*}{\sqrt{\tfrac{2\ln(p_j-s_j)}{n}}+\sqrt{\tfrac{\pi}{2}\tfrac{\ln(s_j)}{n}}}= \frac{\beta_{\min,j}^*-\tau_j+\tau_j}{\sqrt{\tfrac{2\ln(p_j-s_j)}{n}}+\sqrt{\tfrac{\pi}{2}\tfrac{\ln(s_j)}{n}}}$$ and \eqref{eq:necbetaminblockthresh} implies that we have either $\lim_{n\to \infty}\sqrt{n}\tau_j/\sqrt{2\ln(p_j-s_j)} <1$ or else $\lim_{n\to \infty}\sqrt{n}(\beta_{\min,j}^*-\tau_j)/\sqrt{(\pi/2)\ln(s_j)} <1$. 
    
    If $\lim_{n\to \infty}\sqrt{n}\tau_j/\sqrt{2\ln(p_j-s_j)} <1$ by Proposition~\ref{prop:suffnecblockthresh} (ii), $\lim_{n \to \infty} P(\hat{S}^{b} = S) < 1$ and recovery is not possible asymptotically. 
    % We next examine the consequences of $\lim_{n\to \infty}\sqrt{n}(\beta_{\min,j}^*-\tau_j)/\sqrt{(\pi/2)\ln(s_j)} <1$. To do so first, observe that if $\beta_{\min,j}^*<\tau_j$ then $ \lim_{n\to \infty} \sqrt{n}(\beta_{\min,j}^*-\tau_j)\; <\; \infty$, and by Lemma~\ref{lem:gennecIIblockthresh}, $\lim_{n \rightarrow \infty} P(\hat S^{b}\supseteq S  ) < 1$ and recovery is not possible asymptotically. 
    %Indeed, denote by $i_{\circ}\in B_j$ an entry such that $|\beta_{i_{\circ}}^*|=\beta^*_{\min,\,j}$. In the proof of Lemma~\ref{lemma:3tailbounds} (iii), we show that if $y\sim N(\mu, \sigma^2)$, then for any $a<|\mu|$ , $P(|y|>a)  = P(z>\frac{a-|\mu|}{\sigma}) + P(z<-\frac{a+|\mu|}{\sigma})$ where $z\sim N(0,1)$. We have then
    % $$P(|y_{i_{\circ}}/\sqrt{n}|>\tau_j) 
    % = P\left(z>\sqrt{n}\left(\tau_j-\beta_{\min,\,j}^*\right)\right) + P\left(z<-\sqrt{n}\left(\tau_j+\beta_{\min,\,j}^*\right)\right)$$ where $z\sim N(0,1)$. Since $\tau_j>\beta^*_{\min,j}$, $\sqrt{n}\left(\tau_j-\beta_{\min,\,j}^*\right)\to\infty$ and we have that $P(z>\sqrt{n}\big(\tau_j+\beta_{\min,\,j}^*\big)) \to 0$. Next, since $\sqrt{n}\left(\tau_j+\beta_{\min,\,j}^*\right)\geq0$, we have that $P\left(z<-\sqrt{n}\left(\tau_j+\beta_{\min,\,j}^*\right)\right)<1$ and $P(|y_{i_{\circ}}/\sqrt{n}|>\tau_j)<1$. Hence, $\underset{n \to \infty}{\lim}P\left(\hat{S}^{b} \supseteq S \right)<1$.
    Suppose now that $\lim_{n\to \infty}\sqrt{n}\tau_j/\sqrt{2\ln(p_j-s_j)} \geq 1$, then it holds that $\lim_{n\to \infty}\sqrt{n}(\beta_{\min,j}^*-\tau_j)\sqrt{\pi\ln(s_j)/2} <1$. By Proposition~\ref{prop:suffnecblockthresh} (iv), $\lim_{n \to \infty} P(\hat{S}^{b} \supseteq S ) < 1$ and recovery is not possible asymptotically.

\subsection{Proof of Theorem~\ref{theo:fullrecovGS}} We start by showing the bound in \eqref{eq:rateconvgsm}. By the union bound,
\begin{equation}\label{eq:grpihj}
    P(S\neq \hat{S}^b) \; \leq\; P(\hat{S}^b\not\supseteq S)+P(\hat{S}^b\not\subseteq S)
\end{equation}
By the union bound and by Lemma~\ref{lemma:3tailbounds} (i), for $n$ large enough,
\begin{equation}\label{eq:type1convrate2}
P\left(\hat{S}^{b} \not\subseteq S\right)\;\leq\; \sum_{j=1}^b P\Big(\max_{i\in B_j\setminus S_j}|y_i/\sqrt{n}|>\tau_{j}\Big)\;\leq\;\sum_{j=1}^b\frac{e^{-\tfrac{n}{2}\left(\tau_j^2-\tfrac{2\ln(p_j-s_j)}{n}\right)}}{\sqrt{\pi\ln(p_j-s_j)}} 	.
 \end{equation}
where the assumption of Lemma~\ref{lemma:3tailbounds} (i) is met because Assumption~\hyperlink{lab:A4}{A4} is assumed to hold. Moreover, by the union bound, 
	$P\big(\hat{S}^{b} \not\supseteq S\big)\leq \sum_{j=1}^b  P\left(\min_{i\in S_j}|y_i/\sqrt{n}|\leq\tau_j\right)$. 
	By Assumption~\hyperlink{lab:A5}{A5}, $\beta_{\min,j}>\tau_j$ and by Lemma~\ref{lemma:3tailbounds} (ii), for each $j$,	
	$$P\left(\min_{i\in S_j}|y_i/\sqrt{n}|\leq\tau_j\right)\;\leq\; P\left(\max_{i\in S_j}|y_i/\sqrt{n}-\beta_i| \geq\beta_{\min,j}^*-\tau_j \right).$$ By Assumption~\hyperlink{lab:A5}{A5}, $\beta_{\min,j}-\tau_j\geq \sqrt{2\ln(s_j)/n}$ and by Lemma~\ref{lemma:3tailbounds} (i)
	\begin{equation}\label{eq:type2convrate2}
P\left(\hat{S}^{b}\not\supseteq S \right) \;\leq\;\sum_{j=1}^b\frac{e^{-\tfrac{n}{2}\left((\beta^*_{\min,j}-\tau_j)^2-\tfrac{2\ln(s_j)}{n}\right)}}{\sqrt{\pi\ln(s_j)}}.
\end{equation}
Inputting the bounds in \eqref{eq:type1convrate} and \eqref{eq:type2convrate} into \eqref{eq:grpihj} gives
\begin{equation}\label{eq:fhfhfhf}
P(S\neq \hat{S}^b) \; \leq\;\sum_{j=1}^b\frac{e^{-\tfrac{n}{2}\left(\tau_j^2-\tfrac{2\ln(p_j-s_j)}{n}\right)}}{\sqrt{\pi\ln(p_j-s_j)}}+\sum_{j=1}^b\frac{e^{-\tfrac{n}{2}\left((\beta^*_{\min,j}-\tau_j)^2-\tfrac{2\ln(s_j)}{n}\right)}}{\sqrt{\pi\ln(s_j)}}.
\end{equation}
which shows \eqref{eq:rateconvgsm}.

We continue with the second part of the theorem and show that if \eqref{eq:betamin} holds then $\tau^*_j$ satisfies Assumptions~\hyperlink{lab:A4}{A4} and ~\hyperlink{lab:A5}{A5}. Under \eqref{eq:betamin}, there exists $c>1$ such that $\beta^*_{\min,j}=c(\sqrt{\tfrac{2\ln(p_j-s_j)}{n}} + \sqrt{\tfrac{2\ln(s_j)}{n}})$. Denote $a=\sqrt{\tfrac{2\ln(p_j-s_j)}{n}}$ and $b=\sqrt{\tfrac{2\ln(s_j)}{n}}$. We have
$$
\tau_j^*\;=\;\frac{c}{2} (a+b)+\frac{(a^2-b^2)}{2c(a+b)}\;=\;\frac{c}{2} (a+b)+\frac{a-b}{2c}\;=\;\frac{c^2+1}{2c}a+\frac{c^2-1}{2c}b.
$$
Since $\tfrac{c^2+1}{2c}> 1$, $\tfrac{c^2-1}{2c}> 0$ and $b \geq 0$, $\tau^*_j>a=\sqrt{\tfrac{2\ln(p_j-s_j)}{n}}$ and so Assumption~\hyperlink{lab:A4}{A4} holds for the $\tau^*_j$. To show that the upper bound in Assumption~\hyperlink{lab:A5}{A5} also holds, observe that
$$
\beta_{\min,j}^* - \tau^*_j \;=\; c(a+b)- \frac{c^2+1}{2c}a-\frac{c^2-1}{2c}b \;= \; \frac{c^2-1}{2c}a+\frac{c^2+1}{2c}b\;>\;  b=\sqrt{\tfrac{2\ln(s_j)}{n}}.
$$
We proceed with the proof of the bound in \eqref{eq:optimrateconvgsm}. Since Assumptions~\hyperlink{lab:A4}{A4} and \hyperlink{lab:A5}{A5} hold for the $\tau^*_j$'s, for all sufficiently large $n$, \eqref{eq:fhfhfhf} holds for these oracle penalties. Moreover, by assumption $p_j-s_j>1$ and $s_j>1$, then $\sqrt{\pi\ln(p_j-s_j)}>1$, and $\sqrt{\pi\ln(s_j)}>1$ and we get the bound:
$$P(S\neq \hat{S}^b) \; \leq\;\sum_{j=1}^be^{-\tfrac{n}{2}\left({\tau^*_j}^2-\tfrac{2\ln(p_j-s_j)}{n}\right)}+\sum_{j=1}^be^{-\tfrac{n}{2}\left((\beta^*_{\min,j}-\tau^*_j)^2-\tfrac{2\ln(s_j)}{n}\right)}.$$
Simple algebra shows that the $\tau_j^*$ satisfies ${\tau_j^*}^2-\tfrac{2\ln(p_j-s_j)}{n}=(\beta^*_{\min,j}-{\tau_j^*})^2-\tfrac{2\ln(s_j)}{n}$ for all $j$. It follows that
\begin{equation}\label{eq:edpo}
P(S\neq \hat{S}^b) \; \leq\;2\sum_{j=1}^be^{-\tfrac{n}{2}\left({\tau^*_j}^2-\tfrac{2\ln(p_j-s_j)}{n}\right)}.
\end{equation}
For convenience, denote $d = \frac{1}{n \beta_{\min,j}^*} \ln (p_j / s_j-1)$ such that $\tau^*_j= \beta_{\min,j}^*/2 + d$. Then
$$e^{-\frac{n}{2}\left({\tau^*_j}^2-\frac{2\ln(p_j-s_j)}{n}\right)}
=e^{-\left[\frac{n}{8}{\beta_{\min,j}^*}^2+\frac{n}{2}d^2 - \frac{1}{2}\left(\ln (p_j- s_j)+\ln (s_j)\right) \right]}.$$
By considering separately the two possible maxima in $\ln \max\{p_j-s_j,s_j\}$, we get that
$$\frac{e^{-\frac{n}{2}\left({\tau^*_j}^2-\frac{2\ln(p_j-s_j)}{n}\right)}}{e^{-[\frac{n}{8}{\beta_{\min,j}^*}^2-\ln \max\{p_j-s_j,s_j\}]}}= 
e^{-\left[\frac{n}{2}d^2 + \frac{1}{2}\left|\ln (p_j- s_j)-\ln (s_j)\right|\right]} <1.$$
From \eqref{eq:edpo}, we then have
$$P(\hat{S}^{b} \neq S) \; \leq \; 2 \sum_{j=1}^be^{-[\frac{n}{8}{\beta^*_{\min,j}}^2-\ln \max\{p_j-s_j,s_j\}]}.$$
which proves \eqref{eq:optimrateconvgsm}.

\subsection{Proof of Corollary~\ref{cor:selconstimpossible}}
Since $\hat{S}$ is $\hat{S}^b$ with $b=1$, the assumptions of Lemma~\ref{lem:necbetaminblockthresh} for $\hat{S}$ are met and $\lim_{n \to \infty} P(\hat{S} = S) < 1$. Since Assumptions~\hyperlink{lab:A4}{A4} and \hyperlink{lab:A5}{A5} hold, by Proposition~\ref{prop:suffnecblockthresh}, $\lim_{n \to \infty} P(\hat{S}^b \subseteq S)= 1$ and $\lim_{n \to \infty} P(\hat{S}^b \supseteq S)= 1$, and then $\lim_{n \to \infty} P(\hat{S}^b = S)= 1$.\color{black}

\section{Proofs of Section~\ref{sec:properties_legression}}\label{suppsec:proofseclinreg}
\subsection{Proof of Lemma~\ref{lem:useful-reformulation}}
	First note that by discarding constant terms 
 $$\arg\max_{M\in \M} \left\{\max_{\beta\in \mathcal{L}_M}\ell(\bs y; \bs \beta)-\sum_{j=1}^b\kappa_j |M_j| \right\}=\arg\min_{M\in \M} \left\{\min_{\beta\in \mathcal{L}_M}\bigg\{ \tfrac{1}{2}\|\bs y-\bs X\bs\beta\|^2\right\}+\sum_{j=1}^b\kappa_j |M_j| \bigg\}.$$
 
 We also have that 
 $$\min_{\beta\in \mathcal{L}_M} \tfrac{1}{2}\|\bs y-\bs X\bs\beta\|^2=\tfrac{1}{2}\|\bs y-\bs X_{\!\!M}\tilde{\bs \beta}^{(M)}\|^2=\tfrac{1}{2}\|\bs y\|^2-\tfrac{1}{2}\|\bs X_{\!\!M}\tilde{\bs\beta}^{(M)}\|^2,$$
where in the last equality we used that $\tilde{\bs \beta}^{(M)}= (\bs X_M^T \bs X_M)^{-1} \bs X_M^T \bs y$.
 The maximization in \eqref{eq:Mhat} can be then replaced with the maximization of $C(M)=\tfrac{1}{2}
 \|\bs X_{\!\!M}\tilde{\bs\beta}^{(M)}\|^2-\sum_{j=1}^b \kappa_j|M_j|$ which is equivalent to maximizing $NC(M)$.

\subsection{Proof of Lemma~\ref{lem:L0toL1convergence}}
Part (i) follows directly from Lemma~\ref{lem:L0toL1convergencegen} by taking $\{M_1,\ldots,M_k\}=\{S\}$.
Part (ii) follows from
$$NC(M)=\Big(1+\sum_{N \neq M} e^{C(N)-C(M)}\Big)^{-1}
< \big(1+ e^{C(M')-C(M)}\big)^{-1}.$$ 
%Since for any $M,N\in \M$ and $M\neq N$, $C(N) / C(M)>0$, $NC(M) \leq (1+e^{C(M')-C(M)})^{-1}$.

\subsection{Proof of Lemma~\ref{lemma:s7}}
    This result follows directly from Lemma~S7 in \cite{rossell2022concentration} taking $\phi^*=1$.

\subsection{Proof of Lemma~\ref{lem:boundrho}}
This result follows directly Lemma~\ref{lem:boundrhogen} by taking $T=S$, and observing that $\min_{i\in S_j\setminus M_j}{\beta^*_i}^2\geq {\beta^*_{\min,j}}^2$.

\subsection{Proof of Theorem~\ref{theo:suffcondlinearmodel}} 
The proof strategy is to first use Lemma~\ref{lem:boundexpnc} with $T=S$ to show that for every $M\neq S$, $\E(NC(M)) \leq e^{-\psi A_S}$ for every large enough $n$ and any $\psi \in (0,1)$, where $A_S=\gamma\Delta_{MS}+\tfrac{1-\gamma}{6}\mu_{Q_SM}$ (\emph{cf} \eqref{eq:DeltaM} and \eqref{eq:muQM}), $\gamma\in(1/2,1)$ is defined in Assumption~\hyperlink{lab:A7}{A7} and $Q_S= M\cup S$. Assumption~\hyperlink{lab:A6}{A6} and the fact that $M\setminus S\subseteq S^C$ ensure the 
assumptions of Lemma~\ref{lem:boundexpnc} are met. 
The second step is to obtain a lower bound for $A_S$, which gives a new upper bound for $\E(NC(M))$.
The final step is to use these bounds to get an upper-bound on $\sum_{M \in \M\setminus \{S\}} \E\left(NC(M)\right)$ that asymptotically vanishes under Assumptions~\hyperlink{lab:A6}{A6} and \hyperlink{lab:A7}{A7}. We then use Lemma~ \ref{lem:L0toL1convergence} to conclude on the vanishing of $P(\hat S^b \neq S)$. 

First, to show that $\E(NC(M)) \leq e^{-\psi A_S}$
for any $M\in\mathcal{M}\setminus \{S\}$, we show that $A_S$ satisfies the conditions of Lemma~\ref{lem:boundexpnc}, taking $T=S$. 
That is, we wish to show that, $A_S>0$, $|M\setminus S|=o(A_S)$, and $\mu_{Q_S S}=o(A_S)$. 
Observe that $\Delta_{MS}$, defined in \eqref{eq:DeltaM}, can be rewritten as $\Delta_{MS}=\sum_{j=1}^b(|M_j\setminus S_j|-|S_j\setminus M_j|)\kappa_j$. By  Lemma~\ref{lem:boundrho}, for every $n\in \N$ we have 
\begin{equation}\label{eq:nonoverfittingcond}
\begin{aligned}
     A_S & \;=\; \gamma \Delta_{MS}+\frac{1-\gamma}{6}\mu_{Q_SM}\\
     & \;\geq\;  \gamma\sum_{j=1}^b|M_j \setminus S_j|\kappa_j + \sum_{j=1}^b |S_j\setminus M_j |\left(\tfrac{1-\gamma}{6} n \rho(\bs X) {\beta^*_{\min,j}}^{2} - \gamma\kappa_j\right) 
\end{aligned}
\end{equation} 
Since $M\neq S$, $|M\setminus S|\neq0$ or $|S\setminus M|\neq0$, then by Assumptions~\hyperlink{lab:A6}{A6} and \hyperlink{lab:A7}{A7}, for every $n$ large enough, $A_S>0$. We immediately have $\mu_{Q_S S}=o(A_S)$ because $\beta^*_{Q_S\setminus S}=\beta^*_{M\setminus S}=0$ (any parameter outside the true support $S$ is by definition 0) and hence $\mu_{Q_S S}=0$.
If $|M\setminus S|=0$, $|M\setminus S|=o(A_S)$ also immediately. Consider now the case $|M\setminus S|\neq 0$. By Assumption~\hyperlink{lab:A7}{A7}, the last term in \eqref{eq:nonoverfittingcond} is nonnegative, and hence
$$\frac{|M\setminus S|}{A_S}
= \frac{|M\setminus S|}{\gamma \Delta_{MS}+\frac{1-\gamma}{6}\mu_{Q_SM}}
\leq \bigg[ \gamma\sum_{j=1}^b\frac{|M_j \setminus S_j|}{|M\setminus S|}\kappa_j \bigg]^{-1} \leq \Big[ \gamma\min_{j=1,\ldots,b}\kappa_j \Big]^{-1} $$
where the last inequality follows from $\sum_{j=1}^b\frac{|M_j \setminus S_j|}{|M\setminus S|}=1$. By Assumption~\hyperlink{lab:A6}{A6} we have that $\min_j \kappa_j \to \infty$ as $n \rightarrow \infty$, and hence $|M\setminus S|=o(A_S)$. Thus, %We conclude that for every $M\in \mathcal M\setminus \{S\}$, 
by Lemma~\ref{lem:boundexpnc}, for any $\psi\in(0,1)$ and all $n$ large enough,
$
\E(NC(M))\;\leq\;e^{-\psi A_S}.
$

For the second step of the proof, let $A_S^*$ be the lower bound for $A_S$ given in \eqref{eq:nonoverfittingcond}. That is 
$$A_S^*\;:=\;\gamma\sum_{j=1}^b|M_j \setminus S_j|\kappa_j + \sum_{j=1}^b |S_j\setminus M_j |\left(\tfrac{1-\gamma}{6} n \rho(\bs X) {\beta^*_{\min,j}}^{2} - \gamma\kappa_j\right).$$
By \eqref{eq:nonoverfittingcond}, we have, for all $n$ large enough,
\begin{equation}\label{eq:boundexp}
    \E(NC(M))\;\leq\;e^{-\psi A_S^*}.
\end{equation}
Assumption~\hyperlink{lab:A7}{A7} implies there exist $g'_j \to \infty$ such that
\begin{equation}\label{eq:ogbea}
    \frac{(1-\gamma) n\rho(\bs X)}{6}{\beta_{\min,j}^*}^2 \,- \,\kappa_j \;=\; \ln(s_j) + g'_{j}.
\end{equation}
Let $\delta \in (0,1)$ and denote $\bar{m}_j=\max\big\{\frac{2\ln(p_j-s_j)}{f_j},\frac{2\ln(s_j)}{g'_{j}}\big\}$, where $f_j$ is given in Assumption~\hyperlink{lab:A6}{A6}. 
Take $\psi=\max_{j=1,\ldots,b}\frac{\xi + \delta + \bar{m}_j}{1+\bar{m}_j}$ for some $\xi \in(0, 1- \delta)$ then $\psi \in (0,1)$ and we have, for every $j=1,\ldots,b$,
\begin{align}
    &\psi \;>\; \frac{\delta + \frac{2\ln(p_j-s_j)}{f_j}}{1+\frac{2\ln(p_j-s_j)}{f_j}} \;=\; \frac{\delta f_j/2 + \ln(p_j-s_j)}{f_j/2 +\ln(p_j-s_j)} \label{eq:psif}\\
    &\psi \;>\; \frac{\delta + \frac{2\ln(s_j)}{g'_j}}{1+\frac{2\ln(s_j)}{g'_j}} \;=\;  \frac{\delta g'_j/2 + \ln(s_j)}{g'_j/2 +\ln(s_j)} \;\geq\; \frac{\delta g'_j/2 + \ln(s_j)}{g'_j +\ln(s_j)}.\label{eq:psig}
\end{align}
In Assumption~\hyperlink{lab:A7}{A7}, $\gamma$ is defined as $\gamma:=\tfrac12(1+\max_j\ln(p_j-s_j)/\kappa_j)$, we then have 
$$
\gamma\kappa_j\;\geq\;\frac12 \Big(1+\frac{\ln(p_j-s_j)}{\kappa_j}\Big)\kappa_j\;=\;\ln(p_j-s_j)+\frac12(\kappa_j-\ln(p_j-s_j)) \;=\; \ln(p_j-s_j)+ \frac12f_j.
$$
Hence, by \eqref{eq:psif}, we have
\begin{equation}\label{eq:ooofjfj}
\psi\gamma\kappa_j \;\geq\; \psi\big(\ln(p_j-s_j)+ \frac12f_j\big) \;\geq\;  \ln(p_j-s_j)+ \delta\frac12f_j.
\end{equation}
Further,
%By \hyperlink{lab:A7}{A7} and by our choice of $\psi$ and $\delta$,
\begin{equation}\label{eq:ooofjfj2}
\psi\left(\tfrac{1-\gamma}{6} n \rho(\bs X) {\beta^*_{\min,j}}^{2} - \gamma\kappa_j\right) \;\geq\; \psi\left(\ln(s_j)+g_j'\right) \;\geq\; \ln(s_j)+\delta \frac{1}{2}g'_j.
\end{equation}
where the first inequality follows from \eqref{eq:ogbea} and the second inequality from \eqref{eq:psig}.

In \eqref{eq:boundexp}, $\psi A^*_S=\sum_{j=1}^b|M_j \setminus S_j|\psi\gamma\kappa_j + \sum_{j=1}^b |S_j\setminus M_j |\psi\left(\tfrac{1-\gamma}{6} n \rho(\bs X) {\beta^*_{\min,j}}^{2} - \gamma\kappa_j\right)$. Then by \eqref{eq:ooofjfj} and \eqref{eq:ooofjfj2} , we get
\begin{equation}\label{eq:fyv}
\E(NC(M))\;\leq\;\exp\left\{- \sum_{j=1}^b |M_j\setminus S_j|(\ln(p_j-s_j)+\delta \tfrac{f_j}{2})-\sum_{j=1}^b |S_j\setminus M_j|(\ln(s_j)+\delta \tfrac{g'_j}{2})\right\}.    
\end{equation}

For the final step of the proof, denote $\mathcal{S}= \sum_{M \in \M\setminus \{S\}} \E\left(NC(M)\right)$ for convenience. By \eqref{eq:fyv} we have
\begin{align*}
    \mathcal{S} &\leq \sum_{M \in \M\setminus \{S\}}e^{- \sum_{j=1}^b |M_j\setminus S_j|\big(\ln(p_j-s_j)+\delta \tfrac{f_j}{2}\big)-\sum_{j=1}^b |S_j\setminus M_j|\big(\ln(s_j)+\delta \tfrac{g'_j}{2}\big)}.
\end{align*}
Observe that if $|M_j\setminus S_j|=0$ and $|S_j\setminus M_j|=0$ for all $j$, then $M=S$ and the summand in the right-hand side above is 1. Then by adding and resting 1 we get
$$\mathcal{S} \leq \sum_{M \in \M}e^{- \sum_{j=1}^b |M_j\setminus S_j|\big(\ln(p_j-s_j)+\delta \tfrac{f_j}{2}\big)-\sum_{j=1}^b |S_j\setminus M_j|\big(\ln(s_j)+\delta \tfrac{g'_j}{2}\big)}-1.$$
We can split the sum in the right-hand side above into sums over the models that have the same number of inactive variables and missing the same number of truly active variables in every block. That is, the models $M$ such that for all $j$, $|M_j\setminus S_j|=u_j$ and $|S_j\setminus M_j|=w_j$ with $u_j \in \{0,\ldots,p_j-s_j\}$ and $w_j \in \{0,\ldots,s_j\}$. Denote
$$S^{\bs u}_{\bs w}=\sum_{M \in \M:\forall \,j |M_j\setminus S_j|=u_j, |S_j\setminus M_j|=w_j } e^{- \sum_{j=1}^b u_j\big(\ln(p_j-s_j)+\delta \tfrac{f_j}{2}\big)-\sum_{j=1}^b w_j\big(\ln(s_j)+\delta \tfrac{g'_j}{2}\big)}.$$
We get
\begin{equation}\label{eq:uygfew}
    \mathcal{S} \leq -1+\sum_{w_1=0}^{s_1}\cdots\sum_{w_b=0}^{s_b}\sum_{u_1=0}^{p_1-s_1}\cdots\sum_{u_b=0}^{p_b-s_b}S^{\bs u}_{\bs w}.
\end{equation}
The number of models having, for all $j$, $u_j$ inactive parameters and missing $w_j$ out of the $s_j$ active parameters %with given collection of $u_j$'s and $w_j$'s 
is $\prod_{j=1}^b {\binom{p_j-s_j}{u_j}}{\binom{s_j}{w_j}}$. We thus have
\begin{align*}
    S^{\bs u}_{\bs w}&= \bigg(\prod_{j=1}^b {\binom{p_j-s_j}{u_j}}{\binom{s_j}{w_j}}\bigg)e^{- \sum_{j=1}^b u_j\big(\ln(p_j-s_j)+\delta \tfrac{f_j}{2}\big)-\sum_{j=1}^b w_j\big(\ln(s_j)+\delta \tfrac{g'_j}{2}\big)} \\
    &= \prod_{j=1}^b {\binom{p_j-s_j}{u_j}}e^{- u_j\big(\ln(p_j-s_j)+\delta \tfrac{f_j}{2}\big)}{\binom{s_j}{w_j}}e^{-w_j\big(\ln(s_j)+\delta \tfrac{g'_j}{2}\big)}.
\end{align*}
Inputting the expression above in \eqref{eq:uygfew} gives
\begin{align*}
    \mathcal{S} 
    &\leq -1+\sum_{w_1=0}^{s_1}\cdots\sum_{w_b=0}^{s_b}\sum_{u_1=0}^{p_1-s_1}\cdots\sum_{u_b=0}^{p_b-s_b}\prod_{j=1}^b {\binom{p_j-s_j}{u_j}}e^{- u_j\big(\ln(p_j-s_j)+\delta \tfrac{f_j}{2}\big)}{\binom{s_j}{w_j}}e^{-w_j\big(\ln(s_j)+\delta \tfrac{g'_j}{2}\big)}\\
    &\leq - 1 + \prod_{j=1}^b\left(1+\sum_{u_j=1}^{p_j-s_j}{\binom{p_j-s_j}{u_j}}e^{- u_j(\ln(p_j-s_j)+\delta \tfrac{f_j}{2})}\right)\left(1+
    \sum_{w_j=1}^{s_j}{\binom{s_j}{w_j}}
    e^{-w_j(\ln(s_j)+\delta \tfrac{g'_j}{2})}\right) 
\end{align*}
where the second inequality follows from first factorizing over terms in $u_j$ and $w_j$ and then taking the term in 0 out of every sum. A standard bound on binomial coefficient for $1\leq k\leq n$ is
\begin{equation}\label{eq:upperboudbinom}
    \binom{n}{k} \leq \left(\frac{n\,e}{k}\right)^k \leq \left(n\,e\right)^k = e^{k(\ln(n)+1)}.
\end{equation}
Then
\begin{equation}\label{eq:upperbound_sumprob}
    \mathcal{S} \;\;\leq\;\; - 1+\prod_{j=1}^b\left(1+\sum_{u_j=1}^{p_j-s_j}e^{-u_j\left(\delta \tfrac{f_j}{2}-1\right)}\right)\left(1+
    \sum_{w_j=1}^{s_j}    e^{-w_j\left(\delta \tfrac{g'_j}{2}-1\right)}\right) .
\end{equation}
Denote
$$
d_j\;=\;e^{1-\delta \tfrac{f_j}{2}},\qquad h_j\;=\;e^{1-\delta \tfrac{g'_j}{2}}
$$
where both expressions go to zero as $n$ increases since $f_j\to\infty$ and $g'_j\to\infty$.
For every $j$, by the properties of geometric sums, we have
\begin{align*}
    &1+\sum_{u_j=1}^{p_j-s_j}e^{-u_j\left(\delta \tfrac{f_j}{2}-1\right)}
    =\frac{1-d_j^{p_j-s_j+1}}{1-d_j} \\
    &1+\sum_{w_j=1}^{s_j}    e^{-w_j\left(\delta \tfrac{g'_j}{2}-1\right)} =\frac{1-h_j^{s_j+1}}{1-h_j}.
\end{align*}
Since both expressions converge to $1$ as $n$ grows, we get that $$\lim_{n \to \infty} S=\lim_{n \to \infty} \sum_{M \in \M\setminus \{S\}}\E\left(NC(M)\right)= 0.$$ By Lemma~\ref{lem:L0toL1convergence}, $P(\hat S^b\neq S)\leq \,2\, S$ and then $\lim_{n \to \infty} P(\hat S^b= S)=1$.

 \subsection{Proof of Theorem~\ref{prop:fullrecovreg}}

First, we prove the upper bound on $P(\hat{S}^b \neq S)$ in \eqref{eq:rateconvreg}. We assume here \hyperlink{lab:A1}{A1}, \hyperlink{lab:A6}{A6}, and \hyperlink{lab:A7}{A7}. Under the same assumptions, in the proof of Theorem~\ref{theo:suffcondlinearmodel}, the following bound was shown in \eqref{eq:upperbound_sumprob}:
% \piotr{I wrote this earlier already. Be careful when you use parts of proofs of some results rather than the results directly - this is not a common practice. In particular,  Theorem~\ref{theo:suffcondlinearmodel} assumed in addition \hyperlink{lab:A2}{A2}, which this one does not. }
\begin{equation}\label{eq:prodconstreg}
 \sum_{M \in \M\setminus \{S\}} \E\left(NC(M)\right)
    %\mathcal{S} 
    \;\;\leq\;\; - 1+\prod_{j=1}^b\left(1+\sum_{u_j=1}^{p_j-s_j}e^{-u_j\left(\delta \tfrac{f_j}{2}-1\right)}\right)\left(1+
    \sum_{w_j=1}^{s_j}    e^{-w_j\left(\delta \tfrac{g'_j}{2}-1\right)}\right) .
\end{equation}
Denote $S(u_j)=\sum_{u_j=1}^{p_j-s_j}e^{-u_j\left(\delta \tfrac{f_j}{2}-1\right)}$,  $S(w_j)=\sum_{w_j=1}^{s_j}    e^{-w_j\left(\delta \tfrac{g'_j}{2}-1\right)}$, $d_j\;=\;e^{1-\delta \tfrac{f_j}{2}}$, and $h_j\;=\;e^{1-\delta \tfrac{g'_j}{2}}$. For every $j$, we have, by the properties of geometric sums:
\begin{align*}
    &S(u_j)
    =d_j\,\frac{1-d_j^{p_j-s_j}}{1-d_j} \\
    &S(w_j) =h_j\,\frac{1-h_j^{s_j}}{1-h_j}.
\end{align*}
Developing the product in the right-hand side in \eqref{eq:prodconstreg} and reordering the resulting terms gives
$$\sum_{M \in \M\setminus \{S\}} \E\left(NC(M)\right)\;\;\leq\;\;-1+1+\sum_{j=1}^b\big[S(u_j)+S(w_j)\big]+\mathcal{R}$$
where all the terms in $\mathcal{R}$ are product of two or more of the sums $S(u_1),\ldots,S(u_b),S(w_1),\ldots,S(w_b)$. Given that $\delta>0$, $f_j \to \infty$ and $g_j' \to \infty$ by assumption, and hence $d_j \to 0$ and $h_j \to 0$, the $S(u_j)$ and $S(w_j)$ are smaller than 1 for all sufficiently large $n$ for all $j$. Then each of the $2^{2b}-2b-1$ terms in $\mathcal{R}$ is bounded above by $\sum_{j=1}^b\big[S(u_j)+S(w_j)\big]$ and we get, for every $n$ large enough,
 \begin{equation} \label{eq:fierg}
 \sum_{M \in \M\setminus \{S\}} \E\left(NC(M)\right)
 %\mathcal{S} 
 \leq (2^{2b}-2b)\sum_{j=1}^b \bigg[d_j\,\frac{1-d_j^{p_j-s_j}}{1-d_j}+ h_j\,\frac{1-h_j^{s_j}}{1-h_j}\bigg ]. \end{equation} 
 We have $d_j \to 0$ and $h_j \to 0$, then for every $n$ large enough $\frac{1-d_j^{p_j-s_j}}{1-d_j}\to1$ and $\frac{1-h_j^{s_j}}{1-h_j}\to1$ for all $j$. Since $r>1$, we get, for every $n$ large enough,
\begin{equation} \label{eq:fierg2}r>\max_{j=1,\ldots,b}\bigg\{\frac{1-d_j^{p_j-s_j}}{1-d_j} \,,\,\frac{1-h_j^{s_j}}{1-h_j} \bigg\}.\end{equation} 
 By Lemma~\ref{lem:L0toL1convergence}, \eqref{eq:fierg} and \eqref{eq:fierg2}, we then obtain:
$$P(\hat{S}^{b} \neq S)\;\leq\;  2 \sum_{M \in \M\setminus \{S\}} \E\left(NC(M)\right)
     \;\leq\; 2 (2^{2b}-2b) r\,e\,\,\sum_{j=1}^b e^{-\delta \tfrac{f_j}{2}} +e^{-\delta \tfrac{g'_j}{2}}$$
By the definition of $f_j$ in Assumption~\hyperlink{lab:A6}{A6}, that of $g'_j$ in \eqref{eq:ogbea}, and the fact that $2e<6$,  we get, for every $n$ large enough, that
\begin{equation}\label{eq:gpkr}
    P(\hat{S}^{b} \neq S)\;\leq\; (2^{2b}-2b)6r\,\,\sum_{j=1}^b\,\, e^{-\tfrac{\delta}{2}\big[\kappa_j-\ln(p_j-s_j)\big]} + e^{-\tfrac{\delta}{2}\big[\frac{(1-\gamma) n\rho(\bs X)}{6}{\beta_{\min,j}^*}^2 - \kappa_j -\ln(s_j) \big]}.
\end{equation}

We have $\frac{(1-\gamma) n\rho(\bs X)}{6}{\beta_{\min,j}^*}^2 - \kappa_j > \Big(\sqrt{\frac{(1-\gamma) n\rho(\bs X)}{6}}{\beta_{\min,j}^*} - \sqrt{\kappa_j}\Big)^2$ and then:
$$P(\hat{S}^{b} \neq S)\;\leq\; (2^{2b}-2b)6r\,\,\sum_{j=1}^b\,\, e^{-\tfrac{\delta}{2}\big[\kappa_j-\ln(p_j-s_j)\big]} + e^{-\tfrac{\delta}{2}\big[(\sqrt{\frac{(1-\gamma) n\rho(\bs X)}{6}}{\beta_{\min,j}^*} - \sqrt{\kappa_j})^2 -\ln(s_j) \big]}
$$
which proves \eqref{eq:rateconvreg}. 

Second, we prove that, if for all $j=1,\ldots,b$, it holds that
\begin{equation}\label{betamincond:minimalreg}
    \lim_{n\to \infty}\frac{\sqrt{(1-\gamma) n\rho(\bs X)/3}\;\beta_{\min,j}^*}{\sqrt{2\ln(p_j-s_j)} + \sqrt{2\ln(s_j)}} >1
\end{equation}
then $\kappa^*_j$ defined in \eqref{eq:oraclepenreg} satisfies Assumptions~\hyperlink{lab:A6}{A6} and \hyperlink{lab:A7}{A7}. The proof is essentially the same as the proof of the second part of Theorem~\ref{theo:fullrecovGS}, replacing $\tau^*_j$ by $\sqrt{\kappa^*_j}$.  Under \eqref{betamincond:minimalreg}, for every $j$, there exists a sequence $c$ such that $\lim_{n\to\infty} c>1$ and $\sqrt{\frac{(1-\gamma)n}{6}}\beta^*_{\min,j}=c\big(\sqrt{\ln(p_j-s_j)} + \sqrt{\ln(s_j)}\big)$. Denote $a=\sqrt{\ln(p_j-s_j)}$ and $b=\sqrt{\ln(s_j)}$. Proceeding as in the proof of Theorem~\ref{theo:fullrecovGS} shows that:
\begin{equation*}
\sqrt{\kappa^*_j} =\frac{c^2+1}{2c}a+\frac{c^2-1}{2c}b  \geq  \Big(\frac{c^2+1}{2c}+1-1\Big)a = \Big(1+\frac{(c-1)^2}{2c}\Big)\sqrt{\ln(p_j-s_j)} 
\end{equation*}
which implies Assumption~\hyperlink{lab:A7}{A7} since $\lim_{n\to\infty} c>1$. We also have
\begin{equation*}
\frac{\sqrt{(1-\gamma) n\rho(\bs X)}}{6}{\beta_{\min,j}^*}-\sqrt{\kappa^*} = \frac{c^2-1}{2c}a+\frac{c^2+1}{2c}b \geq \Big(1+\frac{(c-1)^2}{2c}\Big)\sqrt{\ln(s_j)}.
\end{equation*}
which implies Assumption~\hyperlink{lab:A6}{A6} since $\lim_{n\to\infty} c>1$.

Finally, we prove \eqref{eq:oraclerateconvreg}. Since Assumptions~\hyperlink{lab:A6}{A6} and \hyperlink{lab:A7}{A7} hold for the $\kappa^*_j$, and because we assume \hyperlink{lab:A1}{A1}, by the first part of the theorem, for any $r> 1$ and every $n$ large enough,
\begin{equation*}
 P(\hat{S}^{b} \neq S)\;\leq\; 6(2^{2b}-2b)r\,\sum_{j=1}^b\,\, e^{-\tfrac{\delta}{2}\big[\kappa_j^*-\ln(p_j-s_j)\big]} + 
    e^{-\tfrac{\delta}{2}\big[(\sqrt{\frac{(1-\gamma) n\rho(\bs X)}{6}}{\beta_{\min,j}^*} - \sqrt{\kappa_j^*})^2 -\ln(s_j) \big]}.
\end{equation*}
Simple algebra shows that the $\kappa_j^*$ satisfy $\kappa^*_j-\ln(p_j-s_j)=\big(\sqrt{\frac{(1-\gamma) n\rho(\bs X)}{6}}{\beta_{\min,j}^*} - \sqrt{\kappa_j^*}\big)^2 -\ln(s_j)$ for all $j$. We then have
\begin{equation}\label{eq:afi}
 P(\hat{S}^{b} \neq S)\;\leq\; 12(2^{2b}-2b)r\,\sum_{j=1}^b\,\, e^{-\tfrac{\delta}{2}\big[\kappa_j^*-\ln(p_j-s_j)\big]}.
\end{equation}
 Proceeding as in the proof of Theorem~\ref{theo:fullrecovGS}, denote \\$d =\frac12\sqrt{\frac{6}{(1-\gamma) n\rho(\bs X)}}\frac{1}{\beta_{\min,j}^*}\big(\ln({p_j-s_j})-\ln({s_j})\big)$ such that $\sqrt{\kappa^*_j}=  \frac12\sqrt{\frac{(1-\gamma) n\rho(\bs X)}{6}}{\beta_{\min,j}^*} + d$. Then
$$e^{-\frac{\delta}{2}\left(\kappa^*_j-\ln(p_j-s_j)\right)}
=e^{-\frac{\delta}{2}\left[\frac{(1-\gamma) n\rho(\bs X)}{24}{\beta_{\min,j}^*}^2+d^2 - \frac{1}{2}\left(\ln (p_j- s_j)+\ln (s_j)\right)\right]}.$$
By considering separately the two possible maxima in $\ln \max\{p_j-s_j,s_j\}$, we get that
\begin{equation}\label{eq:oig}
    \frac{e^{-\frac{\delta}{2}\left(\kappa^*_j-\ln(p_j-s_j)\right)}}{e^{-\frac{\delta}{2}\left[\frac{(1-\gamma) n\rho(\bs X)}{24}{\beta_{\min,j}^*}^2-\ln \max\{p_j-s_j,s_j\}\right]}}= 
e^{-\frac{\delta}{2}\left[d^2 + \frac{1}{2}\left|\ln (p_j- s_j)-\ln (s_j)\right|\right]} <1.
\end{equation}
It follows that, by \eqref{eq:afi} and \eqref{eq:oig},
$$P(\hat{S}^{b} \neq S) \; \leq \; 12(2^{2b}-2b)r\,\sum_{j=1}^b\,\, e^{-\frac{\delta}{2}\left[\frac{(1-\gamma) n\rho(\bs X)}{24}{\beta_{\min,j}^*}^2-\ln \max\{p_j-s_j,s_j\}\right]}.$$
which proves \eqref{eq:oraclerateconvreg}.

\subsection{Proof of Proposition~\ref{prop:nectypeIreg}}

The event $\hat{S}^b=S$ (correct recovery of $S$) requires the event 
%Recovery of $S$ requires that 
$\max_{M\in O_j}\frac{NC(M)}{NC(S)}<1$ ($S$ is preferred to any model in $O_j$, i.e. over-fitting $S$ by 1 variable in block $j$). % with probability going to 1 as n grows. 
Using the definition of the normalized criterion $NC$ in \eqref{eq:CM}, and that $C(S)-C(M)=L_{SM}/2 + \Delta_{MS}$ for $L_{SM}$ and $\Delta_{MS}$ defined in \eqref{eq:DeltaM}, we obtain
%    NC(M)\;:=\;\frac{e^{{\rm C}(M)}}{\underset{M' \in \M}{\sum} e^{{\rm C}(M')}}.
\begin{align}
P(\hat{S}^b=S) 
\leq P\Big(\max_{M\in O_j}\frac{NC(M)}{NC(S)}<1\Big)
= P \left( \max_{M\in O_j} e^{C(M)-C(S)} <1 \right)
\nonumber \\
= P \left( \max_{M\in O_j} e^{\frac{1}{2} L_{MS} + \Delta_{SM}} < 1 \right)
=
P\big(\max_{M\in O_j}\sqrt{L_{MS}}<\sqrt{2\kappa_j}\Big).
\nonumber
\end{align}
By Lemma~\ref{lemma:s7}, for every $M\in O_j$, $L_{MS}\in \chi^2_1$ and there exists $Z_M\sim N(0,1)$ such that $\sqrt{L_{MS}}=|Z_M|$. Since for every $M\in O_j$, $|Z_M|\geq Z_M$, we have
\begin{equation}\label{eq:upperboundrecov}
    P(\hat{S}^b=S) \leq  1-P\big(\max_{M\in O_j}Z_M\geq\sqrt{2\kappa_j}\Big).
\end{equation}

The set $O_j$ has cardinality $p_j-s_j$, then by Theorem 3.4 in \cite{hartigan} and our Assumption~\hyperlink{lab:A2}{A2}, for any $\varepsilon>0$,
$$P\big(\max_{M\in O_j}Z_M \geq \underline{\lambda}_j \sqrt{2 \ln(p_j-s_j)}(1-\varepsilon)\big) \to 1,$$
where $\underline{\lambda}_j$ is as defined prior to the statement of Proposition~\ref{prop:nectypeIreg}.
If $\lim_{n\to\infty}\frac{\kappa_j}{\underline{\lambda}_j^2\ln(p_j-s_j)} <1$, then $\lim_{n \to \infty} P\big(\max_{M\in O_j}Z_M\geq\sqrt{2\kappa_j}\big) = 1$. Hence, by \eqref{eq:upperboundrecov} we have that $\lim_{n \to \infty} P(\hat{S}^b=S) = 0$, as we wished to prove.

\subsection{Proof of Lemma~\ref{lem:upperboundnoncentral}}
The result follows directly from Lemma~\ref{lem:upperboundnoncentralgen}, with $M=S$. 


\subsection{Proof of Proposition~\ref{prop:nectypeIIreg}}

Let $M=S\setminus S^{S}(\kappa)$. Since $S^{S}(\kappa)\neq \emptyset$ by assumption and $S^S(\kappa)\subseteq S$, we have that $M\neq S$. With this notation, we aim at proving that $\lim_{n\to \infty}P(NC(M)<NC(S))=\lim_{n\to \infty}P(\hat S^b=S)=0$. The event $\hat{S}^b=S$ (correct selection of $S$) implies that $NC(M)< NC(S)$ (preferring $S$ over $M$), and hence 
\begin{align}\label{eq:nectypeIIreg_bound}
P(\hat{S}^b=S)\;\leq\; P\left( NC(M) < NC(S) \right)
\;=\; P(L_{SM}> 2\Delta_{SM}),
\end{align}
where we used the definition of $NC$ in
\eqref{eq:CM}, and that $C(S)-C(M)=L_{SM}/2 + \Delta_{MS}$ for $L_{SM}$ and $\Delta_{MS}$ defined in \eqref{eq:DeltaM}. It suffices then to show that the right-hand side in \eqref{eq:nectypeIIreg_bound} converges to 0 as $n \to \infty$.
To do this, we note that, by Lemma~\ref{lemma:s7}, we have $L_{SM}\sim \chi_{|S\setminus M|}^2(\mu_{SM})$. We then use the non-central chi-square bound in Lemma~\ref{lemma:s1s3}.
%We  bound next $P(L_{SM}> 2\Delta_{SM})$ with Lemma~\ref{lemma:s1s3}. 

To apply Lemma~\ref{lemma:s1s3}, we first show that the degrees of freedom satisfy $|S\setminus M|=o(2\Delta_{SM})$ and also the non-centrality parameter is such that $\mu_{SM}=o(2\Delta_{SM})$.
We have that $\Delta_{SM}=\sum_{j=1}^b|S_j \setminus M_j|\kappa_j$ and then
$$\frac{|S\setminus M|}{2\Delta_{SM}}\;\leq\; \bigg[\sum_{j=1}^b\frac{2|S_j \setminus M_j|}{|S\setminus M|}\kappa_j \bigg]^{-1} \;\leq\; \Big[ 2\min_{j=1,\ldots,b}\kappa_j \Big]^{-1} $$
where the last inequality follows from $\sum_{j=1}^b\frac{|S_j \setminus M_j|}{|S\setminus M|}=1$. By Assumption~\hyperlink{lab:A6}{A6}, $\kappa_j\to \infty$ for all $j=1,\ldots,b$, then the left-hand side above goes to 0 as $n$ grows and $|S\setminus M|=o(2\Delta_{SM})$. Further, by Lemma~\ref{lem:upperboundnoncentral}, we also have
\begin{equation}\label{eq:AGEoj}
    \frac{\mu_{SM}}{2\Delta_{SM}}\leq \frac{ \sum_{j=1}^b |S_j\setminus M_j|n\bar{\lambda}\max_{i\in S_j\setminus M_j}{\beta_i^*}^2}{\sum_{j=1}^b |S_j\setminus M_j|2 \kappa_j}.
\end{equation}
Let $\tilde{r}:=\sum_{j=1}^b n\bar{\lambda}\max_{i\in S_j\setminus M_j}{\beta_i^*}^2/(2\kappa_j)$. We show next that $\tilde{r}$ is an upper bound on $\mu_{SM}/(2\Delta_{SM})$. By restricting the sum in $\tilde{r}$ to the $j$ such that $|S_j\setminus M_j|\neq 0$ and multiplying the numerator and denominator of the summand by $|S_j\setminus M_j|$, we get the lower bound on $\tilde{r}$
\begin{eqnarray}\label{eq:vfej}
    \tilde{r}
    &\geq &\sum_{j=1,|S_j\setminus M_j|\neq 0}^b\frac{|S_j\setminus M_j|n\bar{\lambda}\max_{i\in S_j\setminus M_j}{\beta_i^*}^2}{|S_j\setminus M_j|2\kappa_j} .\end{eqnarray}
Note that for any collections $(\alpha_j,\delta_j)\in  \R\times \R\setminus\{0\}$, $j=1,\ldots,b$, we have
\begin{equation}\label{eq:truc}
    \sum_{j=1}^b\frac{\alpha_j}{\delta_j}=
\sum_{j=1}^b\frac{\alpha_j \frac{1}{\delta_j} (\delta_j + \sum_{l \neq j} \delta_l)}{\sum_{j=1}^b \delta_j}=
    \frac{\sum_{j=1}^b\alpha_j(1+\sum_{l\neq j} \frac{\delta_l}{\delta_j})}{\sum_{j=1}^b\delta_j}
\end{equation}
Using the above in the right-hand side of \eqref{eq:vfej} gives
\begin{eqnarray*}
    \tilde{r}
    &\geq& \frac{\sum_{j=1,|S_j\setminus M_j|\neq 0 }^b|S_j\setminus M_j|n\bar{\lambda}\max_{i\in S_j\setminus M_j}{\beta_i^*}^2\Big(1+\sum_{l\neq j}\frac{|S_l\setminus M_l|2\kappa_l}{|S_j\setminus M_j|2\kappa_j}\Big)}{\sum_{j=1,|S_j\setminus M_j|\neq 0 }^b|S_j\setminus M_j|2\kappa_j} \\
    &\geq & \frac{\sum_{j=1}^b|S_j\setminus M_j|n\bar{\lambda}\max_{i\in S_j\setminus M_j}{\beta_i^*}^2}{\sum_{j=1}^b \big|S\setminus M_j\big|2\kappa_j}
\end{eqnarray*}
where last inequality follows from $\Big(1+\sum_{l\neq j}\frac{|S_l\setminus M_l|2\kappa_l}{|S_j\setminus M_j|2\kappa_j}\Big)\geq 1$ for all $j$ and that $\sum_{j=1}^b \big|S\setminus M_j\big|2\kappa_j=\sum_{j=1,|S_j\setminus M_j|\neq 0 }^b|S_j\setminus M_j|2\kappa_j$. Then by \eqref{eq:AGEoj}
$$\frac{\mu_{SM}}{2\Delta_{SM}}\;\leq \;\tilde{r}\;=\;\sum_{j=1}^b\frac{n\bar{\lambda}\max_{i\in S_j\setminus M_j}{\beta_i^*}^2}{2\kappa_j}.$$
For every $j=1,\ldots,b$, $S_j\setminus M_j \subset S_j^{S}(\kappa)$ , then by definition of the $S_j^{S}(\kappa)$, the right-hand 
side above goes to 0 as $n \to \infty$ and  $\mu_{SM}=o(2\Delta_{SM})$. We can now use Lemma~\ref{lemma:s1s3}. %Lemma~\ref{lemma:s7}, 
For any $\phi \in (0,1)$ and every $n$ large enough,
\begin{align}
    P(L_{SM}> 2\Delta_{SM}) \leq e^{-\phi2\Delta_{SM}} = e^{-\phi2\sum_{j=1}^b|S(\kappa)_j^S|\kappa_j}
\nonumber
\end{align} 
where the right hand-side goes to 0 since for all $j=1,\ldots,b$, $\kappa_j \to \infty$. It follows by \eqref{eq:nectypeIIreg_bound} that $\lim_{n\to\infty}P\left( NC(M)/NC(S) < 1  \right)=\lim_{n\to\infty}P(\hat S^b=S)=0$ as we wished to prove.

\subsection{Proof of Corollary~\ref{cor:necrecovreg}}
By assumption we have that, for some $j\in\{1,\ldots,b\}$, 
$$\lim_{n\to\infty}\frac{\sqrt{n\bar{\lambda}}\beta_{\min,j}^*}{\sqrt{\underline{\lambda}_j\ln(p_j-s_j)}}= \lim_{n\to\infty}\frac{\sqrt{n\bar{\lambda}}\beta_{\min,j}^*}{\sqrt{\kappa_j}}\sqrt{\frac{\kappa_j}{\underline{\lambda}_j\ln(p_j-s_j)}} = 0.$$
If $\lim_{n\to\infty}\frac{\sqrt{n\bar{\lambda}}\beta_{\min,j}^*}{\sqrt{\kappa_j}}>0$, then $\lim_{n\to\infty}\sqrt{\frac{\kappa_j}{\underline{\lambda}_j\ln(p_j-s_j)}} = 0$ and, by Proposition~\ref{prop:nectypeIreg}, it follows that $\lim_{n \to \infty} P(\hat{S}^b=S) < 1$. 
If $\lim_{n\to\infty}\frac{\sqrt{n\bar{\lambda}}\beta_{\min,j}^*}{\sqrt{\kappa_j}}=0$, then by Proposition~\ref{prop:nectypeIIreg} $\lim_{n \to \infty} P(\hat{S}^b=S) < 1$, as we wished to prove.

\subsection{Proof of Corollary~\ref{cor:selconstimpossiblereg}}
Since $\hat{S}$ is $\hat{S}^b$ with $b=1$, the assumptions of Corollary~\ref{cor:necrecovreg} for $\hat{S}$ are met and $\lim_{n \to \infty} P(\hat{S} = S) < 1$. Since Assumptions~\hyperlink{lab:A6}{A6} and \hyperlink{lab:A7}{A7} hold, by Theorem~\ref{theo:suffcondlinearmodel} we have that $\lim_{n \to \infty} P(\hat{S}^b = S)= 1$.

\subsection{Proof of Theorem~\ref{thm:convtoT}}
The proof strategy is the same as that of Theorem~\ref{theo:suffcondlinearmodel}, with suitable adjustments. 
The first step is to use Lemma~\ref{lem:boundexpnc} to bound $\E(NC(M))$ for every $M\not\in \T(\kappa)$. The main difference is that %we replace $S$ 
in the proof of Theorem~\ref{theo:suffcondlinearmodel} 
we took $T=S$ in Lemma~\ref{lem:boundexpnc}, whereas now we take  a model $T=T^M\in \T(\kappa)$
%by a model in $T^M\in \T(\kappa)$ 
that depends on $M$. 
Intuitively, $T^M$ contains large truly non-zero parameters that are missed by $M$, and hence $T^M$ should be chosen over $M$ asymptotically. More precisely,
we choose $T^M\in \T(\kappa)$ such that %for every $M\not\in \T(\kappa)$, 
$T^M\setminus M\subseteq S^L(\kappa)$ and the elements in $M\setminus T^M$ are either inactive or in $S^S(\kappa)$. The latter condition %observation 
and Assumption~\hyperlink{lab:A6}{A6} ensure that the assumptions of Lemma~\ref{lem:boundexpnc} are met. We then get a bound $\E(NC(M)) \leq e^{-\psi A_{T^M}}$ for every large enough $n$ and any $\psi \in (0,1)$, where $A_{T^M}=\gamma\Delta_{MT^M}+\tfrac{1-\gamma}{6}\mu_{Q_{T^M}M}$ (\emph{cf} \eqref{eq:DeltaM} and \eqref{eq:muQM}), $\gamma\in(1/2,1)$ is defined in \eqref{eq:defSes} and $Q_{T^M}= M\cup T^M$. The second step is to obtain a lower bound for $A_{T^M}$, which gives an upper bound for $\E(NC(M))$ (distinct to that obtained in the proof of Theorem~\ref{theo:suffcondlinearmodel}).
The final step is to %use these bounds to 
get an upper-bound $\sum_{M \in \M\setminus \T(\kappa)} \E\left(NC(M)\right)$ that vanishes (as $n$ grows) under the assumptions of Theorem~\ref{thm:convtoT}. Then Lemma~\ref{lem:L0toL1convergencegen} immediately implies that $P(\hat S^b\not \in \T(\kappa))$ also vanishes.

For the first step of the proof, recall that the set $\T(\kappa)$ contains all the models that are the union of $S^L(\kappa)$ (large signals) and some subset of $S^{I}(\kappa)$ (intermediate signals). For any $M \not\in\T(\kappa)$, take the unique $T^M\in\T(\kappa)$ such that $T^M = \big[M\cap S^{I}(\kappa)\big]\cup S^L(\kappa) $, which implies $M\cap S^{I}(\kappa)=T^M\cap S^{I}(\kappa)$. 
That is, $T^M$ contains all the large signals plus the intermediate signals in $M$.
To show that $\E(NC(M)) \leq e^{-\psi A_{T^M}}$ we show that $A_{T^M}$ satisfies the conditions of Lemma~\ref{lem:boundexpnc}, taking $T=T^M$. That is, we wish to show that three conditions hold: $A_{T^M}>0$, $|M\setminus T^M|=o(A_{T^M})$, and $\mu_{Q_{T^M} T^M}=o(A_{T^M})$.

Observe that $\Delta_{MT^M}$, defined in \eqref{eq:DeltaM}, can be rewritten as $\Delta_{MT^M}=\sum_{j=1}^b(|M_j\setminus T^M_j|-|T^M_j\setminus M_j|)\kappa_j$ and then:
$$ A_{T^M}=\sum_{j=1}^b|M_j\setminus T^M_j|\gamma\kappa_j+\tfrac{1-\gamma}{6}\mu_{Q_{T^M}M}-\sum_{j=1}^b|T^M_j\setminus M_j|\gamma\kappa_j $$
Since $\gamma <1$, $(1-\gamma)/6 >0$, by  Lemma~\ref{lem:boundrhogen}, for every $n\in \N$ we have 
\begin{equation}\label{eq:nonoverfittingcond2}
   A_{T^M}\;\geq\; \sum_{j=1}^b|M_j \setminus T^M_j|\gamma\kappa_j + \sum_{j=1}^b |T^M_j\setminus M_j |\Big(\tfrac{1-\gamma}{6} n \rho(\bs X) \min_{i\in T^M_j\setminus M_j}{\beta^*_i}^{2} - \gamma\kappa_j\Big).
\end{equation} 
We have that $T^M\subseteq \big(S^I(\kappa) \cup S^L(\kappa)\big)$ since $T^M\in\T(\kappa)$ and that $T^M\cap S^{I}(\kappa) = M\cap S^{I}(\kappa)$. It follows that $T^M\setminus M \subseteq S^L(\kappa)$. By definition of $S^L(\kappa)$, the rightmost component in \eqref{eq:nonoverfittingcond2} is nonnegative and, if $|T^M\setminus M|\neq0$, it is positive, for every $n$ large enough. By Assumption~\hyperlink{lab:A6}{A6}, component $\gamma\sum_{j=1}^b|M_j \setminus T^M_j|\kappa_j$ is nonnegative, and, if $|M\setminus T^M|\neq0$, positive. Now, $M\neq T^M$ implies that necessarily $|M\setminus T^M|\neq0$ or $|T^M\setminus M|\neq0$ and then, for every $n$ large enough, $A_{T^M}>0$, establishing the first condition required by Lemma~\ref{lem:boundexpnc}. Regarding its second condition, if $|M\setminus T^M|=0$, we immediately have $|M\setminus T^M|=o(A_{T^M})$. If $|M\setminus T^M|\neq 0$, since the rightmost component in \eqref{eq:nonoverfittingcond2} is nonnegative, we have
$$\frac{|M\setminus T^M|}{A_{T^M}}\leq \bigg[ \gamma\sum_{j=1}^b\frac{|M_j \setminus T^M_j|}{|M\setminus T^M|}\kappa_j \bigg]^{-1} \leq \bigg[ \gamma\min_{j=1,\ldots,b}\kappa_j \bigg]^{-1} $$
where the last inequality follows from $\sum_{j=1}^b\frac{|M_j \setminus T^M_j|}{|M\setminus T^M|}=1$. By Assumption~\hyperlink{lab:A6}{A6}, $\min_{j=1,\ldots,b}\kappa_j\to \infty$ as $n\to\infty$ and hence $|M\setminus T^M|=o(A_{T^M})$ when $|M\setminus T^M|\neq 0$ too.

Finally, consider the third condition in Lemma~\ref{lem:boundexpnc}.
In the proof of Theorem~\ref{theo:suffcondlinearmodel}, $\mu_{Q_SS}=o(A_S)$ is immediate because $ Q_S \setminus S= M\setminus S\subseteq S^C$, i.e. since all parameters in $M \setminus S$ are truly zero we have that $\mu_{Q_SS}=0$.
Here $M\setminus T^M$ is not necessarily a subset of $S^C$, hence $\mu_{Q_T^M T^M} \geq 0$. Note that, since $T^M\in\T(\kappa)$, we have that $S^L(\kappa)\subseteq T^M$. Moreover we have $M\cap S^{I}(\kappa)=T^M\cap S^{I}(\kappa)$. It follows that $M\setminus T^M\subseteq (S^{I}(\kappa) \cup S^L(\kappa))^C$, that is the elements of $M\setminus T^M$ are either inactive or belong to $S^S(\kappa)$. If $M\cap S^{S}(\kappa)=\emptyset$ then $M\setminus T^M\subseteq S^C$ and we immediately get $\mu_{Q_{T^M} T^M}=o(A_{T^M})$ because $\bs\beta^*_{M \setminus T^M}=\mu_{Q_{T^M} T^M}=0$. Assume now that $M\cap S^{S}(\kappa)\neq\emptyset$. Using that the rightmost component in \eqref{eq:nonoverfittingcond2} is nonnegative and Lemma~\ref{lem:upperboundnoncentralgen}, we have
$$\frac{\mu_{Q_{T^M} T^M}}{ A_{T^M}} \leq \frac{\mu_{Q_{T^M} T^M}}{ \gamma\sum_{j=1}^b|M_j \setminus T^M_j|\kappa_j} \leq \frac{n\,\bar{\lambda}\,\sum_{j=1}^b|(S_j\cap M_j )\setminus T^M_j|\,\max_{i\in (S_j\cap M_j )\setminus T^M_j}{\beta_{i}^*}^2}{ \gamma\sum_{j=1}^b|M_j \setminus T^M_j|\kappa_j}.$$
Observe that for all $j=1,\ldots,b$, $M_j \setminus T^M_j \subseteq (S_j\cap M_j )\setminus T^M_j$ and then $|M_j \setminus T^M_j| \geq |(S_j\cap M_j )\setminus T^M_j|$. We then get
$$\frac{\mu_{Q_{T^M} T^M}}{ A_{T^M}} \leq \frac{n\,\bar{\lambda}\,\sum_{j=1}^b|(S_j\cap M_j )\setminus T^M_j|\,\max_{i\in(S_j\cap M_j )\setminus T^M_j}{\beta_{i}^*}^2}{ \gamma\sum_{j=1}^b|(S_j\cap M_j ) \setminus T^M_j|\kappa_j}.$$
Moreover, since  $M\setminus T^M\subseteq (S^{I}(\kappa) \cup S^L(\kappa))^C$ as discussed earlier, we have, for all $j=1,\ldots,b$, $(S_j\cap M_j )\setminus T^M_j \subseteq S^{S}_j(\kappa)$. It follows $\max_{i\in(S_j\cap M_j )\setminus T^M_j}{\beta_{i}^*}^2 \leq \max_{i\in S_j^S}{\beta_{i}^*}^2$ and
\begin{equation}\label{eq:onbfsa}
  \frac{\mu_{Q_{T^M} T^M}}{ A_{T^M}} \leq \frac{n\,\bar{\lambda}\,\sum_{j=1}^b|(S_j\cap M_j )\setminus T^M_j|\,\max_{i\in S^{S}_j(\kappa)}{\beta_{i}^*}^2}{ \gamma\sum_{j=1}^b|(S_j\cap M_j ) \setminus T^M_j|\kappa_j}.  
\end{equation}
Let $\bar{r}:=\sum_{j=1}^b n\bar{\lambda}\max_{i\in S^{S}_j(\kappa)}{\beta_i^*}^2/(\gamma\kappa_j)$. We show next that $\bar{r}$ is an upper bound on $\mu_{Q_{T^M} T^M}/ A_{T^M}$. By restricting the sum in $\bar{r}$ to the $j$ such that $|(S_j\cap M_j )\setminus T^M_j|\neq 0$ and multiplying the numerator and denominator of the summand by $|(S_j\cap M_j )\setminus T^M_j|$, we get 
\begin{eqnarray}\label{eq:ovrwS}
    \bar{r}
    &\geq &\sum_{j=1,|(S_j\cap M_j )\setminus T^M_j|\neq 0}^b\frac{|(S_j\cap M_j )\setminus T^M_j|n\bar{\lambda}\max_{i\in S^{S}_j(\kappa)}{\beta_i^*}^2}{|(S_j\cap M_j )\setminus T^M_j|\gamma\kappa_j}
\end{eqnarray}
Using the property of \eqref{eq:truc} in the right-hand side of \eqref{eq:ovrwS}, we get
\begin{align*}
     \bar{r}
    &\geq \frac{\sum_{j=1,|(S_j\cap M_j )\setminus T^M_j|\neq 0 }^b|(S_j\cap M_j )\setminus T^M_j|n\bar{\lambda}\max_{i\in S^{S}_j(\kappa)}{\beta_i^*}^2\Big(1+\sum_{l\neq j}\frac{|(S_l\cap M_l)\setminus T^M_l|\gamma\kappa_l}{|(S_j\cap M_j )\setminus T^M_j|\gamma\kappa_j}\Big)}{\sum_{j=1,|(S_j\cap M_j )\setminus T^M_j|\neq 0}^b|(S_j\cap M_j )\setminus T^M_j|\gamma\kappa_j} \\
    &\geq \frac{\sum_{j=1}^b|(S_j\cap M_j )\setminus T^M_j|n\bar{\lambda}\max_{i\in S^{S}_j(\kappa)}{\beta_i^*}^2}{\sum_{j=1}^b \big|(S\cap M_j)\setminus T^M_j\big|\gamma\kappa_j}
\end{align*}
where last inequality follows from $\Big(1+\sum_{l\neq j}\frac{|(S_l\cap M_l)\setminus T^M_l|\gamma\kappa_l}{|(S_j\cap M_j )\setminus T^M_j|\gamma\kappa_j}\Big)\geq 1$ for all $j$ and from the identity $\sum_{j=1,|(S_j\cap M_j )\setminus T^M_j|\neq 0}^b|(S_j\cap M_j )\setminus T^M_j|\gamma\kappa_j=\sum_{j=1}^b \big|(S\cap M_j)\setminus T^M_j\big|\gamma\kappa_j$. Then, %when $M\cap S^{S}(\kappa)\neq\emptyset$, 
by \eqref{eq:onbfsa},
$$\frac{\mu_{Q_{T^M} T^M}}{ A_{T^M}} \;\leq\; \bar{r} \;= \;\sum_{j=1}^b\frac{n\bar{\lambda}\max_{i\in S^{S}_j(\kappa)}{\beta_i^*}^2}{\gamma\kappa_j}.$$
By definition of $S_j^{S}(\kappa)$, $n\bar{\lambda}\max_{i\in S^{S}_j(\kappa)}{\beta_i^*}^2=o(\kappa_j)$ for all $j$ and $\mu_{Q_{T^M} T^M}=o(A_{T^M})$ when $M\cap S^{S}(\kappa)\neq\emptyset$ too.
We can now apply Lemma~\ref{lem:boundexpnc} and get that for every $M\in \mathcal M\setminus \T(\kappa)$, for any $\psi\in(0,1)$ and every $n$ large enough, $\E(NC(M))\;\leq\;e^{-\psi A_{T^M}}$.

The second step of the proof is to lower-bound $A_{T^M}=\gamma\Delta_{MT^M}+\tfrac{1-\gamma}{6}\mu_{Q_{T^M}M}$. %For the second step of the proof, 
Let
$$A_{T^M}^*\;:=\;\gamma\sum_{j=1}^b|M_j \setminus T^M_j|\kappa_j + \sum_{j=1}^b |T^M_j\setminus M_j |\left(\tfrac{1-\gamma}{6} n \rho(\bs X) \min_{i\in S^{L}_j(\kappa)}{\beta^*_i}^{2} - \gamma\kappa_j\right) $$
Recall that, for every $j$, $T^M_j\setminus M_j\subseteq S^{L}_j(\kappa)$. We have then $\min_{i\in T^M_j\setminus M_j}{\beta^*_i}^{2}\geq \min_{i\in S^{L}_j(\kappa)}{\beta^*_i}^{2}$, and by \eqref{eq:nonoverfittingcond2}, $A_{T^M} \geq A_{T^M}^*$. It follows that for any $\psi\in(0,1)$ and for all $n$ large enough,
\begin{equation}\label{eq:boundexp2}
    \E(NC(M))\;\leq\;e^{-\psi A_{T^M}^*}.
\end{equation}

To conclude the second part of the proof we lower-bound $\psi A^*_{T^M}=\sum_{j=1}^b|M_j \setminus T^M_j|\psi\gamma\kappa_j + \sum_{j=1}^b |T^M_j \setminus M_j|\psi\left(\tfrac{1-\gamma}{6} n \rho(\bs X) \min_{i\in S^{L}_j(\kappa)}{\beta^*_i}^{2} - \gamma\kappa_j\right)$. To do this, we obtain a lower bound for $\psi \gamma \kappa_j$ and for $\psi\bigg(\tfrac{1-\gamma}{6} n \rho(\bs X) \min_{i\in S^{L}_j(\kappa)}{\beta^*_i}^{2} - \gamma\kappa_j\bigg)$.


The definition of $S^{L}_j(\kappa)$ implies there exists some $g'_j \to \infty$ such that
\begin{equation}\label{eq:bngrws}
  \frac{(1-\gamma) n\rho(\bs X)}{6}\min_{i\in S^{L}_j(\kappa)}{\beta^*_i}^{2}\,- \,\kappa_j \;=\; \ln(s_j) + g'_{j}.  
\end{equation}
Let $\delta \in (0,1)$ and denote $\bar{m}_j=\max\big\{\frac{2\ln(p_j-s_j)}{f_j},\frac{2\ln(s_j)}{g'_{j}}\big\}$, where $f_j$ is given in Assumption~\hyperlink{lab:A6}{A6}. 
Take $\psi=\max_{j=1,\ldots,b}\frac{\xi + \delta + \bar{m}_j}{1+\bar{m}_j}$ for some $\xi \in(0, 1- \delta)$ then $\psi \in (0,1)$ and we have, for every $j=1,\ldots,b$,
\begin{align}
    &\psi > \frac{\delta + \frac{2\ln(p_j-s_j)}{f_j}}{1+\frac{2\ln(p_j-s_j)}{f_j}} = \frac{\delta f_j/2 + \ln(p_j-s_j)}{f_j/2 +\ln(p_j-s_j)} \label{eq:psif2}\\
    &\psi > \frac{\delta + \frac{2\ln(s_j)}{g'_j}}{1+\frac{2\ln(s_j)}{g'_j}} = \frac{\delta g'_j/2 + \ln(s_j)}{g'_j/2 +\ln(s_j)} \geq \frac{\delta g'_j/2 + \ln(s_j)}{g'_j +\ln(s_j)}.\label{eq:psig2}
\end{align}
 Recall that Assumptions~\hyperlink{lab:A6}{A6}-\hyperlink{lab:A7}{A7} define $f_j= \kappa_j - \ln(p_j-s_j)$ and $\gamma= \frac{1}{2}(1 + \max_j \frac{\ln(p_j-s_j)}{\kappa_j})$ respectively. Hence,
%By definition of $\gamma= \frac{1}{2}(1 + \max_j \frac{\ln(p_j-s_j)}{\kappa_j})$ in \eqref{eq:defSes} we have
$$
\gamma\kappa_j\;\geq\;\frac12 \Big(1+\frac{\ln(p_j-s_j)}{\kappa_j}\Big)\kappa_j\;=\;\ln(p_j-s_j)+\frac12(\kappa_j-\ln(p_j-s_j)) \;=\; \ln(p_j-s_j)+ \frac12f_j.
$$
Hence, by \eqref{eq:psif2}, we have
\begin{equation}\label{eq:ooofjfj3}
\psi\gamma\kappa_j \;\geq\; \psi\big(\ln(p_j-s_j)+ \frac12f_j\big) \;\geq\;  \ln(p_j-s_j)+ \delta\frac12f_j.
\end{equation}
Further,
%By \hyperlink{lab:A7}{A7} and by our choice of $\psi$ and $\delta$,
\begin{equation}\label{eq:ooofjfj4}
\psi\bigg(\tfrac{1-\gamma}{6} n \rho(\bs X) \min_{i\in S^{L}_j(\kappa)}{\beta^*_i}^{2} - \gamma\kappa_j\bigg) \;\geq\; \psi\left(\ln(s_j)+g_j'\right) \;\geq\; \ln(s_j)+\delta \frac{1}{2}g'_j.
\end{equation}
where the first inequality follows from \eqref{eq:bngrws} and the second inequality from \eqref{eq:psig2}. In \eqref{eq:boundexp2}, $\psi A^*_{T^M}=\sum_{j=1}^b|M_j \setminus T^M_j|\psi\gamma\kappa_j + \sum_{j=1}^b |T^M_j \setminus M_j|\psi\left(\tfrac{1-\gamma}{6} n \rho(\bs X) \min_{i\in S^{L}_j(\kappa)}{\beta^*_i}^{2} - \gamma\kappa_j\right)$. Then by \eqref{eq:ooofjfj3} and \eqref{eq:ooofjfj4}, we get
\begin{equation}\label{eq:fyv2}
\E(NC(M))\;\leq\;\exp\left\{- \sum_{j=1}^b |M_j\setminus T^M_j|(\ln(p_j-s_j)+\delta \tfrac{f_j}{2})-\sum_{j=1}^b |T^M_j\setminus M_j|(\ln(s_j)+\delta \tfrac{g'_j}{2})\right\}.
\end{equation}

For the final step of the proof, denote $\mathcal{S}= \sum_{M \in \M\setminus \T(\kappa)} \E\left(NC(M)\right)$ for convenience. By \eqref{eq:fyv2} we have
\begin{align*}
    \mathcal{S} &\leq \sum_{M \in \M\setminus \T(\kappa)}e^{- \sum_{j=1}^b |M_j\setminus T^M_j|\big(\ln(p_j-s_j)+\delta \tfrac{f_j}{2}\big)-\sum_{j=1}^b |T^M_j\setminus M_j|\big(\ln(s_j)+\delta \tfrac{g'_j}{2}\big)}.
\end{align*}
We split the sum in the right-hand side above into sums over models $M$ %$M\in \M\setminus \T(\kappa)$ 
such that $T^M=T$ for some common $T\in \T(\kappa)$. Denote for any $T\in \T$,\, $\M(T):=\{M \in \M\setminus \T(\kappa) \,| T^M=T\}$, then
\begin{align}\label{eq:gepihw}
    \mathcal{S} &\leq \sum_{T \in \T(\kappa)}\sum_{M \in \M(T)}e^{- \sum_{j=1}^b |M_j\setminus T_j|\big(\ln(p_j-s_j)+\delta \tfrac{f_j}{2}\big)-\sum_{j=1}^b |T_j\setminus M_j|\big(\ln(s_j)+\delta \tfrac{g'_j}{2}\big)}.
\end{align}
The right hand-side of \eqref{eq:gepihw} is composed of a double sum over $T \in \T(\kappa)$ and over $M \in \M(T)$. Consider the sum over $M \in \M(T)$, add $T$ to it, and denote it
\begin{equation}\label{eq:defST}
  \mathcal{S}(T)=\sum_{M \in \M(T) \cup T}e^{- \sum_{j=1}^b |M_j\setminus T_j|\big(\ln(p_j-s_j)+\delta \tfrac{f_j}{2}\big)-\sum_{j=1}^b |T_j\setminus M_j|\big(\ln(s_j)+\delta \tfrac{g'_j}{2}\big)}.
\end{equation}
In the summand in the right-hand side of \eqref{eq:defST}, the case $M=T$ correspond to $|M_j\setminus T^M_j|=|T^M_j\setminus M_j|=0$ for all $j$ and the summand is then 1.
By \eqref{eq:gepihw}, we then get that
\begin{align}\label{eq:gepihw2}
    \mathcal{S} &\leq \sum_{T \in \T(\kappa)}\big(\mathcal{S}(T)-1\big).
\end{align}
For each $T=T^M$, we further split $\mathcal{S}(T)$ into sums over subsets of models $M$ that have $u_j$ more parameters than $T^M$ in block $j$, and are missing $w_j$ parameters from $T^M$.
Specifically, consider models $M$ such that, for all $j$, $|M_j\setminus T^M_j|=u_j$ and $|T^M_j\setminus M_j|=w_j$ with $u_j \in \{0,\ldots,p_j-s_j,\ldots,p_j-s_j+|S_j^{S}(\kappa)|\}$ and $w_j \in \{0,\ldots,|S^{L}_j(\kappa)|\}$. Denote by
$$S^{\bs u}_{\bs w}(T)=\sum_{M \in \M(T) \cup T:\forall \,j |M_j\setminus S_j|=u_j, |S_j\setminus M_j|=w_j } e^{- \sum_{j=1}^b u_j\big(\ln(p_j-s_j)+\delta \tfrac{f_j}{2}\big)-\sum_{j=1}^b w_j\big(\ln(s_j)+\delta \tfrac{g'_j}{2}\big)}.$$
We get
\begin{equation}\label{eq:uygfew2}
    \mathcal{S}(T)=\sum_{w_1=0}^{|S^L_1(\kappa)|}\cdots\sum_{w_b=0}^{|S^L_b(\kappa)|}\sum_{u_1=0}^{p_1-s_1+|S_1^S|}\cdots\sum_{u_b=0}^{p_b-s_b+|S_b^S|}S^{\bs u}_{\bs w}(T).
\end{equation}
The number of models missing, for all $j$, $w_j$ out of the $|S^L_j(\kappa)|$ large active parameters and having $u_j$ inactive or small active parameters from $B_j$ is $\prod_{j=1}^b {\binom{p_j-s_j+|S^S_j(\kappa)|}{u_j}}{\binom{|S^L_j(\kappa)|}{w_j}}$. We thus have
\begin{align*}
    S^{\bs u}_{\bs w}(T)&= \bigg(\prod_{j=1}^b {\binom{p_j-s_j+|S^S_j(\kappa)|}{u_j}}{\binom{|S^L_j(\kappa)|}{w_j}}\bigg)e^{- \sum_{j=1}^b u_j\big(\ln(p_j-s_j)+\delta \tfrac{f_j}{2}\big)-\sum_{j=1}^b w_j\big(\ln(s_j)+\delta \tfrac{g'_j}{2}\big)} \\
    &= \prod_{j=1}^b {\binom{p_j-s_j+|S^S_j(\kappa)|}{u_j}}e^{- u_j\big(\ln(p_j-s_j)+\delta \tfrac{f_j}{2}\big)}{\binom{|S^L_j(\kappa)|}{w_j}}e^{-w_j\big(\ln(s_j)+\delta \tfrac{g'_j}{2}\big)}.
\end{align*}
Inputting the expression above in \eqref{eq:uygfew2} and factorizing over terms in $u_j$ and $w_j$ gives
\begin{align*}
    \mathcal{S}(T)\leq\prod_{j=1}^b & \Bigg(\sum_{u_j=0}^{p_j-s_j+|S^S_j(\kappa)|}{\binom{p_j-s_j+|S^S_j(\kappa)|}{u_j}}e^{- u_j(\ln(p_j-s_j)+\delta \tfrac{f_j}{2})}\Bigg)\\&.\Bigg(
    \sum_{w_j=0}^{|S^L_j(\kappa)|}{\binom{|S^L_j(\kappa)|}{w_j}}
    e^{-w_j(\ln(s_j)+\delta \tfrac{g'_j}{2})}\Bigg).
\end{align*}
By the bound in \eqref{eq:upperbound_sumprob} and taking the terms in $u_j=0$ and $w_j=0$ out of the sums above, %taking the term in 0 out of the sum above, 
we have
\begin{align}\label{eq:ogr}
    \mathcal{S}(T)\leq\prod_{j=1}^b & \left(1+\sum_{u_j=1}^{p_j-s_j+|S_j^{S}(\kappa)|}e^{-u_j\left(\delta \tfrac{f_j}{2}-\ln\big(1+\frac{|S^{S}_j(\kappa)|}{p_j-s_j}\big)-1\right)}\right)\\&.\left(1+
    \sum_{w_j=1}^{|S_j^{L}(\kappa)|}    e^{-w_j\left(\delta \tfrac{g'_j}{2}+\ln\big(\frac{s_j}{|S_j^{L}(\kappa)|}\big)-1\right)}\right).\nonumber
\end{align}
Denote
$$
d_j\;=\;e^{1+\ln\big(1+\frac{|S^{S}_j(\kappa)|}{p_j-s_j}\big)-\delta \tfrac{f_j}{2}},\qquad h_j\;=\;e^{1-\ln\big(\frac{s_j}{|S_j^{L}(\kappa)|}\big)-\delta \tfrac{g'_j}{2}}.
$$
By assumption $|S^{S}_j(\kappa)|=O(p_j-s_j)$, and by definition of $|S_j^{L}(\kappa)|$ we have $s_j\geq|S_j^{L}(\kappa)|$. By Assumptions~\hyperlink{lab:A6}{A6}-\hyperlink{lab:A7}{A7}, we also have $f_j\to\infty$ and $g'_j\to\infty$, and then $\lim_{n \to \infty}= d_j= \lim_{n \to \infty} h_j=0$. Using the properties of geometric series, 
 for every $j$ we have
\begin{eqnarray*}
    &1+\sum_{u_j=1}^{p_j-s_j+|S_j^{S}(\kappa)|}e^{-u_j\left(\delta \tfrac{f_j}{2}-\ln\big(1+\frac{|S^{S}_j(\kappa)|}{p_j-s_j}\big)-1\right)}
    =\frac{1-d_j^{p_j-s_j+|S_j^{S}(\kappa)|+1}}{1-d_j} \\
    &1+\sum_{w_j=1}^{|S_j^{L}(\kappa)|}    e^{-w_j\left(\delta \tfrac{g'_j}{2}+\ln\big(\frac{s_j}{|S_j^{L}(\kappa)|}\big)-1\right)} =\frac{1-h_j^{|S_j^{L}(\kappa)|+1}}{1-h_j},
\end{eqnarray*}
where both expressions converge to $1$ as $n$ grows. By \eqref{eq:gepihw2} and \eqref{eq:ogr}:
$$\mathcal{S}\leq\sum_{T\in\T}\Bigg(\prod_{j=1}^b\Big(\frac{1-d_j^{p_j-s_j+|S_j^{S}(\kappa)|+1}}{1-d_j}\Big)\Big(\frac{1-h_j^{|S_j^{L}(\kappa)|+1}}{1-h_j}\Big)-1\Bigg).$$
Each of the summand vanishes as $n\to\infty$. Moreover, by assumption $|S^I(\kappa)|=O(1)$ and then $|\T|=2^{|S^I(\kappa)|}=O(1)$. We thus have $\lim_{n\to\infty}\mathcal{S}=\lim_{n\to\infty}\sum_{M \in \M\setminus \T(\kappa)} \E\left(NC(M)\right)=0$.

Further, by Lemma~\ref{lem:L0toL1convergencegen}, 
$P(\hat{S}^{b} \not\in \T(\kappa) ) \,\,\leq \,\, (|\T(\kappa)|+1) \mathcal{S} = (2^{|S^{I}(\kappa)|}+1) \mathcal{S}$. Since $\mathcal{S}$ vanishes and $|S^{I}(\kappa)|=O(1)$, $\lim_{n\to\infty}P(\hat{S}^{b} \in \T(\kappa) )=1$ as we wished to prove.


\section{Proofs of Section~\ref{sec:informed_l0}}\label{suppsec:proofsempbayes}

\subsection{Proof of Proposition~\ref{prop:shatconsistence}}
Denote
$$A_j:=\frac{1}{p_j}\sum_{i \in B_j}\sum_{M \in \T(\kappa)|i\in M}NC(M) \quad\text{and}\quad  C_j:=\frac{1}{p_j}\sum_{i \in B_j}\sum_{M \in \M\setminus \T(\kappa)|i\in M}NC(M).$$
For every $j=1,\ldots,b$, we have the decomposition
\begin{gather}\label{eq:decompshat}
    \frac{\hat{s}_j}{p_j}
    =\frac{\sum_{i \in B_j} \sum_{M \in \M | i \in M} NC(M)}{p_j}
   =A_j\,+\, C_j.
\end{gather}
To show the lower bound on $\hat{s}_j/p_j$, we decompose $A_j$
\begin{eqnarray}
  A_j
   &=& \sum_{M \in \T(\kappa)}NC(M)\sum_{i \in B_j}\frac{I(i\in M_j)}{p_j}\nonumber\\
   &=& \sum_{M \in \T(\kappa)}NC(M)\sum_{i \in S^{L}_j(\kappa)}\frac{I(i\in M_j)}{p_j}+\sum_{M \in \T(\kappa)}NC(M)\sum_{i \in B_j\setminus S^{L}_j(\kappa)}\frac{I(i\in M_j)}{p_j}\nonumber\\
   &=& \frac{|S^{L}_j(\kappa)|}{p_j}\sum_{M \in \T(\kappa)}NC(M)+\sum_{M \in \T(\kappa)}NC(M)\sum_{i \in B_j\setminus S^{L}_j(\kappa)}\frac{I(i\in M_j)}{p_j}\label{eq:wigr}.
\end{eqnarray}
where the last equality follows from $I(i\in M_j)=1$ for all $i\in S^L(\kappa)$ when $M\in\T(\kappa)$. The righmost term above and $C_j$ are nonnegative, then by the linearity of the expectation
$$\E\Big(\frac{\widehat{s_j}}{p_j}\Big)\geq\frac{|S^{L}_j(\kappa)|}{p_j}\sum_{M \in \T(\kappa)}\E(NC(M)).$$
By Theorem~\ref{thm:convtoT}, $\lim_{n\to\infty}\sum_{M \in \T(\kappa)}\E(NC(M))=1$. It follows that $\lim_{n\to\infty}\E\Big(\frac{\widehat{s_j}}{p_j}\Big)\geq \frac{|S^{L}_j(\kappa)|}{p_j}$ for every $j=1,\ldots,b$.

We now prove the upper bound. Recall that $\T(\kappa)$ by definition includes models that have no small signals, i.e. all parameters are in $S^L(\kappa) \cup S^I(\kappa)$. That is,
%Since 
for all $M\in\T(\kappa)$, we have that $I(i\in M_j)=0$ for all $i\in B_j\setminus(S^{L}_j(\kappa)\cup S^{I}(\kappa)_j)$. Hence, $A_j$ in \eqref{eq:wigr} satisfies
\begin{align*}
    A_j
   &= \frac{|S^{L}_j(\kappa)|}{p_j}\sum_{M \in \T(\kappa)}NC(M)+\sum_{M \in \T(\kappa)}NC(M)\sum_{i \in S^{I}_j(\kappa)}\frac{I(i\in M_j)}{p_j}\\
   &\leq\frac{|S^{L}_j(\kappa)|}{p_j}\sum_{M \in \T(\kappa)}NC(M)+ \frac{|S^{I}(\kappa)_j|}{p_j}\sum_{M \in \T(\kappa)}NC(M)
\end{align*}
where the inequality follows from $\sum_{i \in S^{I}_j(\kappa)} I(i\in M_j) \leq |S^{I}(\kappa)_j|$  for all $M$. By \eqref{eq:decompshat}, we then have
\begin{equation}\label{eq:uuuuu}
    \frac{\hat{s}_j}{p_j}\leq \frac{|S^{L}_j(\kappa)|+|S^{I}(\kappa)_j|}{p_j}\sum_{M \in \T(\kappa)}NC(M) + C_j.
\end{equation}
Moreover, for every $j=1,\ldots,b$, $C_j$ satisfies
\begin{equation}\label{eq:uuuuu2}
    C_j = \sum_{M \in \M\setminus \T(\kappa)}NC(M)\sum_{i \in B_j}\frac{I(i\in M_j)}{p_j}\leq \sum_{M \in \M\setminus \T(\kappa)}NC(M)
\end{equation}
where the inequality follows from $\sum_{i \in B_j}\frac{I(i\in M_j)}{p_j}\leq 1$ for all $M$.
Taking expectations in \eqref{eq:uuuuu} and \eqref{eq:uuuuu2} gives
\begin{equation*}
    \E\bigg(\frac{\widehat{s_j}}{p_j}\bigg) 
    \leq \frac{|S^{L}_j(\kappa)|+|S^{I}(\kappa)_j|}{p_j}\sum_{M \in \T(\kappa)}\E(NC(M)) + \sum_{M \in \M\setminus \T(\kappa)}\E(NC(M)).
\end{equation*}
 By Theorem~\ref{thm:convtoT}, we have on one hand $\lim_{n\to \infty}\sum_{M \in \T(\kappa)}\E(NC(M))= 1$ and, on the other hand, $\lim_{n\to \infty}\sum_{M \in \M\setminus\T(\kappa)}\E(NC(M))= 0$. It follows that $ \lim_{n\to \infty}\E\big(\frac{\widehat{s_j}}{p_j}\big) \leq  \frac{|S^{L}_j(\kappa)|+|S^{I}(\kappa)_j|}{p_j}= \frac{s_j-|S^{S}_j(\kappa)|}{p_j}$ for every $j=1,\ldots,b$, which proves the upper bound.


\subsection{Proof of Theorem~\ref{theo:empbayesselconsist}} The proof strategy is to show that Assumptions~\hyperlink{lab:A1}{A1}, \hyperlink{lab:A6}{A6} and \hyperlink{lab:A7}{A7} hold to apply Theorem~\ref{theo:suffcondlinearmodel}. Recall that Theorem~\ref{theo:empbayesselconsist} makes Assumptions~\hyperlink{lab:A1}{A1}, \hyperlink{lab:A8}{A8} and \hyperlink{lab:A9}{A9}, hence it suffices to show that \hyperlink{lab:A6}{A6}-\hyperlink{lab:A7}{A7} hold.
We first derive a convenient decomposition of the empirical Bayes penalties $\kappa^{EB}_j$. The second step of the proof is to show that, with probability going to 1, these $\kappa^{EB}_j$ satisfy Assumption~\hyperlink{lab:A6}{A6}. The third step consists in showing that Assumption~\hyperlink{lab:A9}{A9} implies Assumption~\hyperlink{lab:A7}{A7} for $\kappa^{EB}_j$. The consistency of $\hat S^{EB,b}$ then follows from Theorem~\ref{theo:suffcondlinearmodel}.

Denote for any $M\in\M$, $NC^{\circ}(M)$, the normalized criterion value for model $M$ under Step 1 penalty $\kappa^{\circ}$. For this choice of penalty and every $j=1,\ldots,b$, we have
\begin{align*}
    \frac{\widehat{s_j}}{p_j}
    &\;=\; \frac{1}{p_j}\sum_{i\in B_j}\,\,\sum_{M\in \M}NC(M)I(i\in M)\\
   &\;=\; \frac{1}{p_j}\sum_{M \in \M}NC^{\circ}(M)\sum_{i \in B_j}I(i\in M) \\
   &\;=\;\frac{s_j}{p_j}NC^{\circ}(S) + \sum_{M \in \M | M \neq S}\frac{|M_j|}{p_j}NC^{\circ}(M).
\end{align*}
Using that $NC^{\circ}(S)=1-\sum_{M \in \M | M \neq S}NC^{\circ}(M)$, we get
\begin{equation*}
   \frac{\hat{s}_j}{p_j}
   =\frac{s_j}{p_j} + \sum_{M \in \M | M \neq S}\frac{|M_j|-s_j}{p_j}NC^{\circ}(M).
\end{equation*}
Consider the decomposition of the sum in the right-hand side above between the sum over models $M$ that contain more parameters than $S$ in block $j$ and the sum over those that contain less parameters than $S$. Denote
$$
    O_j^{\circ} := \sum_{M \in \M | |M_j| > s_j}\frac{|M_j|-s_j}{p_j}NC^{\circ}(M)\quad\text{and}\quad
    U_j^{\circ} := \sum_{M \in \M | |M_j| < s_j}\frac{s_j-|M_j|}{p_j}NC^{\circ}(M).
$$
%are overfitting in $B_j$ in the sense $|M_j|> s_j$ and those that are underfitting in $B_j$, in the sense $|M_j|<s_j$, we can write:
We have
\begin{equation}\label{eq:rewriteshat}
  \frac{\hat{s}_j}{p_j}  = \frac{s_j}{p_j}\,\,+\,\, O_j^{\circ}-\,\,U_j^{\circ}.
\end{equation}
Observe that we have the following decomposition of Step 2 penalties
$$\kappa^{EB}_j
    = \ln(p_j-s_j) + \ln\Big(\frac{\sqrt{n}}{s_j}\Big) + \ln\Big(\frac{p_j-\hat{s}_j}{p_j-s_j}\Big) + \ln\Big(\frac{s_j}{\hat{s}_j}\Big). \nonumber$$
By \eqref{eq:rewriteshat}, it follows that
\begin{eqnarray}\label{eq:decompstep2pen}
    \kappa^{EB}_j&=& \ln(p_j-s_j)+\ln\Big(\frac{\sqrt{n}}{s_j}\Big)+\ln\Big(1-\frac{p_j(O^{\circ}_j-U^{\circ}_j)}{p_j-s_j}\Big)+\ln\Big(\frac{s_j}{\hat{s}_j}\Big),
\end{eqnarray}
completing the first step of the proof.

We continue with the second step of the proof: showing that the $\kappa^{EB}_j$'s satisfy Assumption~\hyperlink{lab:A6}{A6} with probability going to 1. Recall that Assumption~\hyperlink{lab:A6}{A6} states that there exists $f_j\to \infty$ (as $n\to \infty$) such that for every sufficiently large $n$,
\begin{equation*}
    \kappa_j \;=\;\ln(p_j-s_j) + f_j.
\end{equation*}
Since $U_j^{\circ}$ is nonnegative, a lower bound on $\kappa^{EB}_j$ is %for every $j=1,\ldots,b$ is:
\begin{equation}\label{eq:decompstep2pen2}
    \kappa^{EB}_j
    \geq  \ln(p_j-s_j)+\ln\Big(\frac{\sqrt{n}}{s_j}\Big)+\ln\Big(1-\frac{p_j O^{\circ}_j}{p_j-s_j}\Big)+\ln\Big(\frac{s_j}{\hat{s}_j}\Big).
\end{equation}
Plugging in the definition of $O^{\circ}_j$, we have that
$$\frac{p_j O^{\circ}_j}{p_j-s_j} =  \sum_{M \in \M | |M_j| > s_j}\frac{|M_j|-s_j}{p_j-s_j}NC^{\circ}(M)\leq \sum_{M \in \M | |M_j| > s_j}NC^{\circ}(M)$$
where the inequality follows from $(|M_j|-s_j)/(p_j-s_j)\leq 1$ for all $M$.
Note that if $M$ is such that $|M_j|>s_j$, then $M\not\in\T(\kappa^{\circ})$ (this follows immediately from the definition of $\T(\kappa)$ in \eqref{eq:Tkappa})  and therefore $\sum_{M \in \M | |M_j| > s_j}NC^{\circ}(M)\leq \sum_{M \in \M \setminus \T(\kappa^{\circ})}NC^{\circ}(M)$. Moreover, $\kappa^{\circ}$ satisfies Assumption~\hyperlink{lab:A6}{A6} and the assumptions of Theorem~\ref{thm:convtoT} are met for $\kappa^{\circ}$.
Then, by Theorem~\ref{thm:convtoT}, $\lim_{n\to\infty} \sum_{M \in \M \setminus \T(\kappa^{\circ})}NC^{\circ}(M)=\lim_{n\to\infty}\sum_{M \in \M | |M_j| > s_j}NC^{\circ}(M) = 0$, $\frac{p_j O^{\circ}_j}{p_j-s_j}$ vanishes in probability and so does $\ln\Big(1-\frac{p_j O^{\circ}_j}{p_j-s_j}\Big)$. By Assumption~\hyperlink{lab:A8}{A8}, we also have that $\ln(\sqrt{n}s_j^{-1})\to\infty$. 
Then, to show that Assumption~\hyperlink{lab:A6}{A6} holds it is enough to show that with probability going to 1, $\ln(s_j/\hat{s}_j)$ is nonnegative. Observe that all assumptions in Proposition~\ref{prop:shatconsistence} are also met for $\kappa^{\circ}$. By \eqref{eq:uuuuu} and \eqref{eq:uuuuu2} in the proof of Proposition~\ref{prop:shatconsistence}, we have
\begin{equation*}
    \frac{\hat{s}_j}{p_j}\leq \frac{|S^{L}_j(\kappa^{\circ})|+|S^{I}(\kappa^{\circ})_j|}{p_j}\sum_{M \in \T(\kappa^{\circ})}NC(M) + \frac{1}{p_j} \sum_{M \in \M\setminus \T(\kappa^{\circ})}NC(M).
\end{equation*}
By Theorem~\ref{thm:convtoT}, $\sum_{M \in \T(\kappa^{\circ})}NC(M)$ and $\sum_{M \in \M\setminus\T(\kappa^{\circ})}NC(M)$ converge in probability to 1 and 0 respectively. We then have that, with probability going to 1,
\begin{equation*}
    \frac{\hat{s}_j}{p_j}\leq \frac{|S^{L}_j(\kappa^{\circ})|+|S^{I}(\kappa^{\circ})_j|}{p_j} \Longrightarrow
    \quad\ln\Big(\frac{s_j}{\hat{s}_j}\Big)\geq \ln\bigg(\frac{s_j}{|S^{L}_j(\kappa^{\circ})|+|S^{I}(\kappa^{\circ})_j|} \bigg) \geq 0.
\end{equation*}
We then obtain that, with probability going to 1,
\begin{equation}\label{eq:eeeeeeeh}
    \kappa^{EB}_j
    \geq  \ln(p_j-s_j)+\ln\Big(\frac{\sqrt{n}}{s_j}\Big)
\end{equation}
and that the $\kappa^{EB}_j$ satisfy Assumption~\hyperlink{lab:A6}{A6}, completing the second part of the proof.\medskip

For the third and final part of the proof, 
we now show that Assumption~\hyperlink{lab:A9}{A9} implies Assumption~\hyperlink{lab:A7}{A7} for the $\kappa_j^{EB}$ with probability going to 1. Recall that Assumption~\hyperlink{lab:A7}{A7} for the $\kappa_j^{EB}$ states that for each block $j$ there exists $g_j\to \infty$ such that for large enough $n$,
    \begin{equation*}
        \sqrt{\frac{(1-\gamma) n\rho(\bs X)}{6}}{{\beta_{\min,j}^*}} \,- \,\sqrt{\kappa^{EB}_j} \;=\; \sqrt{\ln(s_j)} + g_j .
    \end{equation*}
where $\gamma$ takes value
\begin{equation}\label{eq:gammakeb}
    \gamma=\frac12\Big(1+\max_j\frac{\ln(p_j-s_j)}{\kappa^{EB}_j}\Big).
\end{equation}
Observe that Assumption~\hyperlink{lab:A9}{A9} and Assumption~\hyperlink{lab:A7}{A7} take the same form. To show that Assumption~\hyperlink{lab:A9}{A9} implies Assumption~\hyperlink{lab:A7}{A7} for the $\kappa_j^{EB}$ with probability going to 1, it suffices to show that the following two inequalities
\begin{align}
    \sqrt{\frac{ (1-\gamma)n\rho(\bs X)}{6}}{\beta_{\min,j}^*} \,&\geq \,  \sqrt{\frac{ (1-\psi)n\rho(\bs X)}{6}}{\beta_{\min,j}^*}, \label{eq:oaebr} \\
    - \,\sqrt{\kappa^{EB}_j} \,&\geq \,-\sqrt{\ln\bigg(\frac{p}{|S^{L}(\kappa^{\circ})|}-1\bigg)+\frac12\ln(n)}\label{eq:oaebr2}
\end{align}
hold with probability going to 1 for $\gamma$ as in \eqref{eq:gammakeb} and $\psi=\frac12\big(1+\max_j \frac{\ln(p_j-s_j)}{\ln(p_j/s_j-1)+0.5\ln(n)}\big)$ (defined in Assumption~\hyperlink{lab:A9}{A9}). We first show \eqref{eq:oaebr} holds with probability going to 1 and then that \eqref{eq:oaebr2} does too.

By \eqref{eq:eeeeeeeh}, with probability going to 1,
$$ \frac{\ln(p_j-s_j)}{\kappa^{EB}_j}
    \leq \frac{\ln(p_j-s_j)}{\ln(p_j-s_j)+\ln\Big(\frac{\sqrt{n}}{s_j}\Big)} =  \frac{\ln(p_j-s_j)}{\ln(p_j/s_j-1)+0.5\ln(n)} .$$
It follows that:
\begin{equation*}
    \gamma=\frac12\Big(1+\max_{j=1, \ldots,b}\frac{\ln(p_j-s_j)}{\kappa^{EB}_j}\Big) \leq \frac12\Big(1+\max_j \frac{\ln(p_j-s_j)}{\ln(p_j/s_j-1)+0.5\ln(n)}\Big)=\psi
\end{equation*}
and \eqref{eq:oaebr} holds with probability going to 1.

We now upper bound $\kappa^{EB}_j$ to show \eqref{eq:oaebr2} holds with probability going to 1. Observe that 
$$
\ln\Big(\frac{\hat{s}_j}{s_j}\Big)=
 \ln \left( 1 + \frac{\hat{s}_j - s_j}{s_j}\right) =   
\ln\Big(1+\frac{p_j(O^{\circ}_j-U^{\circ}_j)}{s_j}\Big).
$$
where the second equality follows from \eqref{eq:rewriteshat}.
Plugging this expression into \eqref{eq:decompstep2pen}, and using that $O_j^{\circ} \geq 0$, we have that  %Since $O_j^{\circ}$ is nonnegative, from \eqref{eq:decompstep2pen} we have that %, for every $j=1,\ldots,b$, an upper bound on $\kappa^{EB}_j$ is:
\begin{equation}\label{eq:gggggggg}
    \kappa^{EB}_j \leq \ln\big(p_j-s_j\big)+\ln\Big(\frac{\sqrt{n}}{s_j}\Big)+\ln\bigg(\frac{1+\frac{p_j}{p_j-s_j}U^{\circ}_j}{1-\frac{p_j}{s_j}U^{\circ}_j}\bigg).
\end{equation}
We split the sum in $U^{\circ}_j$ between models in $\T(\kappa^{\circ})$ and those not in $\T(\kappa^{\circ})$. 
$$U_j^{\circ} = \sum_{M \in \T(\kappa^{\circ}) | |M_j| < s_j}\frac{s_j-|M_j|}{p_j}NC^{\circ}(M)+\sum_{M \in \M\setminus\T(\kappa^{\circ}) | |M_j| < s_j}\frac{s_j-|M_j|}{p_j}NC^{\circ}(M).
$$
If $M\in \T(\kappa^{\circ})$, then by definition $|M_j|\geq |S^{L}_j(\kappa^{\circ})|$ and thus $s_j-|M_j|\leq s_j-|S^{L}_j(\kappa^{\circ})|$. A bound on $s_j-|M_j|$ for $M\not\in \T(\kappa^{\circ})$ is simply $s_j-|M_j|\leq s_j$. It follows that
\begin{eqnarray*}
    U^{\circ}_j \leq \frac{s_j-|S^{L}_j(\kappa^{\circ})|}{p_j}\sum_{M \in \T(\kappa^{\circ}) | |M_j| < s_j}NC^{\circ}(M)+  \frac{s_j}{p_j}\sum_{M \in \M \setminus \T(\kappa^{\circ})| |M_j| < s_j}NC^{\circ}(M)
\end{eqnarray*}
By Theorem~\ref{thm:convtoT}, $\sum_{M \in  \T(\kappa^{\circ})| |M_j| < s_j}NC^{\circ}(M)$ and $\sum_{M \in \M \setminus \T(\kappa^{\circ})| |M_j| < s_j}NC^{\circ}(M)$ converge in probability to 1 and 0 respectively. We then get that, with probability going to 1,
\begin{align}
    \frac{p_j}{p_j-s_j}U^{\circ}_j &\leq \frac{s_j-|S^{L}_j(\kappa^{\circ})|}{p_j-s_j}\nonumber\\
    \frac{p_j}{s_j}U^{\circ}_j &\leq  \frac{s_j-|S^{L}_j(\kappa^{\circ})|}{s_j}. \label{eq:gewou}
\end{align}
By the bounds above and \eqref{eq:gggggggg}, we have that with probability going to 1,
\begin{align*}
  \kappa^{EB}_j 
  &\leq \ln\big(p_j-s_j\big)+\ln\Big(\frac{\sqrt{n}}{s_j}\Big)+\ln\Bigg(\frac{1+\frac{s_j-|S^{L}_j(\kappa^{\circ})|}{p_j-s_j}}{1- \frac{s_j-|S^{L}_j(\kappa^{\circ})|}{s_j}}\Bigg) \nonumber\\
  &= \ln\big(p_j-s_j\big) + \ln\Big(\frac{\sqrt{n}}{s_j}\Big) + \ln\Bigg(\frac{\frac{p_j-|S^{L}_j(\kappa^{\circ})|}{p_j-s_j}}{ \frac{|S^{L}_j(\kappa^{\circ})|}{s_j}}\Bigg) \\
&=\ln\big(p_j/|S^{L}_j(\kappa^{\circ})|-1\big)+\frac{1}{2}\ln(n)\label{eq:finalupperboundkappaj}
\end{align*}
which shows \eqref{eq:oaebr2} holds with probability going to 1 and that
Assumption~\hyperlink{lab:A9}{A9} implies Assumption~\hyperlink{lab:A7}{A7} holds for the $\kappa^{EB}_j$ with probability going to 1. 

Since Assumptions~\hyperlink{lab:A6}{A6} and \hyperlink{lab:A7}{A7} hold with probability going to 1, by Theorem~\ref{theo:suffcondlinearmodel}, $\lim_{n\to\infty}P(\hat{S}^{EB,b}= S)= 1$, as we wished to prove.

\subsection{Proof of Theorem~\ref{theo:informedl0consist}} The proof strategy is similar to that of Theorem~\ref{theo:empbayesselconsist} and relies on several results therein. The first step is to show that $\kappa^{A}_j$ satisfies Assumption~\hyperlink{lab:A6}{A6} with probability going to 1 as $n$ grows. The second step is to show that Assumption~\hyperlink{lab:A10}{A10} implies Assumption~\hyperlink{lab:A7}{A7} for the $\kappa^{A}_j$ with probability going to 1. The consistency of $\hat S^{A,b}$ then follows from Theorem~\ref{theo:suffcondlinearmodel}.

Observe that $\kappa_j^{A}=\kappa_j^{EB}+\ln(\widehat s_j)$, hence by \eqref{eq:decompstep2pen2} we have that
\begin{equation*}
    \kappa_j^{A}
    \geq  \ln(p_j-s_j)+\ln(\sqrt{n})+\ln\Big(1-\frac{p_j O^{\circ}_j}{p_j-s_j}\Big).
\end{equation*}
In the proof of Theorem~\ref{theo:empbayesselconsist} we showed that, since $\kappa^{\circ}$ satisfies Assumption A6, by Theorem~\ref{thm:convtoT} $\ln\Big(1-\frac{p_j O^{\circ}_j}{p_j-s_j}\Big)$ vanishes in probability as $n$ grows. With probability going to 1, we then have that
\begin{equation}\label{eq:finallowerboundkappaj2}\kappa^{A}_j 
    \geq  \ln(p_j-s_j)+\ln(\sqrt{n})
\end{equation}
and hence that the $\kappa^{A}_j$'s satisfy Assumption~\hyperlink{lab:A6}{A6}.

For the second part of the proof,
we now show that Assumption~\hyperlink{lab:A10}{A10} implies Assumption~\hyperlink{lab:A7}{A7} for the $\kappa^{A}_j$. Assumption~\hyperlink{lab:A7}{A7} for the $\kappa_j^{A}$ states that for each block $j$ there exists $g_j\to \infty$ such that for large enough $n$,
    \begin{equation*}
        \sqrt{\frac{(1-\gamma) n\rho(\bs X)}{6}}{{\beta_{\min,j}^*}} \,- \,\sqrt{\kappa^{EB}_j} \;=\; \sqrt{\ln(s_j)} + g_j .
    \end{equation*}
where $\gamma$ takes value
\begin{equation}\label{eq:gammaka}
    \gamma=\frac12\Big(1+\max_j\frac{\ln(p_j-s_j)}{\kappa^{A}_j}\Big).
\end{equation}
To show that Assumption~\hyperlink{lab:A10}{A10} implies Assumption~\hyperlink{lab:A7}{A7} for the $\kappa_j^{A}$ with probability going to 1, it suffices to show that the following two inequalities
\begin{align}
    \sqrt{\frac{ (1-\gamma)n\rho(\bs X)}{6}}{\beta_{\min,j}^*} \,&\geq \,  \sqrt{\frac{ (1-\xi)n\rho(\bs X)}{6}}{\beta_{\min,j}^*}, \label{eq:oaebr3}\;\;\text{and} \\
    - \,\sqrt{\kappa^{A}_j} \,&\geq \,-\sqrt{\ln\bigg(p-|S^{L}(\kappa^{\circ})|\bigg)+\frac12\ln(n)}\label{eq:oaebr4}
\end{align}
hold with probability going to 1 for $\gamma$ as in \eqref{eq:gammaka} and $\xi=\frac12\big(1+\max_j \frac{\ln(p_j-s_j)}{\ln(p_j-s_j)+0.5\ln(n)}\big)$ (defined in Assumption~\hyperlink{lab:A10}{A10}). We first show \eqref{eq:oaebr3} holds with probability going to 1 and then that \eqref{eq:oaebr4} does too.

By \eqref{eq:finallowerboundkappaj2} we have that with probability going to 1, for any $j$
\begin{equation*}
\frac{\ln(p_j-s_j)}{\kappa^{A}_j} 
    \leq  \frac{\ln(p_j-s_j)}{\ln(p_j-s_j)+\ln(n)/2}.
\end{equation*}
It follows that
\begin{equation*}
    \gamma=\frac12\Big(1+\max_{j=1, \ldots,b}\frac{\ln(p_j-s_j)}{\kappa^ A_j}\Big) \leq \frac12\big(1+\max_j \frac{\ln(p_j-s_j)}{\ln(p_j-s_j)+\ln(n)/2}\big)=\xi
\end{equation*}
and \eqref{eq:oaebr3} holds with probability going to 1.

We now upper bound $\kappa^{A}_j$ to show \eqref{eq:oaebr4} holds with probability going to 1. By \eqref{eq:decompstep2pen}, we can write
    $$\kappa_j^{A}=\kappa_j^{EB}+\ln(\widehat s_j)= \ln(p_j-s_j)+\ln\big(\sqrt{n}\big)+\ln\Big(1-\frac{p_j(O^{\circ}_j-U^{\circ}_j)}{p_j-s_j}\Big).$$
Since $O_j^{\circ} \geq 0$, we obtain that % is nonnegative, for every $j=1,\ldots,b$, an upper bound on $\kappa^{A}_j$ is:
\begin{equation*}
    \kappa^{A}_j \leq \ln\big(p_j-s_j\big)+\ln(\sqrt{n})+\ln\bigg(1+\frac{p_j}{p_j-s_j}U^{\circ}_j\bigg).
\end{equation*}
By \eqref{eq:gewou}, with probability going to 1:
\begin{align}\label{eq:finalupperboundkappaj2}
    \kappa_j^A &\leq  \ln\big(p_j-s_j\big)+\ln(\sqrt{n})+\ln\bigg(1+\frac{s_j-|S^{L}_j(\kappa^{\circ})|}{p_j-s_j}\bigg)\nonumber\\
    &= \ln\big(p_j-|S^{L}_j(\kappa^{\circ})|\big)+\frac{1}{2}\ln(n).
\end{align}
which shows \eqref{eq:oaebr4} holds with probability going to 1 and that
Assumption~\hyperlink{lab:A10}{A10} implies Assumption~\hyperlink{lab:A7}{A7} holds for the $\kappa^{A}_j$ with probability going to 1. 

Since Assumptions~\hyperlink{lab:A6}{A6} and \hyperlink{lab:A7}{A7} hold with probability going to 1, by Theorem~\ref{theo:suffcondlinearmodel}, $\lim_{n\to\infty}P(\hat{S}^{A,b}= S)= 1$, as we wished to prove.

\section{Gaussian sequence model with fixed number of active signals}\label{suppsec:appendixGSMsfinite}

We derive here selection properties of the block $\ell_0$ penalties in the Gaussian sequence model dropping Assumption~\hyperlink{lab:A3}{A3} and focusing instead on regimes where the following assumption holds\\[5pt]
\hypertarget{lab:A11}{(A11)} \quad For all $j$, \;$s_j \;\leq\;k_j$ for some constant $k_j$.\\[5pt]
Changing Assumption~\hyperlink{lab:A3}{A3} for Assumption~\hyperlink{lab:A11}{A11} implies redeveloping results relative to the probability of false negatives and consequently sufficient and necessary betamin assumptions. Proposition~\ref{prop:isthresh} (on the equivalence between block penalties and thresholding in the Gaussian sequence model) as well as Proposition~\ref{prop:suffnecblockthresh} (i) and (ii) (on the probability of false positives) do not assume Assumption~\hyperlink{lab:A3}{A3}, they do not depend on results assuming \hyperlink{lab:A3}{A3}, and hold equally under Assumption~\hyperlink{lab:A11}{A11}.

\subsection{Selection based on block thresholds}

Consider the betamin assumption\\[5pt]
\hypertarget{lab:A12}{(A12)} \quad For all $j$, \; $\sqrt{n}(\beta_{\min,j}^*-\tau_j) \to \infty$.
\begin{prop}\label{prop:suffnecIIblockthreshbis}
 In the sequence model \eqref{eq:GSM}, assume \hyperlink{lab:A1}{A1}, \hyperlink{lab:A2}{A2}, \hyperlink{lab:A11}{A11}, and \hyperlink{lab:A12}{A12}.
 \begin{enumerate}[label=(\roman*)]
  \item Then $\lim_{n\to\infty}P(\hat{S}^{b}\supseteq S ) \;=\; 1$.
  \item If, in addition, Assumptions~\hyperlink{lab:A4}{A4} holds, then  $\lim_{n\to\infty}P(\hat{S}^{b}=S  ) \;=\; 1$.
 \end{enumerate}
\end{prop}
Under Assumption ~\hyperlink{lab:A11}{A11}, Assumption~\hyperlink{lab:A12}{A12} is then sufficient for $\hat{S}^{b}$ to hold the screening property (i.e., including all truly active parameters asymptotically). When, in addition, Assumption~\hyperlink{lab:A4}{A4}, that requires the block thresholds grow at least as fast as $\sqrt{2\ln(p_j-s_j)/n}$, holds, $\hat{S}^{b}$ is variable selection consistent.

By Lemma~\ref{lem:gennecIIblockthresh}, Assumption~\hyperlink{lab:A12}{A12} is also necessary for $\hat{S}^{b}$ to hold the screening property, independently of assumptions on the $s_j$'s. It follows that, under Assumption ~\hyperlink{lab:A11}{A11}, Assumption~\hyperlink{lab:A12}{A12} is  necessary and sufficient for $\hat{S}^{b}$ to hold the screening property. In Proposition~\ref{prop:suffnecblockthresh} (i) and (ii), we also showed that Assumption~\hyperlink{lab:A4}{A4} is necessary and sufficient for the vanishing of the FWER. We then get the next proposition on necessary assumptions for consistent recovery.

\begin{lem}\label{lem:necbetaminblockthresh2}
    In the sequence model \eqref{eq:GSM}, assume \hyperlink{lab:A1}{A1}, \hyperlink{lab:A2}{A2}, \hyperlink{lab:A11}{A11} and that there exists $j\in\{1,\ldots,b\}$ such that
\begin{equation}\label{eq:necbetaminblockthreshbis}
      \lim_{n\to \infty}\sqrt{n}\beta_{\min,j}^*-\sqrt{2\ln(p_j-s_j)} < \infty . 
    \end{equation}
Then $\lim_{n \to \infty} P(\hat{S}^{b} = S) < 1$.
\end{lem}
By Lemma~\ref{lem:necbetaminblockthresh2}, under Assumption~\hyperlink{lab:A11}{A11}, a necessary assumption for asymptotic support recovery is
\begin{equation}\label{eq:nectypeIIGSMfinite}
\lim_{n \to \infty} \sqrt{n}\beta_{\min,j}^*-\sqrt{2\ln(p_j-s_j)} = \infty.
\end{equation}

The earlier Theorem~\ref{theo:fullrecovGS} on rates of convergence with $\hat S^b$ holds under Assumptions~\hyperlink{lab:A4}{A4} and \hyperlink{lab:A5}{A5}, independently of Assumption~\hyperlink{lab:A3}{A3}. Observe that under Assumption~\hyperlink{lab:A11}{A11}, Assumption~\hyperlink{lab:A12}{A12} implies Assumption~\hyperlink{lab:A5}{A5} for every $n$ large enough. Then Theorem~\ref{theo:fullrecovGS} holds equally under Assumption~\hyperlink{lab:A11}{A11}, assuming  \hyperlink{lab:A4}{A4} and \hyperlink{lab:A12}{A12}.

We next shortly examine the benefits of block penalties in this setting. Assumptions~\hyperlink{lab:A4}{A4} and \hyperlink{lab:A12}{A12} give ranges of thresholds that are necessary and sufficient for asymptotic support recovery. For the standard selector $\hat S$, the single threshold $\tau$ is required to satisfy, for some
sequence $f\to \infty$,
\begin{equation*}
    \sqrt{2\ln(p-s)} \;\leq \;\sqrt{n}\tau\; \leq\; \sqrt{n}\beta_{\min}^* + f.
\end{equation*}
For a block threshold selector $\hat S^b$, the ranges for the $\tau_j$'s are, for some sequences $f_j\to \infty$,
\begin{equation*}
  \sqrt{2\ln(p_j-s_j)} \;\leq\; \sqrt{n}\tau_j \;\leq\; \sqrt{n}\beta_{\min,j}^* + f_j.
\end{equation*}
As in the diverging $s_j$ regime, in the fixed $s_j$ regime the ranges for $\hat S^b$ are wider than that for $\hat S$. The necessary and sufficient assumptions to have variable selection consistency are then milder with block thresholds. The next corollary gives precise conditions under which consistent selection is possible with $\hat S^b$ but not with $\hat S$.
\begin{cor}\label{cor:sfin}
 In the sequence model \eqref{eq:GSM}, assume \hyperlink{lab:A1}{A1}, \hyperlink{lab:A2}{A2}, \hyperlink{lab:A4}{A4}, \hyperlink{lab:A11}{A11} and \hyperlink{lab:A12}{A12}.  If
\begin{eqnarray*}
&\lim_{n\to\infty}\sqrt{n}\beta_{\min}^*-\sqrt{2\ln(p-s)} \;<\;\infty
\end{eqnarray*}
then 
$\lim_{n\to\infty}P(\hat{S} = S)< 1$ and $\lim_{n\to\infty}P(\hat{S}^{b} = S)= 1$.  
\end{cor}

Let $\underline{\beta}_{\min,orth}^{*, b}$ and $\underline{\beta}_{\min , orth}^*$ be the smallest signal recoverable by $\hat{S}^b$ and $\hat{S}$ respectively. Assuming $\beta_{\min}^*$ is in block $b$, Assumptions~\hyperlink{lab:A4}{A4} and \hyperlink{lab:A12}{A12} require that $\underline{\beta}_{\min,orth}^{*, b}$ and $\underline{\beta}_{\min,orth }^*$ satisfy, for some sequences $g, h \rightarrow \infty$,
\begin{align*}
& \sqrt{n}\;\underline{\beta}_{\min ,orth}^{*, b} \geq \sqrt{2\ln \left(p_b-s_b\right)}+g, \quad\text{and}\\
& \sqrt{n}\;\underline{\beta}_{\min,orth}^* \geq \sqrt{2\ln(p-s)}+h.
\end{align*}
These lower bounds are the same as for $\underline{\beta}_{\min ,orth}^{*, b}$ and $ \underline{\beta}_{\min, orth}^*$ in the diverging $s_j$ case, up to logarithmic terms in the number of active signals, and up to $g$ and $h$ which can grow arbitrarily slowly with $n$. Note that in Examples 1, 2 and 4 in Section~\ref{sec:examples}, $\ln(s_j)=o(\ln(p_j-s_j))$ for all $j$. The discussion of the asymptotic behavior of the ratio $\underline{\beta}_{\min ,orth}^{*, b} / \underline{\beta}_{\min, orth}^*$ in those examples hence extend to the fixed $s_j$'s case. Finally, since Theorem~\ref{theo:fullrecovGS} holds both under Assumptions \hyperlink{lab:A3}{A3} and \hyperlink{lab:A11}{A11}, the discussion on the gains in terms of convergence rate in Sections~\ref{sec:prosconsblock} and \ref{sec:examples} remains valid here.

\subsection{Proofs}

\subsubsection{Proof of Proposition~\ref{prop:suffnecIIblockthreshbis}}
 By the union bound, 
$$P\big(\hat{S}^{b} \not\supseteq S\big)\leq \sum_{j=1}^b  P\left(\min_{i\in S_j}|y_i/\sqrt{n}|\leq\tau_j\right).$$ 
By Lemma~\ref{lemma:3tailbounds} (ii), for each $j$,	
$$P\left(\min_{i\in S_j}|y_i/\sqrt{n}|\leq\tau_j\right)\;\leq\; P\left(\max_{i\in S_j}|y_i/\sqrt{n}-\beta_i| \geq\beta_{\min,j}^*-\tau_j \right).$$ By Lemma~\ref{lemma:3tailbounds} (i), we have;
\begin{equation}
P\left(\hat{S}^{b}\not\supseteq S \right) \;\leq\;\sum_{j=1}^b\frac{e^{-\tfrac{n}{2}\left((\beta^*_{\min,j}-\tau_j)^2-\tfrac{2\ln(s_j)}{n}\right)}}{\sqrt{\pi\ln(s_j)}}.		
\end{equation}
Under Assumptions \hyperlink{lab:A11}{A11} and \hyperlink{lab:A12}{A12}, %under Assumption~\hyperlink{lab:A11}{A11}, 
the right-hand side vanishes, which proves part (i). 

We now prove part (ii). By Proposition~\ref{prop:suffnecblockthresh} (i), since \hyperlink{lab:A4}{A4} is assumed to hold, we also have $\lim_{n\to\infty}P\big(\hat{S}^{b} \subseteq S\big)=1$. This implies that $\lim_{n\to\infty}P\big(\hat{S}^{b} = S\big)= 1$.

\subsubsection{Proof of Lemma~\ref{lem:necbetaminblockthresh2}}
 First, we re-write 
 $$\sqrt{n}\beta_{\min,j}^*-\sqrt{2\ln(p_j-s_j)} \,=\, \sqrt{n}\beta_{\min,j}^*-\sqrt{n}\tau_j+\sqrt{2\ln(p_j-s_j)}\Big(\tfrac{\sqrt{n}\tau_j}{\sqrt{2\ln(p_j-s_j)}}-1\Big).$$ 
 Condition \eqref{eq:necbetaminblockthreshbis} implies that there exists $c\in\R^+$ such that
  \begin{equation}\label{eq:proofnecbetaminbis}
      \lim_{n\to \infty}\sqrt{n}\beta_{\min,j}^*-\sqrt{n}\tau_j+\sqrt{2\ln(p_j-s_j)}\Big(\tfrac{\sqrt{n}\tau_j}{\sqrt{2\ln(p_j-s_j)}}-1\Big)\leq c
  \end{equation}

Consider the case $\lim_{n\to \infty}\tfrac{\sqrt{n}\tau_j}{\sqrt{2\ln(p_j-s_j)}}-1<0$. Then by Proposition~\ref{prop:suffnecblockthresh} (ii),
$\lim_{n\to \infty}P(\hat{S}^{b} \subseteq S)< 1$ and $\lim_{n\to \infty}P(\hat{S}^{b} = S)< 1$.  Now consider the case $\lim_{n\to \infty}\tfrac{\sqrt{n}\tau_j}{\sqrt{2\ln(p_j-s_j)}}-1 \geq0$. Condition \eqref{eq:proofnecbetaminbis} then implies that $\lim_{n\to \infty}\sqrt{n}\beta_{\min,j}^*-\sqrt{n}\tau_j \leq c$. By Lemma~\ref{lem:gennecIIblockthresh}, we have that $\lim_{n\to \infty}P(\hat{S}^{b} \supseteq S)< 1$ and $\lim_{n\to \infty}P(\hat{S}^{b} = S)< 1$.

\subsubsection{Proof of Corollary~\ref{cor:sfin}} Observe that the conditions of Lemma~\ref{lem:necbetaminblockthresh2} hold for $\hat{S}$ ($\hat{S}^b$ for $b=1$), and then $\lim_{n\to \infty}P(\hat{S} = S)<1$. Since Assumptions~\hyperlink{lab:A4}{A4} and \hyperlink{lab:A12}{A12} hold, by Proposition~\ref{prop:suffnecblockthresh} (i) and Proposition~\ref{prop:suffnecIIblockthreshbis}, $\lim_{n\to \infty}P(\hat{S}^b \subseteq 1)= 1$ and $\lim_{n\to \infty}P( \hat{S}^b \supseteq S)=1$, and then $\lim_{n\to \infty}P(\hat{S}^b = S)=1$.

\section{Non linear block $\ell_0$ penalties in high-dimensional linear regression}\label{suppsec:appendixnonlinear}

\subsection{Selection properties of nonlinear block $\ell_0$ penalties}

In this section we show the variable selection consistency of the block $\ell_0$ penalties in linear regression without assuming linearity of the penalty functions. Results holds equally for fixed or diverging $p_j-s_j$ and $s_j$. We let, for all $j=1,\ldots,b$, $\eta_j$ be any non-negative and increasing functions on the natural numbers. The selector based on those block penalties is:
\begin{equation}
  \hat S^b\;\in\;\arg\max_{M\in \M} \left\{\max_{\beta\in \mathcal{L}_M} \ell(\bs y;\bs \beta)-\sum_{j=1}^b\eta_j (|M_j|)\right\}.
\end{equation}
Rewrite the difference in penalty between any model $M$ and $S$ as
\begin{equation}\label{eq:DeltaMbis}
\Delta_{MS} \; := \; \sum_{j=1}^b \eta_j(|M_j|)-\eta_j(|S_j|).
\end{equation}
We define the average block penalty when comparing $M$ and $S$ as
%in block penalty between any model $M$ and $S$ as the function:
\begin{equation}\label{eq:defkappatilde}
    \tilde{\kappa}_{j}(M)= \begin{cases}
     \frac{\eta_j(|M_j|)-\eta_j(|S_j|)}{|M_j|\,-\,|S_j|} \;\;\; \text{if}\;\;\; |M_j|\,\neq\,|S_j| \\
     0 \;\;\;\text{if}\;\;\; |M_j|\,=\,|S_j|.
 \end{cases}
\end{equation}
The function $\tilde{\kappa}_{j}$ plays a similar role to $\kappa_j$ in the linear penalty case. The quantity $\tilde{\kappa}_{j}(M)$ is the average penalty incurred for adding a variable from block $B_j$ to model $M$. If $\eta_j(|M_j|)$ is linear in $|M_j|$ for every $M$, then $\tilde{\kappa}_{j}= \kappa_j$.

To show the consistency of $\hat S^b$, we replace Assumption~\hyperlink{lab:A6}{A6} by an assumption on the $\tilde{\kappa}_{j}$'s, and we require a new betamin assumption.\\[5pt]
\hypertarget{lab:A13}{(A13)} \quad For each block $j$, there exists $f_j\to \infty$ (as $n\to \infty$) such that, for all $M \in \M$ such that  $|M_j| \neq |S_j|$ and $|M_j \setminus S_j| >0$, and for all sufficiently large $n$,
\begin{equation*}
    \tilde{\kappa}_{j}(M) \;=\;\ln\bigg(\frac{p_j-s_j}{|M_j \setminus S_j| }\bigg) + f_j(M).
\end{equation*}
\hypertarget{lab:A14}{(A14)}\quad For each block $j$, there exists $l_j\to \infty$ such that for all sufficiently large $n$,
    \begin{equation*}
        \sqrt{\frac{(1-\gamma) n\rho(\bs X)}{6}}\beta_{\min,j}^* \,- \,\sqrt{\max_{M\in \M}\tilde{\kappa}_{j}(M)} \;=\; \sqrt{\ln(s_j)} + l_j
    \end{equation*}
where $\gamma:=\frac12(1+\max_j \frac{\ln(p_j-s_j)}{\ln(p_j-s_j) +\min_{M: |M_j\setminus S_j|>0}f_j(M)})\in (\tfrac12,1)$.\\[5pt] 
We can now state the main result of this section.
\begin{thm}\label{theo:suffcondlinearmodelnonlinpen}

Under Assumptions~\hyperlink{lab:A1}{A1}, \hyperlink{lab:A13}{A13} and, \hyperlink{lab:A14}{A14}, we have
$$\sum_{M \in \M\setminus \{S\}}\E\left(NC(M)\right)\to 0 \qquad \mbox{and} \qquad P(\hat{S}^b = S) \to 1.$$
\end{thm}
Theorem~\ref{theo:suffcondlinearmodelnonlinpen} is consistent with results in the literature. A popular nonlinear penalty in high-dimensional variable selection is the EBIC penalty \cite{EBIC}, which sets for some $\zeta\geq 0$,
\begin{equation}\label{eq:EBICpen}
    \eta(|M|)=\zeta\ln{\binom{p}{|M|}} + \tfrac{1}{2}\ln(n).
\end{equation}
The penalty can be shown to satisfy Assumption~\hyperlink{lab:A13}{A13} under a restriction on the number of active signals. A corollary of Theorem~\ref{theo:suffcondlinearmodelnonlinpen} is as follows.
\begin{cor}\label{cor:ebic}
    Suppose that Assumption~\hyperlink{lab:A1}{A1} holds, $\eta$ is as in \eqref{eq:EBICpen} with $\zeta\geq1$, $s^{\zeta+1}=o(\sqrt{n})$, and that there exists $l\to \infty$ such that, for sufficiently large $n$,
    \begin{equation}\label{eq:betaminebic}
        \sqrt{\frac{ \big(\ln \big(\tfrac{\sqrt{n}}{(1+s)^{\zeta}}\big)+ k(s)\big)n\rho(\bs X)\beta_{\min}^*}{12\big(\ln(p-s)+ \ln \big(\tfrac{\sqrt{n}}{(1+s)^{\zeta}}\big)+k(s)\big)}}\,- \,\sqrt{\zeta\ln\big(p-s+1\big)+\zeta-1+\ln(\sqrt{n})} \;=\; \sqrt{\ln(s)} + l .
    \end{equation}
    where $k(s)=\zeta s\ln(1-s^{-1})$. Then Assumptions~\hyperlink{lab:A13}{A13} and \hyperlink{lab:A14}{A14} hold for $\eta$ and, $$\sum_{M \in \M\setminus \{S\}}\E\left(NC(M)\right)\to 0 \qquad \mbox{and} \qquad P(\hat{s}^b = S) \to 1.$$    
\end{cor}
Corollary~\ref{cor:ebic} shows the consistency of the EBIC penalty under milder betamin conditions than the literature \cite{EBIC,ebicdiverging}. It also shows that EBIC achieves \emph{strong} selection consistency and in a $L_1$ sense.

Observe that the assumptions and proof strategy of Theorem~\ref{theo:suffcondlinearmodelnonlinpen} are similar to those of Theorem~\ref{theo:suffcondlinearmodel}. We do not develop them here but results analogous to Theorem~\ref{prop:fullrecovreg} on convergence rates and the necessary conditions in Section~\ref{ssec:neccond_linreg} can be obtained for the nonlinear penalties. We also expect the benefits of linear block penalties to extend to the nonlinear ones.

\subsection{Proofs}
\subsubsection{Proof of Theorem~\ref{theo:suffcondlinearmodelnonlinpen}}
The proof is essentially the same as the proof of Theorem~\ref{theo:suffcondlinearmodel} replacing $\kappa_j$ by $\tilde{\kappa}_{j}$ for every $j=1,\ldots,b$. We first use Lemma~\ref{lem:boundexpnc} with $T = S$ to show that for every $M\neq S$, $\E(NC(M)) \leq e^{-\psi A_S}$ for every large enough $n$ and any $\psi \in (0,1)$, where $A_S=\gamma\Delta_{MS}+\tfrac{1-\gamma}{6}\mu_{Q_SM}$ (\emph{cf} \eqref{eq:DeltaMbis} and \eqref{eq:muQM}), $\gamma\in(1/2,1)$ is defined in Assumption~\hyperlink{lab:A14}{A14} and $Q_S= M\cup S$.
The second step is to obtain a lower bound for $A_S$, which gives a new upper bound for $\E(NC(M))$.
The final step is to use these bounds to get an upper-bound on $\sum_{M \in \M\setminus \{S\}} \E\left(NC(M)\right)$ that vanishes under Assumptions~\hyperlink{lab:A13}{A13} and \hyperlink{lab:A14}{A14}. We then use Lemma~\ref{lem:L0toL1convergence} to conclude on the vanishing of $P(\hat S^b \neq S)$. 

First, to show that $\E(NC(M)) \leq e^{-\psi A_S}$
for any $M\in\mathcal{M}\setminus \{S\}$, we show that $A_S$ satisfies the conditions of Lemma~\ref{lem:boundexpnc}, taking $T=S$. 
That is, we wish to show that, $A_S>0$, $|M\setminus S|=o(A_S)$, and $\mu_{Q_S S}=o(A_S)$. 
Observe that $\Delta_{MS}$, defined in \eqref{eq:DeltaMbis}, can be rewritten as $\Delta_{MS}=\sum_{j=1}^b(|M_j\setminus S_j|-|S_j\setminus M_j|)\tilde{\kappa}_j(M)$. By  Lemma~\ref{lem:boundrho}, for every $n\in \N$ we have 
\begin{equation}\label{eq:nonoverfittingcondbis}
\begin{aligned}
   A_S\;&=\; \gamma \Delta_{MS}+\frac{1-\gamma}{6}\mu_{Q_SM}\\
   &\geq\; \gamma\sum_{j=1}^b|M_j \setminus S_j|\tilde{\kappa}_j(M) + \sum_{j=1}^b |S_j\setminus M_j |\left(\tfrac{1-\gamma}{6} n \rho(\bs X) {\beta^*_{\min,j}}^{2} - \gamma\tilde{\kappa}_j(M)\right) 
\end{aligned}
\end{equation} 
Since $M\neq S$, $|M\setminus S|\neq0$ or $|S\setminus M|\neq0$, then by Assumptions~\hyperlink{lab:A13}{A13} and \hyperlink{lab:A4}{A14}, for every $n$ large enough, $A_S>0$. We immediately have $\mu_{Q_S S}=o(A_S)$ because $\beta^*_{Q_S\setminus S}=\beta^*_{M\setminus S}=0$ (any parameter outside the true support $S$ is by definition 0) and hence $\mu_{Q_S S}=0$.
If $|M\setminus S|=0$, $|M\setminus S|=o(A_S)$ also immediately. Consider now the case $|M\setminus S|\neq 0$. By Assumption~\hyperlink{lab:A4}{A14} the last term in \eqref{eq:nonoverfittingcondbis} is nonnegative, and hence
$$\frac{|M\setminus S|}{A_S}
= \frac{|M\setminus S|}{\gamma \Delta_{MS}+\frac{1-\gamma}{6}\mu_{Q_SM}}
\leq \bigg[ \gamma\sum_{j=1}^b\frac{|M_j \setminus S_j|}{|M\setminus S|}\tilde{\kappa}_j(M) \bigg]^{-1} \leq \Big[ \gamma\min_{j=1,\ldots,b|\tilde{\kappa}_j(M) \neq0}\tilde{\kappa}_j(M) \Big]^{-1} $$
where the last inequality follows from $\sum_{j=1}^b\frac{|M_j \setminus S_j|}{|M\setminus S|}=1$. By Assumption~\hyperlink{lab:A13}{A13}, $\min_{j|\tilde{\kappa}_j(M) \neq0} \tilde{\kappa}_j(M) \to \infty$ as $n \rightarrow \infty$, and hence $|M\setminus S|=o(A_S)$. Thus, %We conclude that for every $M\in \mathcal M\setminus \{S\}$, 
by Lemma~\ref{lem:boundexpnc}, for any $\psi\in(0,1)$ and all $n$ large enough,
$
\E(NC(M))\;\leq\;e^{-\psi A_S}.
$

For the second step of the proof, let $A_S^*$ be the lower bound for $A_S$ given in \eqref{eq:nonoverfittingcondbis}. That is 
$$A_S^*\;:=\;\gamma\sum_{j=1}^b|M_j \setminus S_j|\tilde{\kappa}_j(M) + \sum_{j=1}^b |S_j\setminus M_j |\left(\tfrac{1-\gamma}{6} n \rho(\bs X) {\beta^*_{\min,j}}^{2} - \gamma\tilde{\kappa}_j(M)\right).$$
By \eqref{eq:nonoverfittingcondbis}, we have for all $n$ large enough,
\begin{equation}\label{eq:boundexpbis}
    \E(NC(M))\;\leq\;e^{-\psi A_S^*}.
\end{equation}
Assumption~\hyperlink{lab:A14}{A14} implies that there exists $g'_j \to \infty$ such that
\begin{equation}\label{eq:ogbeabis}
    \frac{(1-\gamma) n\rho(\bs X)}{6}{\beta_{\min,j}^*}^2 \,- \,\tilde{\kappa}_j(M) \;=\; \ln(s_j) + g'_{j}.
\end{equation}
Let $\delta \in (0,1)$ and denote $\bar{m}_j=\max\big\{\frac{2\ln(p_j-s_j)}{f_j(M)},\frac{2\ln(s_j)}{g'_{j}}\big\}$ where $f_j$ is given in Assumption~\hyperlink{lab:A13}{A13}. 
Take $\psi=\max_{j=1,\ldots,b}\frac{\xi + \delta + \bar{m}_j}{1+\bar{m}_j}$ for some $\xi \in(0, 1- \delta)$ then $\psi \in (0,1)$ and we have, for every $j=1,\ldots,b$,
\begin{align}
    &\psi > \frac{\delta + \frac{2\ln(p_j-s_j)}{f_j(M)}}{1+\frac{2\ln(p_j-s_j)}{f_j(M)}} = \frac{\delta f_j(M)/2 + \ln(p_j-s_j)}{f_j(M)/2 +\ln(p_j-s_j)} \geq \frac{\delta f_j(M)/2 + \ln[(p_j-s_j)/|M_j\setminus S_j|]}{f_j(M)/2 +\ln[(p_j-s_j)/|M_j\setminus S_j|]}\label{eq:psifbis}\\
    &\psi > \frac{\delta + \frac{2\ln(s_j)}{g'_j}}{1+\frac{2\ln(s_j)}{g'_j}} = \frac{\delta g'_j/2 + \ln(s_j)}{g'_j/2 +\ln(s_j)} \geq \frac{\delta g'_j/2 + \ln(s_j)}{g'_j +\ln(s_j)}.\label{eq:psigbis}
\end{align}
By definition of $\gamma$ in Assumption~\hyperlink{lab:A14}{A14}, for all $M$ such that $|M_j\setminus S_j|>0$, $$\gamma \geq \frac12\Bigg(1+ \frac{\ln[(p_j-s_j)/|M_j\setminus S_j|]}{\ln[(p_j-s_j)/|M_j\setminus S_j|] + f_j(M)}\Bigg),$$ and it follows that 
\begin{align*}
    \gamma\tilde{\kappa}_j(M)
    &\;\geq\;\frac12 \Big(1+\frac{\ln[(p_j-s_j)/|M_j\setminus S_j|]}{\tilde{\kappa}_j(M)}\Big)\tilde{\kappa}_j(M)\\
    &\;=\;\ln\Big(\frac{p_j-s_j}{|M_j\setminus S_j|}\Big)+\frac12\Big(\tilde{\kappa}_j(M)-\ln\Big(\frac{p_j-s_j}{|M_j\setminus S_j|}\Big)\Big)\\
    & \;=\; \ln\Big(\frac{p_j-s_j}{|M_j\setminus S_j|}\Big)+ \frac12f_j(M).
\end{align*}
Hence, by \eqref{eq:psifbis}, when $|M_j\setminus S_j|>0$ we have
\begin{equation*}
\psi\gamma\tilde{\kappa}_j(M) \;\geq\; \psi\Big(\ln\Big(\frac{p_j-s_j}{|M_j\setminus S_j|}\Big)+ \frac12f_j(M)\Big)\\ \;\geq\; \ln\Big(\frac{p_j-s_j}{|M_j\setminus S_j|}\Big)+ \delta\frac12f_j(M).
\end{equation*}
Taking the minimum of $f_j(M)$ over $M\in \M$ such that $|M_j\setminus S_j|>0$, we get
\begin{equation*}
\psi\gamma\tilde{\kappa}_j(M) \;\geq\;  \ln\Big(\frac{p_j-s_j}{|M_j\setminus S_j|\vee 1}\Big)+ \delta\frac12 \min_{M: |M_j\setminus S_j|>0}f_j(M),
\end{equation*}
and then, for any $|M_j \setminus S_j| \geq 0$,
\begin{equation}\label{eq:ooofjfjbis}
    |M_j\setminus S_j|\psi\gamma\tilde{\kappa}_j(M) \geq |M_j\setminus S_j|\Big(\ln\Big(\frac{p_j-s_j}{|M_j\setminus S_j|\vee 1}\Big)+ \delta\frac12\min_{M: |M_j\setminus S_j|>0}f_j(M)\Big).
\end{equation}
Further using that $\gamma \in(0,1)$,
%By \hyperlink{lab:A14}{A14} and by our choice of $\psi$ and $\delta$,
\begin{equation}\label{eq:ooofjfj2bis}
\psi\left(\tfrac{1-\gamma}{6} n \rho(\bs X) {\beta^*_{\min,j}}^{2} - \gamma\tilde{\kappa}_j(M)\right) \;\geq\; \psi\left(\ln(s_j)+g_j'\right) \;\geq\; \ln(s_j)+\delta \frac{1}{2}g'_j.
\end{equation}
where the first inequality follows from \eqref{eq:ogbeabis} and the second inequality from \eqref{eq:psigbis}. In \eqref{eq:boundexpbis}, $\psi A^*_S=\sum_{j=1}^b|M_j \setminus S_j|\psi\gamma\tilde{\kappa}_j(M) + \sum_{j=1}^b |S_j\setminus M_j |\psi\left(\tfrac{1-\gamma}{6} n \rho(\bs X) {\beta^*_{\min,j}}^{2} - \gamma\tilde{\kappa}_j(M)\right)$. Then by \eqref{eq:ooofjfjbis} and \eqref{eq:ooofjfj2bis} , we get
\begin{align}\label{eq:fyvbis}
\E(NC(M))\;\leq\;\exp  \Bigg\{ & - \sum_{j=1}^b |M_j\setminus S_j|\big(\ln\big(\tfrac{p_j-s_j}{|M_j\setminus S_j|\vee 1}\big)+\delta \tfrac{\min_{M: |M_j\setminus S_j|>0}f_j(M)}{2}\big) \\ & -\sum_{j=1}^b |S_j\setminus M_j|(\ln(s_j)+\delta \tfrac{g'_j}{2})\Bigg\}.    \nonumber
\end{align}

For the final step of the proof, denote $\mathcal{S}= \sum_{M \in \M\setminus \{S\}} \E\left(NC(M)\right)$ for convenience. By \eqref{eq:fyvbis} we have that
\begin{align*}
    \mathcal{S} &\leq \sum_{M \in \M\setminus \{S\}}e^{- \sum_{j=1}^b |M_j\setminus S_j|\big(\ln\big(\tfrac{p_j-s_j}{|M_j\setminus S_j|\vee 1}\big)+\delta \tfrac{\min_{M: |M_j\setminus S_j|>0}f_j(M)}{2}\big)-\sum_{j=1}^b |S_j\setminus M_j|\big(\ln(s_j)+\delta \tfrac{g'_j}{2}\big)}.
\end{align*}
Observe that if $|M_j\setminus S_j|=0$ and $|S_j\setminus M_j|=0$ for all $j$, then $M=S$ and the summand in the right-hand side above is 1. Then by adding and resting 1 we get that
$$\mathcal{S} \leq -1+\sum_{M \in \M}e^{- \sum_{j=1}^b |M_j\setminus S_j|\big(\ln\big(\tfrac{p_j-s_j}{|M_j\setminus S_j|\vee 1}\big)+\delta \tfrac{\min_{M: |M_j\setminus S_j|>0}f_j(M)}{2}\big)-\sum_{j=1}^b |S_j\setminus M_j|\big(\ln(s_j)+\delta \tfrac{g'_j}{2}\big)}.$$
We can split the sum in the right-hand side above into sums over the models that have the same number of inactive variables and missing the same number of truly active variables in every block. That is, the models $M$ such that for all $j$, $|M_j\setminus S_j|=u_j$ and $|S_j\setminus M_j|=w_j$ with $u_j \in \{0,\ldots,p_j-s_j\}$ and $w_j \in \{0,\ldots,s_j\}$. Denote
$$S^{\bs u}_{\bs w}=\sum_{\substack{M \in \M:\;\forall \,j\;\; |M_j\setminus S_j|=u_j,\\ |S_j\setminus M_j|=w_j }} e^{- \sum_{j=1}^b u_j\big(\ln\big(\tfrac{p_j-s_j}{u_j\vee 1}\big)+\delta \tfrac{\min_{M: |M_j\setminus S_j|>0}f_j(M)}{2}\big)-\sum_{j=1}^b w_j\big(\ln(s_j)+\delta \tfrac{g'_j}{2}\big)}.$$
We get that
\begin{equation}\label{eq:uygfewbis}
    \mathcal{S} \leq -1+\sum_{w_1=0}^{s_1}\cdots\sum_{w_b=0}^{s_b}\sum_{u_1=0}^{p_1-s_1}\cdots\sum_{u_b=0}^{p_b-s_b}S^{\bs u}_{\bs w}.
\end{equation}
The number of models having, for all $j$, $u_j$ inactive parameters and missing $w_j$ out of the $s_j$ active parameters %with given collection of $u_j$'s and $w_j$'s 
is $\prod_{j=1}^b {\binom{p_j-s_j}{u_j}}{\binom{s_j}{w_j}}$. We thus have that
\begin{align*}
    S^{\bs u}_{\bs w}&= \bigg(\prod_{j=1}^b {\binom{p_j-s_j}{u_j}}{\binom{s_j}{w_j}}\bigg)e^{- \sum_{j=1}^b u_j\big(\ln\big(\tfrac{p_j-s_j}{u_j\vee 1}\big)+\delta \tfrac{\min_{M: |M_j\setminus S_j|>0}f_j(M)}{2}\big)-\sum_{j=1}^b w_j\big(\ln(s_j)+\delta \tfrac{g'_j}{2}\big)} \\
    &= \prod_{j=1}^b {\binom{p_j-s_j}{u_j}}e^{- u_j\big(\ln\big(\tfrac{p_j-s_j}{u_j\vee 1}\big)+\delta \tfrac{\min_{M: |M_j\setminus S_j|>0}f_j(M)}{2}\big)}{\binom{s_j}{w_j}}e^{-w_j\big(\ln(s_j)+\delta \tfrac{g'_j}{2}\big)}.
\end{align*}
Plugging the expression above into \eqref{eq:uygfewbis} gives that
\begin{align*}
    \mathcal{S} 
    \;\leq\; -1+\sum_{w_1=0}^{s_1}\cdots\sum_{w_b=0}^{s_b}\sum_{u_1=0}^{p_1-s_1}\cdots\sum_{u_b=0}^{p_b-s_b}\prod_{j=1}^b & {\binom{p_j-s_j}{u_j}}e^{- u_j\big(\ln\big(\tfrac{p_j-s_j}{u_j\vee 1}\big)+\delta \tfrac{\min_{M: |M_j\setminus S_j|>0}f_j(M)}{2}\big)}\\&. {\binom{s_j}{w_j}}e^{-w_j\big(\ln(s_j)+\delta \tfrac{g'_j}{2}\big)}
\end{align*}
and by factorizing,
\begin{align*}
    \mathcal{S}\;\leq\; -1+ \prod_{j=1}^b&\left(1+\sum_{u_j=1}^{p_j-s_j}{\binom{p_j-s_j}{u_j}}e^{- u_j\big(\ln\big(\tfrac{p_j-s_j}{u_j}\big)+\delta \tfrac{\min_{M: |M_j\setminus S_j|>0}f_j(M)}{2}\big)}\right)\\&.\left(1+
    \sum_{w_j=1}^{s_j}{\binom{s_j}{w_j}}
    e^{-w_j\big(\ln(s_j)+\delta \tfrac{g'_j}{2}\big)}\right).
\end{align*}
where the second inequality follows from first factorizing over terms in $u_j$ and $w_j$ and then taking the term in 0 out of every sum.
By the bound on the binomial coefficient in \eqref{eq:upperboudbinom}, we have that
\begin{equation}\label{eq:upperbound_sumprobbis}
    \mathcal{S} \;\;\leq\;\; -1+\prod_{j=1}^b\left(1+\sum_{u_j=1}^{p_j-s_j}e^{-u_j\left(\delta \tfrac{\min_{M: |M_j\setminus S_j|>0}f_j(M)}{2}-1\right)}\right)\left(1+
    \sum_{w_j=1}^{s_j}    e^{-w_j\left(\delta \tfrac{g'_j}{2}-1\right)}\right).
\end{equation}
Denote
$$
d_j=e^{1-\delta \frac{\min_{M: |M_j\setminus S_j|>0}f_j(M)}{2}}, \quad h_j=e^{1-\delta \frac{g_j^{\prime}}{2}}.
$$
where both expressions go to zero as $n$ increases since $\min_{M: |M_j\setminus S_j|>0}f_j(M) \rightarrow \infty$ and $g_j^{\prime} \rightarrow \infty$. For every $j$, by the properties of geometric sums, we have
$$
\begin{aligned}
& 1+\sum_{u_j=1}^{p_j-s_j} e^{-u_j\left(\delta \frac{\min_{M: |M_j\setminus S_j|>0}f_j(M)}{2}-1\right)}=\frac{1-d_j^{p_j-s_j+1}}{1-d_j}, \\
& 1+\sum_{w_j=1}^{s_j} e^{-w_j\left(\delta \frac{g_j^{\prime}}{2}-1\right)}=\frac{1-h_j^{s_j+1}}{1-h_j}.
\end{aligned}
$$
Since both expressions converge to 1 as $n$ grows, we get $$\lim _{n \rightarrow \infty} S=\lim _{n \rightarrow \infty} \sum_{M \in \mathcal{M} \backslash\{S\}}\E(NC(M))=0.$$ By Lemma 4.2, $P\left(\hat{S}^b \neq S\right) \leq 2 S$ and then $\lim _{n \rightarrow \infty} P\left(\hat{S}^b=S\right)=1$.

\subsubsection{Proof of Corollary~\ref{cor:ebic}}
The proof strategy is to first show that $\eta$ satisfies Assumption~\hyperlink{lab:A13}{A13}, then that \eqref{eq:betaminebic} implies Assumption~\hyperlink{lab:A14}{A14}. The consistency results then follow from Theorem~\ref{theo:suffcondlinearmodelnonlinpen}.

To show that $\eta$ satisfies Assumption~\hyperlink{lab:A13}{A13}, we show that, for all $M\neq S$ such that $|M\setminus S|>0$ and $\tilde{\kappa}(M)$ defined in \eqref{eq:defkappatilde}, the function $f(M):=\tilde{\kappa}(M)-\ln(p-s/|M\setminus S|)$ is lower bounded by a diverging sequence. 

Let $M\in\M$ such that $M\neq S$ and $|M\setminus S|>0$. Denote $|M\setminus S|=u$ and $|S\setminus M|=w$. We have $|M|=u+s-w$ and $|M|-|S|=u-w$. Since $M\neq S$ and $|M\setminus S|>0$, we have $u- w\neq0$ and $u>0$. We consider first the case where $u-w>0$, we have
    \begin{equation*}
        \tilde{\kappa}(M)=\frac{\zeta}{u-w}\ln\Bigg[\frac{\binom{p}{|M|}}{\binom{p}{|S|}}\Bigg]+ \frac{1}{2}\ln(n).
    \end{equation*}
    A well-known property of binomial coefficients is that for any positive integers $n$, $h$, $k$ we have
    \begin{equation}\label{eq:idbinom}
        \binom{n}{h}\binom{n-h}{k}=\binom{n}{h+k}\binom{h+k}{h}.
    \end{equation}
    Taking $n=p$, $k=u-w$ and $h=s$ in \eqref{eq:idbinom}, we get
    \begin{equation}\label{eq:identitybinom}
        \tilde{\kappa}(M)=\frac{\zeta}{u-w}\ln\bigg[\frac{\binom{p -s}{u-w}}{\binom{|M|}{s}}\bigg]+ \frac{1}{2}\ln(n).
    \end{equation}
Standard bounds on binomial coefficient for $1\leq k\leq n$ are
\begin{equation}\label{eq:boundbinom}
  \left(\frac{n}{k}\right)^k \leq  \binom{n}{k} \leq \frac{n^n}{k^k(n-k)^{n-k}}.
\end{equation}
Using the bounds in \eqref{eq:boundbinom} in \eqref{eq:identitybinom}, we get
\begin{equation}\label{eq:gggg}
    \frac{\zeta}{u-w}\ln\bigg[\frac{\binom{p}{|M|}}{\binom{p}{s}}\bigg] \geq \zeta\ln\bigg(\frac{p-s}{u-w}\bigg)-\zeta\bigg[\frac{s}{u-w}\ln\bigg(\frac{u-w}{s}+1\bigg) +\ln\bigg(1+\frac{s}{u-w}\bigg) \bigg].
\end{equation}
We have
\begin{equation}\label{eq:gggg2}
\zeta\ln\bigg(\frac{p-s}{u-w}\bigg) \geq \zeta\ln\bigg(\frac{p-s}{u}\bigg) \geq \ln\bigg(\frac{p-s}{u}\bigg).\end{equation}
where the first inequality follows from $u-w\leq u$ and the second from $\zeta\geq 1$.
Observe also that $h(x)= \frac{\ln(x+1)}{x}+\ln(1+{x}^{-1})$ is decreasing for $x>0$ and that $\frac{u-w}{s} \geq s^{-1}$. By \eqref{eq:gggg} and \eqref{eq:gggg2}, we then have 
\begin{equation*}
        \tilde{\kappa}(M) \geq \ln\big(\frac{p-s}{u}\big)-\zeta s\ln(s^{-1}+1)+ \ln \big(\tfrac{\sqrt{n}}{(1+s)^{\zeta}}\big),
\end{equation*}
and for every $M$ such that $u-w>0$,
\begin{equation}\label{eq:lowerboundb1ebic2}
        f(M) \geq -\zeta s\ln(s^{-1}+1)+ \ln \big(\tfrac{\sqrt{n}}{(1+s)^{\zeta}}\big).
\end{equation}
Since $\lim_{n\to\infty}s\ln(s^{-1}+1)=1$ and by assumption $s^{\zeta+1}=o(\sqrt{n})$, $\eta$ satisfies Assumption~\hyperlink{lab:A13}{A13} in the case of $M$ such that $u-w>0$.

Consider now the case where $M$ is such that $u-w<0$. We have
     \begin{equation*}
        \tilde{\kappa}(M)=\frac{\zeta}{w-u}\ln\Bigg[\frac{\binom{p}{s}}{\binom{p}{|M|}}\Bigg]+ \frac{1}{2}\ln n.
    \end{equation*}
 Taking $n=p$, $k=w-u$ and $h=|M|$ in \eqref{eq:idbinom} gives
     \begin{equation}\label{eq:identitybinom2}
     \tilde{\kappa}(M)=\frac{\zeta}{w-u}\ln\Bigg[\frac{\binom{p - |M|}{w-u}}{\binom{s}{|M|}}\Bigg]+ \frac{1}{2}\ln n.
    \end{equation}
Using the bounds in \eqref{eq:boundbinom}, we get
\begin{equation}\label{eq:PVnk}
    \frac{\zeta}{w-u}\ln\bigg[\frac{\binom{p}{s}}{\binom{p}{|M|}}\bigg] \geq \zeta\ln\bigg(\frac{p-|M|}{w-u}\bigg)+\zeta\bigg[\frac{s}{w-u}\ln\bigg(1-\frac{w-u}{s}\bigg) -\ln\bigg(\frac{s}{w-u}-1\bigg) \bigg].
\end{equation}
We have
$$\ln\bigg(\frac{p-|M|}{w-u}\bigg)=\ln\bigg(\frac{p-s}{u}\bigg)+\ln\bigg(\frac{p-|M|}{p-s}\bigg)+\ln\bigg(\frac{u}{w-u}\bigg).$$
Since $u-w<0$, we have $|M|<s$ and $\ln((p-|M|)/(p-s))\geq 0$. Since $w\leq s$ and $u\geq 1$, $\ln(u/(w-u))\geq\ln(1/(s-1))$. Using also that $\zeta\geq 1$, we have
\begin{equation}\label{eq:feIH}
    \zeta\ln\bigg(\frac{p-|M|}{w-u}\bigg) \geq \ln\bigg(\frac{p-s}{u}\bigg) - \ln(s-1)
\end{equation}
where the right-hand side is well defined because since $u>0$ and $u-w<0$, we have $2\leq w\leq s$. Observe that $g:x \mapsto \frac{\ln(1-x)}{x}-\ln({x}^{-1}-1)$ for $x\in(0,1)$ is increasing and that $1>\frac{w-u}{s} \geq s^{-1}$. By \eqref{eq:PVnk} and \eqref{eq:feIH}, we then get for all $M\in \M$ such that $u-w<0$, 
\begin{equation*}
        \tilde{\kappa}(M) \geq \ln\big(\frac{p-s}{u}\big)+\zeta s\ln(1-s^{-1})+ \ln \big(\tfrac{\sqrt{n}}{(s-1)^{\zeta+1}}\big).
\end{equation*}
and for every $M\in\mathcal{M}$ such that $w-u>0$,
\begin{equation}\label{eq:bounbkappaebicunderfit}
        f(M) \geq \zeta s\ln(1-s^{-1})+ \ln \big(\tfrac{\sqrt{n}}{(s-1)^{\zeta+1}}\big).
\end{equation}
Since $\lim_{n\to\infty}s\ln(1-s^{-1})=-1$ and by assumption $s^{\zeta+1}=o(\sqrt{n})$, $\eta$ satisfies Assumption~\hyperlink{lab:A13}{A13} in the case $u-w<0$ too.

We now show that if \eqref{eq:betaminebic} holds then Assumption~\hyperlink{lab:A14}{A14} holds.
Assumption~\hyperlink{lab:A14}{A14} for $\eta$ as in \eqref{eq:EBICpen} states that there exists $l_j\to \infty$ such that for large enough $n$,
    \begin{equation*}
        \sqrt{\frac{(1-\gamma) n\rho(\bs X)}{6}}{{\beta_{\min}^*}} \,- \,\sqrt{\max_{M\in\M}\tilde{\kappa}(M)} \;=\; \sqrt{\ln(s)} + l_j .
    \end{equation*}
where $\gamma$ takes value
\begin{equation}\label{eq:gammaebic}
    \gamma=\frac12\Big(1+\frac{\ln(p-s)}{\ln(p-s)+\min_{M: |M\setminus S|>0}f(M)}\Big).
\end{equation}
To show that if \eqref{eq:betaminebic} holds then Assumption~\hyperlink{lab:A14}{A14} holds for $\eta$, it suffices to show that the following two inequalities
\begin{align}
    - \,\sqrt{\max_{M\in\M}\tilde{\kappa}(M)} \,&\geq \,-\sqrt{\zeta\ln\big(p-s+1\big)+\zeta-1+\frac{1}{2}\ln(n)}\label{eq:oaebr4bis},\quad\text{and} \\
    \sqrt{\frac{ (1-\gamma)n\rho(\bs X)}{6}}{\beta_{\min,j}^*} \,&\geq \,  \sqrt{\frac{ \ln \big(\tfrac{\sqrt{n}}{(1+s)^{\zeta}}\big)+\zeta s\ln(1-s^{-1})}{\ln(p-s)+ \ln \big(\tfrac{\sqrt{n}}{(1+s)^{\zeta}}\big)+\zeta s\ln(1-s^{-1})}\frac{ n\rho(\bs X)}{12}}{\beta_{\min,j}^*} \label{eq:oaebr3bis}
\end{align}
hold. We start with \eqref{eq:oaebr4bis}. If $u-w>0$, by \eqref{eq:identitybinom}, the upper bound in \eqref{eq:upperboudbinom}, and the lower bound in \eqref{eq:boundbinom}, then
 \begin{eqnarray*}
     \tilde{\kappa}(M) &\leq& \zeta\ln\Big(\frac{(p-s)e}{u-w}\Big)-\zeta(1+\frac{s}{u-w})\ln\Big(1+\frac{u-w}{s}\Big)+\frac{1}{2}\ln(n).
 \end{eqnarray*}
 Using that, for $x\geq1$, $(1+\frac{1}{x})\ln(1+x)\geq 1$ and $\ln\big(\frac{(p-s)e}{x}\big)\leq \ln(p-s)+1$. We get that, for all $M$ such that $u-w>0$, 
  \begin{eqnarray}\label{eq:pifr} 
     \tilde{\kappa}(M) &\leq& \zeta\ln(p-s)+\zeta-1+\frac{1}{2}\ln(n).
 \end{eqnarray}
If $w-u>0$, by \eqref{eq:identitybinom2}, the upper bound in \eqref{eq:upperboudbinom}, and the lower bound in \eqref{eq:boundbinom}, then
 \begin{eqnarray*}
     \tilde{\kappa}(M) &\leq& 
     \zeta\ln\Big(\frac{(p-s+(w-u))e}{w-u}\Big)+\zeta(\frac{s}{w-u}-1)\ln\Big(1
     -\frac{w-u}{s}\Big)+\frac{1}{2}\ln(n).
 \end{eqnarray*}
Using that, for $x\geq 1$, $(1+\frac{1}{x})\ln(1+x)\geq 1$ and $\ln\Big(\frac{(p-s+x)e}{x}\Big)\leq \ln(p-s+1)+1$, we get that for all $M$ such that $w-u>0$,
\begin{eqnarray}\label{eq:pifr2}
     \tilde{\kappa}(M) &\leq& \zeta\ln\big(p-s+1\big)+\zeta-1+\frac{1}{2}\ln(n).
 \end{eqnarray}
By \eqref{eq:pifr} and \eqref{eq:pifr2},
 \begin{eqnarray*}
     \max_{M\in\M}\tilde{\kappa}(M)  &\leq& \zeta\ln\big(p-s+1\big)+\zeta-1+\frac{1}{2}\ln(n).
 \end{eqnarray*}
which shows \eqref{eq:oaebr4bis}. We now show \eqref{eq:oaebr3bis}. By \eqref{eq:lowerboundb1ebic2} and \eqref{eq:bounbkappaebicunderfit}, we have
$$\min_{M: |M\setminus S|>0}f(M) \geq \zeta s\ln(1-s^{-1})+ \ln \big(\tfrac{\sqrt{n}}{(1+s)^{\zeta}}\big).$$
If follows that
 \begin{align*}
 \gamma&=\frac12\Big(1+\frac{\ln(p-s)}{\ln(p-s)+\min_{M: |M\setminus S|>0}f(M)}\Big) \\&\leq \frac12\Big(1+\frac{\ln(p-s)}{\ln(p-s)+ \ln \big(\tfrac{\sqrt{n}}{(1+s)^{\zeta}}\big)+\zeta s\ln(s^{-1}-1)}\Big).
 \end{align*}
 Simple algebra gives
  \begin{equation*}
 1-\gamma\geq \frac12\frac{ \ln \big(\tfrac{\sqrt{n}}{(1+s)^{\zeta}}\big)+\zeta s\ln(1-s^{-1})}{\ln(p-s)+ \ln \big(\tfrac{\sqrt{n}}{(1+s)^{\zeta}}\big)+\zeta s\ln(1-s^{-1})}
 \end{equation*}
 which shows \eqref{eq:oaebr3bis} and that \eqref{eq:betaminebic} is sufficient for  Assumption~\hyperlink{lab:A14}{A14} to hold.
 
 Since Assumptions~\hyperlink{lab:A13}{A13} and \hyperlink{lab:A14}{A14} hold for $\eta$ as in \eqref{eq:EBICpen}, by Theorem~\ref{theo:suffcondlinearmodelnonlinpen}, $\lim_{n\to\infty}P(\hat{S}^{b}\neq S)= 1$, as we wished to prove.


\section{Tightness of conditions for variable selection consistency in linear regression}\label{suppsec:tightness}

We compare our sufficient conditions for variable selection consistency for standard $\ell_0$ selector $\hat{S}$ to those for an optimal selector that knows $s$ analyzed in \cite{infotheowainwright} and to our necessary conditions. This section is organized follows. We first recall our sufficient conditions, those in \cite{infotheowainwright} and our necessary conditions. We then proceed to compare them.

Theorem~\ref{theo:suffcondlinearmodel} shows variable selection consistency with $\hat{S}^b$ under Assumptions~\hyperlink{lab:A6}{A6} and \hyperlink{lab:A7}{A7}. By \eqref{eq:gpkr} in the proof of Theorem~\ref{prop:fullrecovreg}, an assumption slightly less stringent than \hyperlink{lab:A7}{A7}, and easier to analyze, is sufficient together with \hyperlink{lab:A6}{A6}. That assumption is: \\[5pt]
\noindent\hypertarget{lab:A15}{(A15)} \quad for each block $j$, there exists $g_j\to \infty$ such that for every sufficiently large $n$,
    \begin{equation*}
        \frac{(1-\gamma) n\rho(\bs X)}{6}{\beta_{\min,j}^*}^2 \,- \,\kappa_j \;=\; \ln(s_j) + g_j .
    \end{equation*}
where $\gamma:=\tfrac12(1+\max_j\ln(p_j-s_j)/\kappa_j)\in (\tfrac12,1)$

Consider assumptions~\hyperlink{lab:A7}{A7} and \hyperlink{lab:A15}{A15} for standard $\ell_0$ selector $\hat{S}$ with single penalty $\kappa$. Their combination implies a condition on the quadruplet $(n,p,s,\beta_{\min}^*)$. To simplify the analysis of that condition, we assume $\kappa=(1+\varepsilon)\ln(p-s)$ for some fixed $\varepsilon>0$. This choice guarantees that $\kappa$ meets assumption~\hyperlink{lab:A6}{A6} and that $1-\gamma$ is constant and bounded away from 0. A sufficient condition on $(n,p,s,\beta_{\min}^*)$ that follows from assumptions~\hyperlink{lab:A7}{A7} and \hyperlink{lab:A15}{A15} is then that there exists $t\to \infty$, growing at an arbitrarily slow rate, such that:
\begin{equation}\label{eq:groh}
    n= \frac{12(1+\varepsilon)}{\varepsilon}\frac{(1+\varepsilon)\ln(p-s)+\ln(s)}{\rho(\bs X){\beta_{\min}^*}^2} \;+\;t
\end{equation}

In \cite{infotheowainwright}, it is shown that a sufficient condition on $(n,p,s,\beta_{\min}^*)$ for an optimal selector that knows $s$ to be variable selection consistent is:
\begin{equation}\label{eq:groh2}
    n > (c_1+2048) \max \left\{\log \binom{p-s}{s}, \frac{\log (p-s)}{\rho(\bs X) {\beta_{\min}^*}^2}\right\}
\end{equation}
for some $c_1 >0$. 

Corollary~\ref{cor:necrecovreg} shows condition~\eqref{eq:betaminnecreg} is necessary to get variable selection consistency with $\hat{S}^b$. When applied to $\hat{S}$, it implies the necessary condition on $(n,p,s,\beta_{\min}^*)$,
\begin{equation}\label{eq:gpiR}
\lim_{n\to\infty}\frac{\sqrt{n\bar{\lambda}}\beta_{\min}^*}{\underline{\lambda}\sqrt{\ln(p-s)}} > 0.
\end{equation}

We first observe that for any regime of $(n,p,s,\beta_{\min}^*)$ such that $\ln(s)=O(\rho(\bs X){\beta_{\min}^*}^2)$, \eqref{eq:groh} is less stringent than \eqref{eq:groh2}. It is also the case when $\rho(\bs X){\beta_{\min}^*}^2=\Theta(1)$ and $s<p/2$ for example. Table~\ref{tab:scalings} gives, for some regimes of interest, the scalings of \eqref{eq:groh}, \eqref{eq:groh2}, and \eqref{eq:gpiR} where we assume $\underline{\lambda}$ and $\bar{\lambda}$ are bounded for simplicity.
%\paul{Table S1 takes the \textbf{format and regimes} of a table in "Information-Theoretic Limits on Sparse Signal Recovery: Dense versus Sparse Measurement Matrices" by Wang; Wainwright; Ramchandran. Note that the sufficient and necessary conditions in that paper are for a slightly stronger problem where $S$ is not fixed and considered random and you study the vanishing of $\max_{S \in \text{some set}}P(\hat{S}\neq S)$. They do not apply to our setting although they are similar in quite a few regimes.} 
\begin{table*}[ht]
\caption{Scaling of conditions for variable selection consistency}
\label{tab:scalings}
\begin{tabular}{@{}cccccc@{}}
\hline \\[-1em]
Regime & \thead{Our sufficient \\ condition} & \thead{Sufficient condition\\ as in \cite{infotheowainwright}} & \thead{Our necessary \\ condition} \\
\hline\\[-1em]
\thead{$s=\Theta(p)$\\$\rho(\bs X){\beta_{\min}^*}^2=\Theta(1/s)$}    & $\Theta(p\ln(p))$ & $\Theta(p\ln(p))$ & $\Theta(p\ln(p))$ \\\hline\\[-1em]
\thead{$s=\Theta(p)$\\$\rho(\bs X){\beta_{\min}^*}^2=\Theta\big(\ln(s)/s\big)$}    & $\Theta(p)$ & $\Theta(p)$ & $\Theta(p)$ \\\hline\\[-1em]
\thead{$s=\Theta(p)$\\$\rho(\bs X){\beta_{\min}^*}^2=\Theta(1)$}   & $\Theta(\ln(p))$ & $\Theta(p)$ & $\Theta(\ln(p))$ \\\hline\\[-1em]
\thead{$s=o(p)$\\$\rho(\bs X){\beta_{\min}^*}^2=\Theta(1/s)$}    & $\Theta(s\ln(p))$ & $\Theta(s\ln(p))$ & $\Theta(s\ln(p))$ \\\hline\\[-1em]
\thead{$s=o(p)$\\$\rho(\bs X){\beta_{\min}^*}^2=\Theta(\ln(s)/s)$}    & $\Theta\big(s\ln(p)/\ln(s)\big)$ & $\Theta(s\ln(p))$ & $\Theta\big(s\ln(p)/\ln(s)\big)$ \\\hline\\[-1em]
\thead{$s=o(p)$\\$\rho(\bs X){\beta_{\min}^*}^2=\Theta(1)$}    & $\Theta(\ln(p))$ & $\Theta(s\ln(p))$ & $\Theta(\ln(p))$ \\\hline
\end{tabular}
\end{table*}
The scalings implied by our sufficient conditions match or improve those implied by sufficient conditions of the optimal selector in \cite{infotheowainwright}. The scalings implied by our sufficient conditions also match those implied by our necessary conditions, confirming the tightness of our results. 

\end{document}
